%%
% !TeX root = main.tex
%% This is file `sample-acmsmall-submission.tex',
%% generated with the docstrip utility.
%%
%% The original source files were:
%%
%% samples.dtx  (with options: `all,journal,bibtex,acmsmall-submission')
%% 
%% IMPORTANT NOTICE:
%% 
%% For the copyright see the source file.
%% 
%% Any modified versions of this file must be renamed
%% with new filenames distinct from sample-acmsmall-submission.tex.
%% 
%% For distribution of the original source see the terms
%% for copying and modification in the file samples.dtx.
%% 
%% This generated file may be distributed as long as the
%% original source files, as listed above, are part of the
%% same distribution. (The sources need not necessarily be
%% in the same archive or directory.)
%%
%%
%% Commands for TeXCount
%TC:macro \cite [option:text,text]
%TC:macro~\cite [option:text,text]
%TC:macro \citet [option:text,text]
%TC:envir table 0 1
%TC:envir table* 0 1
%TC:envir tabular [ignore] word
%TC:envir displaymath 0 word
%TC:envir math 0 word
%TC:envir comment 0 0
%%
%% The first command in your LaTeX source must be the \documentclass
%% command.
%%
%% For submission and review of your manuscript please change the
%% command to \documentclass[manuscript, screen, review]{acmart}.
%%
%% When submitting camera ready or to TAPS, please change the command
%% to \documentclass[sigconf]{acmart} or whichever template is required
%% for your publication.
%%
%%
\documentclass[acmsmall,screen,review,anonymous]{acmart}
\usepackage[utf8]{inputenc}
\usepackage{tikz}
\usepackage{booktabs} % 必须：提供三线表命令
\usepackage{tabularx} % 可选：提供自动换行和宽度控制
\usepackage{amsmath}  % 供符号使用
\usepackage{bm}
\usepackage{diagbox}   % 用于 \diagbox 表头分割
\usepackage{multirow}  % 表格跨行
\usepackage{enumitem}  % itemize 选项
%%
%% \BibTeX command to typeset BibTeX logo in the docs
\AtBeginDocument{%
  \providecommand\BibTeX{{%
    Bib\TeX}}}

%% Rights management information.  This information is sent to you
%% when you complete the rights form.  These commands have SAMPLE
%% values in them; it is your responsibility as an author to replace
%% the commands and values with those provided to you when you
%% complete the rights form.
\setcopyright{acmlicensed}
\copyrightyear{2018}
\acmYear{2018}
\acmDOI{XXXXXXX.XXXXXXX}

%%
%% These commands are for a JOURNAL article.
\acmJournal{JACM}
\acmVolume{37}
\acmNumber{4}
\acmArticle{111}
\acmMonth{8}

%%
%% Submission ID.
%% Use this when submitting an article to a sponsored event. You'll
%% receive a unique submission ID from the organizers
%% of the event, and this ID should be used as the parameter to this command.
%%\acmSubmissionID{123-A56-BU3}

%%
%% For managing citations, it is recommended to use bibliography
%% files in BibTeX format.
%%
%% You can then either use BibTeX with the ACM-Reference-Format style,
%% or BibLaTeX with the acmnumeric or acmauthoryear sytles, that include
%% support for advanced citation of software artefact from the
%% biblatex-software package, also separately available on CTAN.
%%
%% Look at the sample-*-biblatex.tex files for templates showcasing
%% the biblatex styles.
%%

%%
%% The majority of ACM publications use numbered citations and
%% references.  The command \citestyle{authoryear} switches to the
%% ``author year'' style.
%%
%% If you are preparing content for an event
%% sponsored by ACM SIGGRAPH, you must use the ``author year'' style of
%% citations and references.
%% Uncommenting
%% the next command will enable that style.
%%\citestyle{acmauthoryear}


%%
%% end of the preamble, start of the body of the document source.
\begin{document}
%%
%% The ``title'' command has an optional parameter,
%% allowing the author to define a ``short title'' to be used in page headers.
\title{A Survey of Embodied AI for Robotic Arms: Embodiments, Interaction Data, and Learning Paradigms}

%%
%% The ``author'' command and its associated commands are used to define
%% the authors and their affiliations.
%% Of note is the shared affiliation of the first two authors, and the
%% ``authornote'' and ``authornotemark'' commands
%% used to denote shared contribution to the research.
\author{Ben Trovato}
\authornote{Both authors contributed equally to this research.}
\email{trovato@corporation.com}
\orcid{1234-5678-9012}
\author{G.K.M. Tobin}
\authornotemark[1]
\email{webmaster@marysville-ohio.com}
\affiliation{%
  \institution{Institute for Clarity in Documentation}
  \city{Dublin}
  \state{Ohio}
  \country{USA}
}

\author{Lars Th{\o}rv{\"a}ld}
\affiliation{%
  \institution{The Th{\o}rv{\"a}ld Group}
  \city{Hekla}
  \country{Iceland}}
\email{larst@affiliation.org}

\author{Valerie B\'eranger}
\affiliation{%
  \institution{Inria Paris-Rocquencourt}
  \city{Rocquencourt}
  \country{France}
}

\author{Aparna Patel}
\affiliation{%
 \institution{Rajiv Gandhi University}
 \city{Doimukh}
 \state{Arunachal Pradesh}
 \country{India}}

\author{Huifen Chan}
\affiliation{%
  \institution{Tsinghua University}
  \city{Haidian Qu}
  \state{Beijing Shi}
  \country{China}}

\author{Charles Palmer}
\affiliation{%
  \institution{Palmer Research Laboratories}
  \city{San Antonio}
  \state{Texas}
  \country{USA}}
\email{cpalmer@prl.com}

\author{John Smith}
\affiliation{%
  \institution{The Th{\o}rv{\"a}ld Group}
  \city{Hekla}
  \country{Iceland}}
\email{jsmith@affiliation.org}

\author{Julius P. Kumquat}
\affiliation{%
  \institution{The Kumquat Consortium}
  \city{New York}
  \country{USA}}
\email{jpkumquat@consortium.net}

%%
%% By default, the full list of authors will be used in the page
%% headers. Often, this list is too long, and will overlap
%% other information printed in the page headers. This command allows
%% the author to define a more concise list
%% of authors' names for this purpose.
\renewcommand{\shortauthors}{Trovato et al.}

%%
%% The abstract is a short summary of the work to be presented in the
%% article.
\begin{abstract}
  A clear and well-documented \LaTeX\ document is presented as an
  article formatted for publication by ACM in a conference proceedings
  or journal publication. Based on the ``acmart'' document class, this
  article presents and explains many of the common variations, as well
  as many of the formatting elements an author may use in the
  preparation of the documentation of their work.
\end{abstract}

%%
%% The code below is generated by the tool at http://dl.acm.org/ccs.cfm.
%% Please copy and paste the code instead of the example below.
%%
\begin{CCSXML}
<ccs2012>
 <concept>
  <concept_id>00000000.0000000.0000000</concept_id>
  <concept_desc>Do Not Use This Code, Generate the Correct Terms for Your Paper</concept_desc>
  <concept_significance>500</concept_significance>
 </concept>
 <concept>
  <concept_id>00000000.00000000.00000000</concept_id>
  <concept_desc>Do Not Use This Code, Generate the Correct Terms for Your Paper</concept_desc>
  <concept_significance>300</concept_significance>
 </concept>
 <concept>
  <concept_id>00000000.00000000.00000000</concept_id>
  <concept_desc>Do Not Use This Code, Generate the Correct Terms for Your Paper</concept_desc>
  <concept_significance>100</concept_significance>
 </concept>
 <concept>
  <concept_id>00000000.00000000.00000000</concept_id>
  <concept_desc>Do Not Use This Code, Generate the Correct Terms for Your Paper</concept_desc>
  <concept_significance>100</concept_significance>
 </concept>
</ccs2012>
\end{CCSXML}

\ccsdesc[500]{Do Not Use This Code~Generate the Correct Terms for Your Paper}
\ccsdesc[300]{Do Not Use This Code~Generate the Correct Terms for Your Paper}
\ccsdesc{Do Not Use This Code~Generate the Correct Terms for Your Paper}
\ccsdesc[100]{Do Not Use This Code~Generate the Correct Terms for Your Paper}

%%
%% Keywords. The author(s) should pick words that accurately describe
%% the work being presented. Separate the keywords with commas.
\keywords{Do, Not, Use, This, Code, Put, the, Correct, Terms, for,
  Your, Paper}

\received{XXX}
\received[revised]{XXXX}
\received[accepted]{XXXX}

%%
%% This command processes the author and affiliation and title
%% information and builds the first part of the formatted document.
\maketitle

\section{Introduction}\label{sec:introduction}

Embodied intelligence argues that cognition is inseparable from the body, its sensors, and its continuous interaction with the physical world~\cite{brooksIntelligenceWithoutRepresentation1991, pfeiferHowBodyShapes2012, clarkBeingTherePutting1998}. Among all embodiments, manipulator robots occupy a unique position: they are the most mature and widely deployed platforms, yet the gap between their mechanical reliability and their \emph{manipulation intelligence} remains vast~\cite{sicilianoSpringerHandbookRobotics2016}. Most industrial arms still execute pre-programmed scripts, unable to generalize beyond narrowly engineered environments. The central bottleneck is no longer hardware availability but the ability to learn robust policies for unstructured, contact-rich tasks~\cite{levineLearningHandEyeCoordination2016, kalashnikovQTOptScalable2018}.

What distinguishes manipulation from other embodied tasks such as navigation is its \emph{physical-interaction nature}: it demands joint reasoning over high-DoF kinematics, contact dynamics, force exchange, and fine-grained object perception. These characteristics make manipulation the defining capability of general-purpose embodied agents---the layer where perception and reasoning are ultimately grounded in physical contact~\cite{khatibInertialPropertiesRobot1987, lakeBuildingMachinesThat2017}.

Recent advances in vision-language foundation models, world models, and diffusion-based policies have reshaped the field. Research focus has shifted from isolated module design (controller, policy network, action representation) toward understanding how hardware, data, and learning paradigms interact as a system. Three persistent challenges drive this shift: (i) the high cost and distribution bias of embodied interaction data; (ii) the sim-to-real gap that degrades transferred policies; and (iii) the lack of unified frameworks connecting imitation learning, vision-language-action models, and world models across heterogeneous embodiments.

A growing body of surveys has addressed individual facets: Bai et al.\ revisit the high-level/low-level division of manipulation~\cite{baiUnifiedUnderstandingRobot2025}; Feng et al.\ and Liu et al.\ survey LLM- and WM-driven embodied systems~\cite{feng2025embodiedaillmsworld,liu2025aligningcyberspacephysical,liuEmbodiedIntelligenceSynergy2025}; multiple works investigate VLA architectures and deployment~\cite{kawaharazukaVisionLanguageActionModelsRobotics2025,maSurveyVisionLanguageActionModels2025,zhongSurveyVisionLanguageActionModels2025,guanEfficientVisionLanguageActionModels2025,xiangParallelsVLAModel2025}; Long et al.\ examine simulators and world models~\cite{longSurveyLearningEmbodied2025}; Ruan et al.\ and Jaquier et al.\ focus on multi-sensor fusion and transfer learning~\cite{ruanSurveyMultisensorFusion2025,jaquierTransferLearningRobotics2024}. However, most existing surveys are organized around model categories or functional modules, leaving the systemic interplay among hardware constraints, data pipelines, and learning paradigms underexplored.

\textbf{Objectives and scope of this survey:}
This survey focuses on \emph{learned manipulation intelligence} for robotic arms and their extensions to mobile and humanoid platforms. We scope the analysis to systems where perception, decision-making, and motor control are jointly shaped by data-driven methods---spanning imitation learning, reinforcement learning, vision-language-action models, and world models. Pure industrial automation (fixed scripts, teach pendants) and non-manipulation embodied tasks (autonomous driving, drone navigation) fall outside our scope.

\textbf{Contributions.}
We organize the field around the chain \emph{embodiment $\to$ data $\to$ learning $\to$ evaluation} (Figure~\ref{fig:survey_structure}), making explicit how hardware constraints propagate through data design into policy capability and deployment feasibility; the three-stage deployment lifecycle (\S4.4) and the four-dimensional evaluation axes (\S\ref{sec:evaluation}) emerge as integral components rather than isolated appendices.
We classify embodied datasets into skill-level, task-level, and foundation-level tiers via a recursive knowledge-state model ($\mathcal{K}_0 \to \mathcal{K}_3$, \S\ref{sec:eai_datasets}) that quantifies the uncertainty each tier reduces.
Rather than treating Vision-Language-Action models as a single category, we decompose manipulation intelligence into cognitive planning (five paradigms, \S\ref{sec:embodied_brain}) and policy learning (four paradigms, \S\ref{sec:policy_learning_low_level}), then propose a Type~I/II/III integration taxonomy (\S4.3) that captures the 2023--2026 structural evolution from RT-2 through $\pi_0$ to GR00T~N1.
Finally, we ground the open challenges (\S\ref{sec:challenges}) in the specific gaps identified across all preceding sections, providing actionable research directions rather than generic desiderata.

\textbf{Structure of the Paper:}
As illustrated in Figure~\ref{fig:survey_structure}, the paper is structured around four technical pillars:

\begin{itemize}[nosep,leftmargin=*]
  \item \textbf{System Architecture} (\S\ref{sec:embodiment_hardware}) examines the physical embodiment---end effectors, stationary platforms, and mobile platforms---together with the software stack (middleware, motion planning, real-time control) that closes the perception-to-actuation loop.
  \item \textbf{Data Ecosystem} (\S3) traces how interaction experience is collected, organized, and scaled: from skill-level and task-level datasets through teleoperation and simulation pipelines to foundation-level cross-embodiment corpora.
  \item \textbf{Learning Paradigms} (\S4) covers the full spectrum from cognitive planning through policy learning to integration architectures and deployment strategies that bridge high-level reasoning with low-level motor execution.
  \item \textbf{Evaluation} (\S\ref{sec:evaluation}) synthesizes emerging metrics across task execution quality, data efficiency, generalization, and energy consumption.
\end{itemize}

Section~\ref{sec:challenges} concludes with open challenges and future directions.

\begin{figure*}[t]
  \centering
  \framebox[\textwidth]{\rule{0pt}{6cm}} % placeholder for survey structure figure
  \caption{Organization of this survey. The four technical pillars---System Architecture, Data Ecosystem, Learning Paradigms, and Evaluation---are connected by the flow from physical embodiment through data to learned policies and their real-world assessment.}\label{fig:survey_structure}
\end{figure*}

\section{System Architecture of Robotic Arms}\label{sec:embodiment_hardware}

The realizable action space of a robotic system is bounded by its electromechanical architecture: the end effector's contact geometry, the manipulator's kinematic redundancy, and the base's mobility jointly determine what tasks a policy can physically execute. This section surveys the hardware landscape along three axes---end effectors, stationary platforms, and mobile platforms---then describes the software stack (middleware, motion planning, and real-time control) that closes the loop from perception to actuation.

\subsection{Physical Embodiment}

The physical embodiment of a manipulation system spans three levels of abstraction (Fig.~\ref{fig:embodiment_taxonomy}): the \emph{end effector} that mediates contact, the \emph{stationary platform} that provides kinematic reach and torque bandwidth, and the \emph{mobile platform} that extends the workspace to unstructured environments.

% TODO: replace with actual figure
\begin{figure}[t]
\centering
% \includegraphics[width=\linewidth]{figures/embodiment_taxonomy.pdf}
\caption{Taxonomy of physical embodiments for robotic manipulation. End effectors are classified by contact modality (jaw gripper, dexterous hand, soft hand) and sensing (vision, tactile); platforms are divided into stationary (single-arm, bimanual) and mobile (wheeled, quadrupedal, humanoid).}\label{fig:embodiment_taxonomy}
\end{figure}

\subsubsection{End Effectors}\label{subsec:end_effectors}

End effectors define the contact geometry, force resolution, and sensory feedback available to the policy. Jaw grippers such as the Robotiq 2F/3F series use mechanical linkages for adaptive form closure, trading in-hand dexterity for grasp robustness. Dexterous hands---Shadow Hand (24 DoF), Allegro Hand (16 DoF), and the D'Claw benchmark---approximate human-level kinematics but require solving high-rank inverse kinematics under non-convex friction cone constraints. Soft hands (RBO Hand 3~\cite{Puhlmann_2022}, Festo BionicSoftHand, SpiRobs) exploit pneumatic compliance to bypass explicit force control, at the cost of state-estimation difficulty typically addressed by FEM or data-driven dynamics models. On the sensing side, eye-in-hand cameras (e.g., RealSense) resolve pre-contact geometry, while tactile sensors such as GelSight~\cite{zhao2023gelsightsveltehumanfingershaped} and DIGIT convert contact geometry into images processable by standard vision backbones, enabling visuo-tactile closed-loop control.

\subsubsection{Stationary Platforms}\label{subsec:manipulators}

Fixed-base manipulators offer the highest control bandwidth and repeatability for contact-rich research. Single-arm systems such as the Franka Panda and KUKA iiwa provide joint-side torque sensing for 1\,kHz impedance control; the Kinova Gen3 adds 7-DoF continuous-rotation redundancy for null-space optimization. Bimanual setups (e.g., dual Franka Pandas, ABB YuMi) introduce closed-chain constraints requiring explicit internal-force coordination. The ALOHA system~\cite{zhao2023learning} trades torque-loop fidelity for teleoperation throughput, enabling large-scale demonstration collection where learned policies compensate for mechanical backlash.

\subsubsection{Mobile Platforms}\label{subsec:mobile_platforms}

Mounting manipulators on mobile bases couples locomotion and manipulation into a single dynamical system. Wheeled mobile manipulators (PR2, TIAGo, Hello Robot Stretch~3) pair differential or holonomic bases with manipulation arms; controllers must reconcile non-holonomic base constraints with end-effector tracking accuracy. Quadrupedal manipulators (Spot, Unitree GO2/B1) add terrain-adaptive stability: the controller maintains the dynamic stability margin while the arm exerts interaction forces, requiring high-frequency re-planning of foot contacts. Humanoids (Atlas, Optimus Gen~2, Figure~02, 1X Neo, Unitree G1) combine bipedal balance with bimanual coordination under centroidal dynamics control, where manipulation forces must be balanced against whole-body angular momentum.

Across these embodiment families, expanding embodiment scope increases the reachable task space but also strengthens coupling between contact, motion, and stability, shifting the bottleneck from hardware capability to real-time coordination and policy execution.

\subsection{SoftWare Framework}

\subsubsection{Modern Robotics Middleware and Operating Systems}\label{subsec:control_middleware}

We structure the software stack into three tiers (Table~\ref{tab:framework_comparison}): the Kernel Layer (L0) for hard real-time motor control, the Transport Layer (L1) for inter-process data flow, and the System Layer (L2) for distributed orchestration.

At L0, commercial kernels such as VxWorks\footnote{\url{https://www.windriver.com/products/vxworks}} and QNX\footnote{\url{https://blackberry.qnx.com}} guarantee bounded interrupt latency through priority-based preemptive scheduling and microkernel fault isolation, respectively. For resource-constrained edge devices, lightweight RTOS options---FreeRTOS\footnote{\url{https://www.freertos.org/}}, Zephyr\footnote{\url{https://zephyrproject.org/}}, and the POSIX-compliant NuttX\footnote{\url{https://nuttx.apache.org/}}---minimize context-switching overhead while providing scalable device tree architectures.

At L1, the transition from fragmented middleware (Player~\cite{gerkey2001player}, YARP~\cite{metta2006yarp}, Orocos~\cite{bruyninckx2001open}) to ROS~1~\cite{quigley2009ros} and subsequently ROS~2~\cite{macenski2022robot}---which adopted decentralized DDS with granular QoS policies---established the current communication backbone.
As large embodied models demand high-bandwidth multimodal fusion, two strategies address the resulting serialization bottleneck: hardware-accelerated middleware such as RobotCore~\cite{mayoral2022robotcore} and NVIDIA Isaac ROS that bypass CPU serialization via FPGA/GPU pipelines, and zero-copy frameworks such as LCM~\cite{huang2010lcm}, Dora-rs~\cite{han2024dora} (Apache Arrow columnar format), Agnocast~\cite{ishikawa2025agnocast} (true zero-copy pub/sub for all ROS~2 message types).

At L2, hardware abstraction frameworks like Viam~\cite{viam2022} encapsulate actuators and sensors into language-agnostic RPC APIs, while declarative dataflow orchestrators such as Zenoh-Flow~\cite{zenoh2024} deploy execution DAGs across edge and cloud. For fleet-scale deployments, Open-RMF~\cite{openrmf2021} provides vendor-neutral traffic deconfliction and task dispatching, and FogROS2~\cite{ichnowski2022fogros2} extends the ROS~2 graph into cloud environments for dynamic compute offloading. More recently, AI-native systems such as RoboOS~\cite{wang2025roboos} propose a hierarchical architecture coupling multimodal LLMs for task planning with modular skill libraries for motor execution, enabling cross-embodiment adaptability across heterogeneous fleets.

\begin{table*}[t]
\centering
\caption{\textbf{Three-Tier Software Architecture for Embodied Intelligence.}
\textbf{L0:} Hard real-time kernels for microsecond-level motor control determinism.
\textbf{L1:} Transport middleware for high-throughput, low-latency inter-process data flow.
\textbf{L2:} System-layer orchestration for distributed deployment, fleet management, and AI-native coordination.
\textbf{Perf.:} Interrupt/IPC latency (L0/L1) or deployment granularity (L2).}\label{tab:framework_comparison}
\renewcommand{\arraystretch}{1.2}
\setlength{\tabcolsep}{3pt}
\scriptsize

\begin{tabular}{
  p{0.11\linewidth}
  p{0.09\linewidth}
  p{0.10\linewidth}
  p{0.34\linewidth}
  p{0.10\linewidth}
  p{0.12\linewidth}
}
\toprule
\textbf{Framework} &
\textbf{Type} &
\textbf{Core Mech.} &
\textbf{Key Contribution} &
\textbf{Perf.} &
\textbf{Scope} \\
\midrule

\multicolumn{6}{c}{\textbf{L0: Kernel Layer --- Hard Real-Time Execution}} \\
\midrule

\textbf{VxWorks} &
Monolithic RTOS &
Priority Preemptive &
Industry-standard deterministic kernel; bounded interrupt latency via priority-based preemptive scheduling for synchronous multi-axis servo control in aerospace and industrial automation. &
$<$10\,$\mu$s &
Industrial / Aerospace \\

\textbf{QNX} &
Microkernel &
Spatial Isolation &
Microkernel architecture isolating drivers and protocols into independent user-space memory segments; memory faults in peripherals cannot crash the core scheduling loop. POSIX-certified. &
$<$5\,$\mu$s &
Safety-critical / Automotive \\

\textbf{Zephyr} &
Modular RTOS &
Device Tree + SMP &
Linux Foundation governed; scalable device tree architecture with native SMP support; enforces strict memory domain protection even on microcontrollers without hardware MMUs. &
$<$20\,$\mu$s &
MCU / Robotic Edge \\

\midrule
\multicolumn{6}{c}{\textbf{L1: Transport Layer --- High-Throughput Data Flow}} \\
\midrule

\textbf{ROS~2}~\cite{macenski2022robot} &
Decentralized &
DDS (RTPS) &
Replaced ROS~1's centralized Master with decentralized DDS\@; granular QoS policies for reliability, deadline, and liveliness. Current \emph{de facto} communication backbone for the robotics community. &
1--10\,ms &
General Robotics \\

\textbf{LCM}~\cite{huang2010lcm} &
Masterless &
UDP Multicast &
Lightweight marshalling with minimal serialization overhead; masterless UDP multicast eliminates broker bottleneck for high-frequency legged locomotion control loops ($>$1\,kHz). &
$<$1\,ms &
Legged Robots \\

\textbf{Dora-rs}~\cite{han2024dora} &
Dataflow &
Apache Arrow &
Rust-based dataflow framework integrating Apache Arrow columnar format for zero-copy shared memory; eliminates serialization jitter between Python deep learning inference and Rust execution loops. &
$<$50\,$\mu$s (IPC) &
Low-latency AI \\

\textbf{Agnocast}~\cite{ishikawa2025agnocast} &
Zero-Copy IPC &
Shared Memory &
True zero-copy pub/sub for \emph{all} ROS~2 message types including dynamically-sized payloads (PointCloud2, Image); rclcpp-compatible. 25\% worst-case latency reduction on Autoware pipelines. &
Const.\ overhead &
ROS~2 Ecosystem \\

\midrule
\multicolumn{6}{c}{\textbf{L2: System Layer --- Distributed Orchestration \& AI-Native OS}} \\
\midrule

\textbf{Zenoh}~\cite{zenoh2024} &
DAG Orchestration &
Zenoh Protocol &
Declarative dataflow framework defining robotic pipelines as DAGs; automatically deploys and routes execution nodes across MCUs, edge GPUs, and cloud servers via the Zenoh protocol. &
Edge-to-Cloud &
Distributed Systems \\

\textbf{FogROS2}~\cite{ichnowski2022fogros2} &
Cloud Extension &
ROS~2 Graph &
Natively extends ROS~2 computational graph into cloud environments; robots offload latency-tolerant semantic reasoning to cloud while maintaining high-frequency reactive control on local edge. &
Hybrid &
Robot-as-a-Service \\

\textbf{RoboOS}~\cite{wang2025roboos} &
Brain--Cerebellum &
MLLM + Skill Lib &
Hierarchical architecture: MLLM-based embodied brain for global perception, task planning, and multi-agent scheduling; modular cerebellum skill library for low-level motor execution. Cross-embodiment adaptability. &
Agentic Logic &
Multi-Agent \\

\bottomrule
\end{tabular}
\end{table*}

\subsubsection{Low-level Motion Planning and Control}

While middleware provides the communication infrastructure, the low-level motion planning and control layer is the execution engine of the system. This layer is responsible for translating trajectory waypoints into high-frequency motor torques, ensuring that the physical execution adheres to the limits of kinematics and dynamics while maintaining stability during contact-rich interactions.

\textbf{Kinematic and Dynamic Foundations} The primary task of this layer is to resolve the mapping between the operational space and the joint space. For a manipulator with $n$ degrees-of-freedom, the differential kinematics are governed by the relationship:
\begin{equation}
    \dot{\mathbf{x}} = \mathbf{J}(\mathbf{q})\dot{\mathbf{q}}
\end{equation}
where $\mathbf{x} \in \mathbb{R}^6$ is the end-effector pose, $\mathbf{q} \in \mathbb{R}^n$ is the joint configuration, and $\mathbf{J}(\mathbf{q})$ is the analytical Jacobian matrix. To achieve precise tracking and disturbance rejection, the controller must account for the full rigid-body dynamics, typically expressed in the joint space as:
\begin{equation}
    \mathbf{M}(\mathbf{q})\ddot{\mathbf{q}} + \mathbf{C}(\mathbf{q}, \dot{\mathbf{q}})\dot{\mathbf{q}} + \mathbf{g}(\mathbf{q}) = \bm{\tau} - \mathbf{J}^T(\mathbf{q})\mathbf{F}_{ext}
\end{equation}
where $\mathbf{M}(\mathbf{q})$ is the inertia matrix, $\mathbf{C}(\mathbf{q}, \dot{\mathbf{q}})$ represents the Coriolis and centrifugal forces, $\mathbf{g}(\mathbf{q})$ is the gravity vector, $\bm{\tau}$ is the vector of joint torques, and $\mathbf{F}_{ext}$ represents external interaction forces~\cite{sicilianoSpringerHandbookRobotics2016}.

\textbf{Optimization-based Planning and Control} Modern embodied systems increasingly favor \textbf{Model Predictive Control (MPC)} and \textbf{Whole-Body Control (WBC)} frameworks to handle high-dimensional constraints. At each control cycle (typically $1$ ms), the system solves a \textbf{Quadratic Programming (QP)} problem to find the optimal control input $\mathbf{u}$ (e.g., accelerations or torques):
\begin{equation}
    \begin{aligned}
        \min_{\mathbf{u}} \quad & \|\mathbf{A}\mathbf{u} - \mathbf{b}\|^2 + \lambda \|\mathbf{u}\|^2 \\
        \text{s.t.} \quad & \mathbf{C}_{ineq}\mathbf{u} \leq \mathbf{d}_{ineq}, \quad \mathbf{C}_{eq}\mathbf{u} = \mathbf{d}_{eq}
    \end{aligned}
\end{equation}
where the constraints account for joint limits, collision avoidance, and the non-linear dynamics defined in Eq. (2). This optimization-based approach allows the robot to reactively adjust its motion in response to environmental disturbances while strictly respecting the physical boundaries of the hardware~\cite{schulman2014optimizing}.

\textbf{Impedance Control for Physical Interaction} Unlike pure navigation, robotic manipulation necessitates intentional contact with the environment. To ensure safety and stability during such contact-rich tasks, \textbf{Impedance Control}~\cite{hogan1985impedance} is employed to regulate the dynamic relationship between the robot and its environment. In the task space, the desired behavior is modeled as a mass-spring-damper system:
\begin{equation}
    \mathbf{M}_d \Delta \ddot{\mathbf{x}} + \mathbf{D}_d \Delta \dot{\mathbf{x}} + \mathbf{K}_d \Delta \mathbf{x} = \mathbf{F}_{ext}
\end{equation}
where $\Delta \mathbf{x} = \mathbf{x} - \mathbf{x}_{ref}$ is the tracking error, and $\mathbf{M}_d, \mathbf{D}_d, \mathbf{K}_d$ are the desired inertia, damping, and stiffness matrices, respectively. By modulating these parameters, the controller can exhibit compliant behavior, essential for tasks requiring a ``soft touch'' such as assembly or human-robot collaboration.

\textbf{Temporal Determinism and RTOS} The performance of these numerical solvers is fundamentally bound by the temporal determinism of the execution environment. High-performance controllers are typically deployed on \textbf{Real-Time Operating Systems (RTOS)} such as \textbf{Preempt-RT} or \textbf{QNX}. This infrastructure guarantees that the entire control loop—including the QP solver and dynamic compensation—completes within a bounded latency, preventing the oscillatory instabilities often caused by timing jitter in high-frequency feedback loops.

Together, the hardware embodiment and software stack define the \emph{action envelope} within which any learned policy must operate. The next question is what data is needed to fill that envelope with intelligent behavior---and how such data should be organized, collected, and scaled.

\section{Data Ecosystem of Robotic Arms}\label{sec:embodiment_dataset}

The system architecture described in \S\ref{sec:embodiment_hardware} determines what a robot \emph{can} do; the data ecosystem determines what it \emph{learns} to do. End-effector morphology constrains grasp geometry, platform kinematics bound the reachable workspace, and middleware latency caps the control frequency---all of which propagate directly into dataset design choices: what modalities to record, what action spaces to annotate, and what diversity to target. This section traces the data pipeline from dataset organization (\S\ref{sec:eai_datasets}) through real-world and simulated acquisition (\S\ref{sec:data_acquisition}) to augmentation (\S\ref{sec:data_augmentation}) and industrial-scale data pipelines (\S\ref{sec:data_factories_robotic_arms}).

\subsection{Embodied Datasets}\label{sec:eai_datasets}

Embodied datasets are the empirical substrate upon which all manipulation intelligence is built. Their design is fundamentally shaped by a single question: \emph{what must the downstream model learn from this data?} We organize the landscape into three tiers based on the nature of knowledge each dataset encodes (see Fig.~\ref{fig:evolution} and Tables~\ref{tab:grasping_datasets}--\ref{tab:foundation_datasets}).

\textbf{Skill-level datasets} provide dense, high-resolution sampling of local interaction manifolds---grasping geometry, contact forces, material compliance---from which models learn how individual actions succeed or fail. \textbf{Task-level datasets} bind these primitives into instruction-aligned temporal sequences spanning diverse environments, from which models learn how tasks are composed and completed. \textbf{Foundation-level datasets} unify experience across heterogeneous morphologies and sensing modalities under standardized protocols, from which foundation models distill the invariant physical laws governing manipulation at scale.

This three-tier progression can be understood through an information-theoretic lens: each tier targets a distinct source of uncertainty about the physical world, and the value of a dataset is measured by how much of that uncertainty it eliminates. We formalize this via a recursive \emph{knowledge state} $\mathcal{K}_i$ that accumulates the information provided by all preceding tiers:
\begin{equation}\label{eq:knowledge_state}
\mathcal{K}_0 = \varnothing, \qquad \mathcal{K}_i = \mathcal{K}_{i-1} \oplus \mathcal{D}_i
\end{equation}
where $\oplus$ denotes knowledge augmentation. The total information gain decomposes by the chain rule into three nested layers, each conditioned on the knowledge already established:
\begin{equation}\label{eq:chain_rule}
I_{\text{total}} = \underbrace{I_{\text{skill}}}_{\text{geometric}} + \underbrace{I_{\text{task}}}_{\text{semantic}} + \underbrace{I_{\text{found}}}_{\text{transfer}}
\end{equation}

\begin{enumerate}[leftmargin=*, nosep]
\item \textbf{Skill-level Datasets} reduce \emph{geometric entropy}---the uncertainty over where and how physical interaction succeeds, starting from a blank slate:
\begin{equation}\label{eq:skill_info}
I_{\text{skill}} = H(\mathbf{a}_{\text{prim}}, \mathbf{o}_{\text{contact}} \mid \mathcal{K}_0) - H(\mathbf{a}_{\text{prim}}, \mathbf{o}_{\text{contact}} \mid \mathcal{K}_1)
\end{equation}
where $\mathbf{a}_{\text{prim}}$ denotes atomic motor primitives (grasping, pushing, folding, assembly) and $\mathbf{o}_{\text{contact}}$ the associated contact observations (force, tactile, geometry). In practice, each real-world trajectory demands physical hardware time, human supervision, and post-hoc annotation---a per-trajectory cost that does not diminish with scale, causing $I_{\text{skill}}$ to saturate well below the theoretical maximum.

\item \textbf{Task-level Datasets}, building on the geometric knowledge $\mathcal{K}_1$ already acquired, reduce the residual \emph{semantic entropy}---the uncertainty over how primitives compose into goal-directed behaviors:
\begin{equation}\label{eq:task_info}
I_{\text{task}} = H(\tau, l \mid \mathcal{K}_1) - H(\tau, l \mid \mathcal{K}_2)
\end{equation}
where $\tau$ denotes trajectories and $l$ language instructions. Yet heterogeneous annotation schemas, action space definitions, and observation configurations across labs and platforms prevent direct data mixing, discounting the effective information gain by format compatibility $\alpha$ and observation-action alignment $\beta$: $I_{\text{task}}^{\text{eff}} = \alpha \cdot \beta \cdot I_{\text{task}}$.

\item \textbf{Foundation-level Datasets}, given the accumulated knowledge $\mathcal{K}_2$ from both geometric and semantic tiers, reduce \emph{transfer entropy}---the uncertainty over whether learned behaviors generalize across heterogeneous embodiments and environments:
\begin{equation}\label{eq:found_info}
I_{\text{found}} = H(s' \mid s, a, e, \mathcal{K}_2) - H(s' \mid s, a, e, \mathcal{K}_3)
\end{equation}
where $e$ indexes embodiment type and $s' \mid s, a$ captures transition dynamics. The fundamental challenge here is measurement itself: no universally accepted metric exists to quantify whether cross-embodiment generalization has truly occurred, making the actual information contribution of each dataset to downstream capability difficult to isolate and compare.
\end{enumerate}

\subsubsection{Skill-level Datasets}\label{sec:skill_datasets_narrative}

Skill-level datasets provide dense sampling of atomic motor primitives---grasping, pushing, folding, assembly---that define the boundaries of physical interaction (Table~\ref{tab:grasping_datasets}).

Early benchmarks such as \textbf{Cornell}~\cite{lenz2014deep} and \textbf{Jacquard}~\cite{depierre2018jacquard} represented graspable regions as 2D oriented rectangles, limiting interaction to top-down approaches. \textbf{Dex-Net 2.0}~\cite{mahler2017dex} and \textbf{ACRONYM}~\cite{eppner2020acronym} addressed data scarcity through analytical synthesis based on force-closure metrics, generating millions of grasp labels. This line has scaled to billion-level real-world benchmarks: \textbf{GraspNet-1B}~\cite{fang2020graspnet} and \textbf{GraspClutter6D}~\cite{back2025graspclutter6d} provide up to 9.3 billion annotated 6-DoF grasps in cluttered scenes.

Multi-fingered systems with $>$20 DoF require fundamentally different data. \textbf{DexGraspNet}~\cite{wang2023dexgraspnetlargescaleroboticdexterous} and \textbf{DexGraspNet 2.0}~\cite{zhang2024dexgraspnet20learninggenerative} (427M grasps in cluttered multi-object scenes, 90.7\% real-world success via diffusion-based synthesis), together with the billion-scale \textbf{Dex1B}~\cite{ye2025dex1b}, target sensorimotor alignment under heavy visual occlusion. Tactile modalities are critical at this level: \textbf{T-DEX}~\cite{guzey2023tdex} shows that self-supervised pre-training on 2.5 hours of unstructured tactile play data outperforms vision-only baselines by 1.7$\times$ on five dexterous tasks, enabling slip detection and in-hand reorientation that are unobservable through vision alone.

Beyond grasping, contact-rich and non-rigid dynamics demand dedicated resources (Table~\ref{tab:grasping_datasets}, Diverse Motor Primitives). \textbf{Omnipush}~\cite{bauza2021omnipush} quantifies planar pushing non-linearities across 250 objects. \textbf{PartNet-Mobility}~\cite{xiang2020sapien} and \textbf{FlowBot3D}~\cite{flowbot3d_2022} provide part-level interaction data for articulated object manipulation. Deformable objects require mesh-level fidelity: \textbf{UniFolding}~\cite{xue2023unifolding} captures local deformation laws for garment folding. For precision assembly, \textbf{FurnitureBench}~\cite{heo2023furniturebench} provides 5,000+ teleoperated trajectories (200+ hours) with a paired MuJoCo simulator, and \textbf{ToolFlowNet}~\cite{seita2022toolflownetroboticmanipulationtools} achieves 82\% real-world success on tool-use tasks by predicting dense per-point flow from point clouds.

A recent direction integrates natural language with geometric affordance (Table~\ref{tab:grasping_datasets}, Language-Grounded Grasping). \textbf{OCID-VLG}~\cite{tziafas2023ocidvlg} pairs cluttered scenes with referring expressions and grasp poses. \textbf{Grasp-Anything}~\cite{vuong2023graspanything} scales this to 1M images with 600M grasp annotations via foundation model synthesis. \textbf{ReasoningGrasp}~\cite{jin2024reasoninggrasp} goes further by requiring models to infer \emph{why} and \emph{what} to grasp from implicit intent, rather than following explicit object references.


\begin{table*}[t]
\centering
\caption{\textbf{Chronological Overview of Skill-level Datasets.} Evolution reflects the shift from planar grasping rectangles to dense 6-DoF poses, dexterous hands, language-grounded reasoning, and diverse motor primitives (pushing, articulation, deformable manipulation, assembly). Note the scarcity of real-world data for non-grasping primitives.}\label{tab:grasping_datasets}
\resizebox{\textwidth}{!}{%
\begin{tabular}{l c c c c c c c c}
\toprule
\textbf{Dataset} & \textbf{Year} & \textbf{Skill Type} & \textbf{Scene Complexity} & \textbf{Objects} & \textbf{Domain} & \textbf{Scale (Size)} & \textbf{Modality} & \textbf{Lang.} \\
\midrule
\multicolumn{9}{c}{\textbf{Early Geometric Grasping}} \\
\midrule
\textbf{Cornell}~\cite{lenz2014deep} & 2011 & Rect.\ Grasp (2D) & Single Object & 240 & Real & 1k imgs & RGB-D & $\times$ \\
\textbf{Dex-Net 2.0}~\cite{mahler2017dex} & 2017 & Parallel-Jaw Grasp & Single Object & 1.5k & Sim & 6.7M grasps & Depth & $\times$ \\
\textbf{Jacquard}~\cite{depierre2018jacquard} & 2018 & Rect.\ Grasp (2D) & Single Object & 11k & Sim & 54k imgs (1.1M grasps) & RGB-D & $\times$ \\
\midrule
\multicolumn{9}{c}{\textbf{Dense 6-DoF Grasping}} \\
\midrule
\textbf{GraspNet-1B}~\cite{fang2020graspnet} & 2020 & 6-DoF Grasp & Heavy Clutter & 88 & Real & 97k imgs (1.2B grasps) & RGB-D & $\times$ \\
\textbf{ACRONYM}~\cite{eppner2020acronym} & 2021 & 6-DoF Grasp & Multi-Object & 8,872 & Sim & 17.7M grasps & RGB-D & $\times$ \\
\textbf{GraspClutter6D}~\cite{back2025graspclutter6d} & 2025 & 6-DoF Grasp & Cluttered & 200 & Real & 52k imgs (9.3B grasps) & RGB-D & $\times$ \\
\midrule
\multicolumn{9}{c}{\textbf{Dexterous Grasping}} \\
\midrule
\textbf{DexGraspNet}~\cite{wang2023dexgraspnetlargescaleroboticdexterous} & 2023 & Dexterous Grasp & Single Object & 5,355 & Sim & 1.32M grasps & Point Cloud & $\times$ \\
\textbf{DexGraspNet 2.0}~\cite{zhang2024dexgraspnet20learninggenerative} & 2024 & Dexterous Grasp & Multi-Object & 5,355 & Sim & 427M grasps & Point Cloud & $\times$ \\
\textbf{Dex1B}~\cite{ye2025dex1b} & 2025 & Dexterous Grasp & Single Object & 1M & Sim & 1B grasps & Point Cloud & $\times$ \\
\midrule
\multicolumn{9}{c}{\textbf{Language-Grounded Grasping}} \\
\midrule
\textbf{OCID-VLG}~\cite{tziafas2023ocidvlg} & 2023 & Rect.\ Grasp & Multi-Object & 89 & Real & 75k grasps & RGB-D & $\checkmark$ \\
\textbf{Grasp-Anything}~\cite{vuong2023graspanything} & 2023 & Rect.\ Grasp + Push & Multi-Object & 3M & Syn & 1M imgs (600M grasps) & RGB & $\checkmark$ \\
\textbf{ReasoningGrasp}~\cite{jin2024reasoninggrasp} & 2024 & 6-DoF Grasp & Multi-Object & 64 & Real & 99M grasps & RGB-D & $\checkmark$ \\
\midrule
\multicolumn{9}{c}{\textbf{Diverse Motor Primitives (Non-Grasping)}} \\
\midrule
\textbf{Omnipush}~\cite{bauza2021omnipush} & 2019 & Pushing (Planar) & Single Object & 250 & Real & 62.5k pushes & RGB-D & $\times$ \\
\textbf{PartNet-Mobility}~\cite{xiang2020sapien} & 2020 & Articulated (Open/Close) & Single Object & 2,346 models & Sim & 2,346 articulated assets & 3D / PC & $\times$ \\
\textbf{FlowBot3D}~\cite{flowbot3d_2022} & 2022 & Articulated (Flow) & Cluttered & 1k+ & Sim & Per-point 3D flow fields & 3D / PC & $\times$ \\
\textbf{FurnitureBench}~\cite{heo2023furniturebench} & 2023 & Assembly (Precise) & Tabletop & Furniture & Sim+Real & 5.1k demos (220 hrs) & RGB & $\times$ \\
\textbf{UniFolding}~\cite{xue2023unifolding} & 2023 & Deformable (Folding) & Multi-Object & Garments & Sim+Real & Sample-efficient ($<$1k) & RGB-D & $\times$ \\
\bottomrule
\end{tabular}
}
\end{table*}

\subsubsection{Task-level Datasets}\label{sec:task_datasets_narrative}

Task-level datasets bind atomic motor primitives into temporally extended, instruction-aligned trajectories. Where skill-level data answers \emph{how} a single action succeeds, task-level data answers \emph{what sequence of actions} completes a goal. Table~\ref{tab:real_world_tasks} summarizes representative collections.

Early fleet-scale collection by \textbf{MT-Opt}~\cite{kalashnikov2021mt} provided initial evidence that data volume directly improves multi-task success rates, but these trajectories lacked language grounding---actions were indexed by task IDs rather than natural language.

The introduction of language-conditioned trajectories marked a qualitative shift. \textbf{BC-Z}~\cite{jang2022bc} paired 26k trajectories with 100 semantic verbs, and the \textbf{RT-1 Dataset}~\cite{brohan2023rt1} scaled to 130k trajectories spanning 700+ instructions on the Everyday Robot platform. These datasets enabled the first instruction-following policies that generalize across task descriptions. Yet both were collected in 1--2 institutional kitchens, limiting environmental diversity.

Subsequent work addressed this scene bottleneck through distributed collection. \textbf{BridgeData V2}~\cite{walke2024bridgedatav2} aggregated 60k trajectories across 24 scenes via multi-site collection. \textbf{Dobb-E}~\cite{shafiullah2023bringing} reduced hardware requirements to an iPhone-mounted Stretch robot, accessing 22 domestic environments. \textbf{DROID}~\cite{khazatsky2024droid} coordinated 18 institutions to collect 76k trajectories in 564 distinct scenes with calibrated extrinsics, representing the most environmentally diverse single-embodiment dataset to date. Complementing visual diversity, \textbf{FMB}~\cite{fu2025fmb} emphasized physical grounding through rigorous force-torque calibration across 22k contact-rich assembly trajectories. More recently, the \textbf{RoboMIND} series~\cite{wu2025robomind, hou2026robomind20} contributes a diagnostic dimension: \textbf{RoboMIND}~(107k trajectories, 4 robots) incorporates 5,000 failure demonstrations annotated with explicit error causes, and \textbf{RoboMIND 2.0}~(310k trajectories, 6 robots, 739 tasks) further scales cross-embodiment coverage, enabling models to learn recovery behaviors rather than only successful executions.

\begin{table*}[t] % 使用 table* 横跨双栏
\centering
\caption{Taxonomy of Task-level Embodied Datasets: Real-world Physical Trajectories.}\label{tab:real_world_tasks}
\footnotesize 
\setlength{\tabcolsep}{2.5pt} 
\resizebox{\textwidth}{!}{ 
\begin{tabular}{lcccccccc}
\toprule
\textbf{Dataset} & \textbf{Year} & \textbf{\# Traj.} & \textbf{Verbs} & \textbf{Scenes} & \textbf{Lang.} & \textbf{Calib.} & \textbf{Robot} & \textbf{Collection} \\
\midrule
MIME~\cite{sharma2018mime} & 2018 & 8.3k & 20 & 1 & $\times$ & $\times$ & Baxter & Teleop. \\
RoboTurk~\cite{mandlekar2018roboturk} & 2018 & 2.1k & 2 & 1 & $\times$ & $\times$ & Panda & Teleop. \\
MT-Opt~\cite{kalashnikov2021mt} & 2021 & 800k & 2 & 1 & $\times$ & $\times$ & Google Arm & Mixed \\
BC-Z~\cite{jang2022bc} & 2021 & 26k & 100 & 1 & $\checkmark$ & $\times$ & EDR & Teleop. \\
BridgeData V1~\cite{ebert2021bridgedata} & 2022 & 7.2k & 4 & 12 & $\checkmark$ & $\times$ & WidowX & Teleop. \\
RT-1 Dataset~\cite{brohan2023rt1} & 2022 & 130k & 700+ & 2 & $\checkmark$ & $\times$ & EDR & Teleop. \\
RoboSet~\cite{bharadhwaj2023roboagent} & 2023 & 98.5k & 12 & 11 & $\checkmark$ & $\times$ & Panda & H/S Hybrid \\
Dobb-E~\cite{shafiullah2023bringing} & 2023 & 1.5k & 109 & 22 & $\times$ & $\times$ & Stretch & iPhone \\
BridgeData V2~\cite{walke2024bridgedatav2} & 2024 & 60k & 82 & 24 & $\checkmark$ & $\times$ & WidowX & H/S Hybrid \\
DROID~\cite{khazatsky2024droid} & 2024 & 76k & 50+ & 564 & $\checkmark$ & $\checkmark$ & Panda & Teleop. \\
FMB~\cite{fu2025fmb} & 2025 & 22k & 10 & 1 & $\checkmark$ & $\checkmark$ & Panda & Mixed \\
RoboMIND~\cite{wu2025robomind} & 2025 & 107k & 479 & \- & $\checkmark$ & $\times$ & 4 types & Teleop. \\
RoboMIND 2.0~\cite{hou2026robomind20} & 2026 & 310k & 739 & \- & $\checkmark$ & $\checkmark$ & 6 types & Teleop. \\
\bottomrule
\end{tabular}
}
\end{table*}

\subsubsection{Foundation-level Datasets}\label{sec:foundation_datasets}

Foundation-level datasets unify experience across heterogeneous robot morphologies under standardized observation-action protocols (Table~\ref{tab:foundation_datasets}). Their defining goal is to enable transfer across embodiments---reducing the $H(s' \mid s, a, e, \mathcal{K}_2)$ term in Eq.~\ref{eq:found_info}---rather than optimizing performance on any single platform.

Cross-platform data standardization laid the groundwork. \textbf{RoboNet}~\cite{dasari2020robonet} initiated multi-institutional data sharing with 162k trajectories across seven robot platforms. \textbf{Open X-Embodiment} (OXE)~\cite{padalkar2023open} scaled this to over 1 million trajectories from 22 morphologies, standardized via the RLDS protocol, enabling the first cross-embodiment foundation model training experiments.

However, naively aggregating heterogeneous data introduces observation misalignment. \textbf{ARIO}~\cite{wang2024ario} addresses this with approximately 3 million episodes featuring synchronized multi-view timestamps, ensuring learned representations remain invariant to camera extrinsics. \textbf{RH-20T}~\cite{fang2023rh} extends the sensory scope beyond vision, integrating force-torque and audio signals across 110k sequences for contact-rich manipulation.

Recent efforts focus on scaling trajectory volume while maintaining physical fidelity. \textbf{AgiBot World Beta}~\cite{agibot2025} provides over one million trajectories ($\sim$43.8 TB) collected from 100+ humanoid robots for bimanual coordination and 6-DoF dexterous manipulation. Yet much of the highest-fidelity real-world data now remains proprietary, as industrial labs such as Physical Intelligence~\cite{black2024pi0} keep their large-scale datasets closed-source. To maintain parity, the open-source community is pivoting towards scalable simulation and low-cost egocentric gathering. \textbf{Ego-Exo4D}~\cite{grauman2024egoexo4d} exemplifies this trend, providing 740 hours of synchronized first- and third-person video across 8 activity domains, bridging abundant human demonstration footage with the exocentric viewpoints typical of robotic systems.

\begin{table*}[t!]
\centering
\caption{Foundation-level Datasets for Cross-Embodiment Manipulation.
\# Traj.\,=\,trajectory count;
\# Embod.\,=\,distinct robot platforms/configurations when explicitly reported;
Embod.\ Types: S\,=\,single-arm, B\,=\,bimanual, M\,=\,mobile manipulator, H\,=\,humanoid;
Modalities beyond RGB: D\,=\,depth, F\,=\,force/torque, A\,=\,audio;
Lang.\,=\,language supervision;
Action: EE\,=\,end-effector, Emb.\,=\,embodiment-specific, WB\,=\,whole-body;
Alignment: how cross-embodiment differences are reconciled.}
\label{tab:foundation_datasets}
\renewcommand{\arraystretch}{1.15}
\setlength{\tabcolsep}{2.5pt}
\scriptsize
\begin{tabular}{lccclclccll}
\toprule
\textbf{Dataset} & \textbf{Year} & \textbf{\# Traj.} & \textbf{\# Embod.} & \textbf{Types} & \textbf{Source} & \textbf{Modalities} & \textbf{Lang.} & \textbf{Action} & \textbf{Alignment} & \textbf{Schema} \\
\midrule
RoboNet~\cite{dasari2020robonet} & 2019 & 162k & 7 & S & Scripted & RGB & $\times$ & EE & Interface-level & Custom \\
RH-20T~\cite{fang2023rh} & 2023 & 110k & 7 & S & Teleop & RGB, D, F, A & $\checkmark$ & Emb. & Calibration-level & Custom \\
OXE~\cite{padalkar2023open} & 2023 & 1M+ & 22 & S, B, M & Aggregated & RGB (varied) & $\checkmark$ & EE & Schema-level & RLDS \\
ARIO~\cite{wang2024ario} & 2024 & 3M & -- & S, B, M & Aggregated (Real+Sim) & RGB, D & $\checkmark$ & Emb. & Temporal & Unified \\
AgiBot World~\cite{agibot2025} & 2025 & 1M+ & 100+ & B, H & Teleop & RGB, D, F & $\checkmark$ & WB & Pipeline-level & Custom \\
\bottomrule
\end{tabular}
\end{table*}

\subsection{Data Acquisition}\label{sec:data_acquisition}

\subsubsection{Real-world Data Gathering}\label{sec:real_acquisition}

Physical world data acquisition provides the indispensable grounding for robotic intelligence to navigate real-world stochasticity. This field has evolved along three orthogonal paradigms: the refinement of teleoperation interfaces, the scaling of decentralized crowdsourcing, and the emergence of autonomous, intervention-based data engines.

\textbf{High-Fidelity Teleoperation and Multi-Embodiment Alignment:} 
The primary challenge in real-world acquisition is the precise mapping of human expertise to robotic joint spaces. Early efforts like \textbf{MIME}~\cite{sharma2018mime} established the foundation for imitation learning, which was subsequently industrialized by the \textbf{ALOHA} ecosystem~\cite{zhao2023learning, fu2024mobile} through low-cost leader-follower teleoperation. To resolve the ``embodiment gap'' across diverse platforms, researchers introduced hardware-agnostic retargeting frameworks such as \textbf{AnyTeleop}~\cite{qin2024anyteleop} and \textbf{OPEN TEACH}~\cite{iyer2024openteach}. This era also witnessed a surge in specialized interfaces for high-dimensional coordination; \textbf{TeleMoMa}~\cite{dass2024telemoma} enables whole-body teleoperation across heterogeneous embodiments, while \textbf{Tilde}~\cite{si2024tilde} provides a kinematic-twin interface for finger-level dexterous hand learning and \textbf{ALPHA-BiACT}~\cite{alphabiact2024} integrates bilateral position-and-force control into a low-cost dual-arm platform for contact-aware demonstration. Wearable exoskeletons further reduce the barrier to in-the-wild data collection: \textbf{AirExo}~\cite{airexo2024} captures whole-arm trajectories with contact dynamics, and its successor \textbf{AirExo-2}~\cite{fang2025airexo2} introduces demonstration adaptors that convert human data into pseudo-robot experience via image inpainting and differentiable-rendering calibration, achieving policy transfer at \$600 hardware cost with teleoperated data.
Bridging teleoperation and full autonomy, \textbf{AutoRT}~\cite{gdm2024autort} deploys foundation models to orchestrate fleets of 20+ robots that autonomously propose and execute diverse data collection tasks in real environments, guided by a VLM-based ``Robot Constitution'' for safety; this pipeline collected 77k episodes spanning 6,650 unique task descriptions with minimal human oversight.

\textbf{Decentralized Acquisition and Perceptual Scaling:} 
To mitigate the constraints of laboratory-based data collection, research has increasingly pivoted toward decentralized acquisition strategies to maximize environmental entropy. Early frameworks such as \textbf{RoboTurk}~\cite{mandlekar2018roboturk} explored \textbf{remote web-based crowdsourcing}, utilizing latency-compensated interfaces to harvest demonstrations from a geographically distributed pool of operators. This approach transitioned toward \textbf{institutional distributed networks}, as exemplified by \textbf{DROID}~\cite{khazatsky2024droid}, which coordinates standardized collection across 18 global institutions to capture ``in-the-wild'' data in 564 distinct scenes. Parallel to institutional scaling, the field has seen a surge in \textbf{portable egocentric systems} that decouple the demonstration phase from the robot's physical presence. \textbf{Universal Manipulation Interface} (UMI)~\cite{chi2024umi} uses handheld grippers paired with consumer-grade cameras to harvest gripper-based datasets directly in domestic environments, while \textbf{DexCap}~\cite{wang2024dexcap} employs SLAM-based localization and electromagnetic tracking to capture high-dimensional dexterous hand motions for transfer to multi-fingered robotic systems. Taking this decoupling further, \textbf{M2R}~\cite{kim2024m2r} demonstrates that policies can be trained entirely from master-device demonstrations without any real robot involvement, enabling deep imitation learning via master-to-robot policy transfer. To ensure the utility of such heterogeneous data, alignment frameworks like \textbf{EgoMimic}~\cite{egomimic2024} and \textbf{AV-ALOHA}~\cite{chuang2025avaloha} focus on maintaining \textbf{perceptual congruency} between the human demonstrator's perspective and the robot's onboard sensors. 
Finally, minimizing hardware requirements to the limit, \textbf{RoboPocket}~\cite{fang2026robopocket} enables ubiquitous data collection by allowing users to fine-tune policies via smartphone AR corrections, effectively removing the need for physical robots or specialized rigs during the iteration phase. Pushing this minimalism further, \textbf{AoE (Always-on Egocentric)}~\cite{yang2026aoe} proposes a passive, always-on data collection paradigm using only smartphones worn by humans during daily activities, harvesting large-scale egocentric interaction videos without any teleoperation or robotic hardware---providing the statistical variance necessary for robust policy generalization at near-zero marginal cost. By systematically expanding both the operator pool and the spatial variety of scenes, these decentralized methodologies provide the statistical variance necessary for robust policy generalization.


\textbf{Interactive Refinement and Industrial Data Engines:} 
The most significant paradigm shift (2024--2026) is the transition from passive recording to \textbf{Active Data Engines} based on Human-in-the-loop (HITL) and Generalized HITL principles. Unlike traditional teleoperation, these pipelines utilize human effort only for \textbf{Interactive Intervention}, where operators correct autonomous rollouts only at critical failure points. This ``correction-as-data'' approach was pioneered by Google DeepMind's \textbf{RoboCat}~\cite{bousmalis2023robocat}, which established a self-improving loop for cross-task scaling. This paradigm has since been industrialized at an unprecedented scale: the \textbf{GEN-0} pipeline~\cite{generalist2025gen0} employs high-throughput manual intervention to generate over 270,000 hours of high-fidelity data, while the \textbf{AgiBot World} ecosystem~\cite{agibot2025} and \textbf{Tesla Optimus Data Pipeline} utilize massive humanoid fleets to refine policies through millions of real-world ``failure-to-success'' transitions. By integrating active learning and corrective feedback---further systematized by \textbf{IntervenGen}~\cite{pinto2024intervengen} and \textbf{CR-DAgger}~\cite{crdagger2024}---these industrial data engines transform the collection process into a dynamic pedagogical loop, effectively mitigating the ``physical hallucinations'' inherent in long-horizon task execution.


\subsubsection{Simulation Data Ecosystems}\label{sec:simulation_infrastructure}

Training data-hungry embodied models requires simulation at scale. We organize the current landscape around three ecosystems that, based on the benchmarks surveyed in this work, underpin the majority of recent manipulation research: \textbf{MuJoCo}, \textbf{SAPIEN}, and \textbf{NVIDIA Isaac Sim} (Table~\ref{tab:sim_classification}).

\textbf{MuJoCo} provides high-fidelity contact dynamics with analytic stability, serving as the backbone for benchmarks such as RoboCasa, LIBERO, Meta-World, and CALVIN. \textbf{SAPIEN} bridges simulation with large-scale articulated asset libraries (e.g., PartNet-Mobility), powering the ManiSkill series and enabling training across diverse unseen object categories. The \textbf{Isaac Ecosystem} leverages GPU-accelerated ray tracing and massive parallelism for photorealistic rendering, supporting Isaac Lab, DexGraspNet, and MimicGen among others. Legacy platforms (\textbf{PyBullet}, \textbf{CoppeliaSim}) and navigation-focused environments (\textbf{AI2-THOR}, \textbf{Habitat}) remain in active use for specialized benchmarks such as RLBench and Ravens.

\begin{table*}[t]
\centering
\caption{\textbf{Landscape of Simulation Infrastructures for large embodied models} We re-categorize platforms into the three dominant ecosystems driving modern manipulation research (MuJoCo, SAPIEN, Isaac) and specialized platforms serving as functional baselines. This taxonomy highlights the transition from task-specific simulators to comprehensive generative data factories.}\label{tab:sim_classification}
\resizebox{\textwidth}{!}{%
\begin{tabular}{@{}l p{4.2cm} p{5.8cm} p{5.5cm}@{}}
\toprule
\textbf{Simulator} & \textbf{Features} & \textbf{Capabilities \& Modalities} & \textbf{Key Benchmarks \& Datasets} \\
\midrule
\multicolumn{4}{c}{\textbf{\textbf{I. The Core Simulation Ecosystems}}} \\
\midrule
\addlinespace[0.5ex]
\textbf{MuJoCo} & Stable analytic contact solver; MJX hardware acceleration (TPU/GPU); Unified differentiable dynamics. & \textbf{Physics-Aligned Data}: Artifact-free trajectories for fine-grained motor control. 

\textbf{Modalities: Proprioception (Joint states), Contact forces, Kinematics, RGB.} & \textbf{RoboCasa}~\cite{nasiriany2024robocasa}, \textbf{LIBERO}~\cite{liu2023libero}, \textbf{RoboCerebra}~\cite{han2025robocerebra}, \textbf{BiGym}~\cite{chernyadev2024bigym}, \textbf{Meta-World}~\cite{yu2021meta}, \textbf{Robomimic} \\ \addlinespace

\textbf{Isaac Suite} & Massive GPU parallelism (PhysX5);RTX-enabled neural rendering; LLM-driven generative synthesis. & \textbf{Generative \& Scalable Data}: Billion-scale synthetic distributions and high-DoF control. 

\textbf{Modalities: RGB-D, LiDAR, Segmentation, Bounding boxes, Force/Torque, Event camera.} & \textbf{Isaac Lab}~\cite{mittal2023orbit}, \textbf{HumanoidSim}, \textbf{DexGraspNet}, \textbf{MimicGen}, \textbf{ExBody},  \textbf{FurnitureBench}~\cite{heo2023furniturebench}, \textbf{Arnold}~\cite{gong2023arnold} \\ \addlinespace

\textbf{SAPIEN} & High-fidelity articulated dynamics; Part-level mobility (PartNet); Parallelized Vulkan rendering. & \textbf{Semantic Interaction Data}: Interactions across articulated categories and part-level reasoning. 

\textbf{Modalities: RGB-D, Semantic/Instance masks, Contact forces, Articulated object states.} & \textbf{ManiSkill3}~\cite{tao2025maniskill3}, \textbf{SIMPLER}~\cite{li2025simpler}, \textbf{GAPartNet}, \textbf{RoboTwin}, \textbf{RoboFactory} \\ \addlinespace

\midrule

\multicolumn{4}{c}{\textbf{\textbf{II\@. Specialized Task Platforms}}} \\
\midrule
\addlinespace[0.5ex]
\textbf{PyBullet / CoppeliaSim} & Lightweight legacy physics; Stable IK solvers; Broad hardware and ROS abstraction. & \textbf{Algorithmic Baseline Data}: Logic planning and action primitives. 

\textbf{Modalities: RGB-D, Joint states, Contact forces, Proximity sensors, Force/Torque.} & \textbf{CALVIN}~\cite{mees2022calvin}, \textbf{RLBench}~\cite{james2020rlbench}, \textbf{GEM-Bench}, \textbf{Ravens}, \textbf{Colosseum} \\ \addlinespace

\textbf{Habitat / AI2-THOR} & Procedural indoor scene generation; High-throughput visual navigation kernels. & \textbf{Embodied Navigation Data}: Semantic search and indoor reasoning. 

\textbf{Modalities: RGB-D, Semantic/Instance segmentation, Object states, Agent pose, Audio.} & \textbf{OpenEQA}~\cite{majumdar2024openeqa}, \textbf{ALFRED}, \textbf{HM3D} \\ \addlinespace

\textbf{SoftGym / Genesis} & Unified MPM/FEM/PBD solvers; Differentiable multi-physics; Automated program generation. & \textbf{Generative Multi-Physics Data}: Coupled interactions (rigid/soft/fluid). 

\textbf{Modalities: RGB-D, Proprioceptive haptics, Particle states, Differentiable gradients.} & \textbf{ClothFunnels}~\cite{canberk2022clothfunnels}, \textbf{Difftaichi}~\cite{hu2020difftaichi}, \textbf{DexDeform}~\cite{li2023dexdeform} \\
\bottomrule
\end{tabular}%
}
\end{table*}

\textbf{MuJoCo\@: From Analytic Dynamics to Generative Generalization}

As one of the foundational physics kernels in embodied AI research, the evolutionary trajectory of \textbf{MuJoCo}~\cite{todorov2012mujoco} reflects the broader paradigm shift in robot learning---from state-driven atomic skills to vision-driven semantic generalization. 
Distinct from graphics-first engines, MuJoCo employs a soft-contact model and an analytic inverse dynamics solver, yielding strong numerical stability for contact-rich manipulation. 
Its ecosystem has progressed through three recognizable phases---algorithmic verification, data-driven learning, and generative simulation---tracking the growing demand for data scale and task complexity in embodied foundation models.
During the early phase of deep reinforcement learning integration, MuJoCo primarily served as a verification platform for state-space exploration. 
\textbf{Adroit}~\cite{rajeswaran2017learning} utilized MuJoCo's contact stability to establish high-dimensional dexterous manipulation benchmarks with a 24-DoF Adroit hand model (e.g., pen twirling, door opening, tool use), laying the groundwork for complex motor control. 
Subsequently, \textbf{Meta-World}~\cite{yu2021meta} expanded the scope to multi-task learning by defining 50 parameterized manipulation tasks on a Sawyer arm. 
However, benchmarks from this era generally suffered from \textbf{visual-semantic poverty}: scenes were constructed from textureless geometric primitives (geoms), with task variations limited to parametric coordinate randomization. 
While sufficient for state-based algorithms, this design left a significant Sim-to-Real visual gap, limiting their utility for training the visual representations required by modern embodied foundation models.
With the research focus shifting towards imitation learning and long-horizon planning, the MuJoCo ecosystem underwent a dual transformation in data scale and logical complexity, supported by the standardized \textbf{robosuite}~\cite{zhu2020robosuite} middleware. 
To address the scarcity of high-quality demonstration data, \textbf{RoboTurk}~\cite{mandlekar2018roboturk} leveraged MuJoCo's computational efficiency to build a crowdsourced teleoperation platform based on mobile devices, 
enabling the collection of large-scale human demonstrations that formed the foundation of \textbf{RoboMimic}~\cite{mandlekar2022matters}. 
Concurrently, the ecosystem bifurcated to address higher complexity: on the physical side, \textbf{Franka Kitchen}~\cite{gupta2019frankakitchen} evaluates long-horizon manipulation sequences requiring coordinated execution of multiple contact-rich skills (e.g., opening cabinets, turning burners, and manipulating switches), thereby stressing temporal compositionality and hierarchical policy learning; on the cognitive side, 
\textbf{LIBERO}~\cite{liu2023libero} introduced procedural task generation across 130 tasks in four task suites to evaluate \textbf{lifelong learning} and \textbf{compositional generalization}. 
More recently, \textbf{VLABench}~\cite{zhang2024vlabench} further extended this line by providing 100 task categories with over 2,000 objects and language-conditioned instructions, establishing the first large-scale MuJoCo-based benchmark for evaluating embodied foundation models on long-horizon reasoning and fine-grained language grounding. 
This phase marked MuJoCo's transition from a physics simulator to a comprehensive platform for assessing both physical dexterity and cognitive architecture.
Entering the era of Foundation Models, MuJoCo has undergone a shift towards generative and parallel simulation. 
On the computational front, Google DeepMind introduced \textbf{MJX}\footnote{\url{https://mujoco.readthedocs.io/en/stable/mjx.html}}, re-implementing the physics engine as a JAX computation graph. 
This enables end-to-end differentiable, massively parallel simulation on TPUs and GPUs, substantially improving throughput and alleviating efficiency bottlenecks for large-scale data generation. 
Building on MJX, \textbf{MuJoCo Playground}~\cite{deepmind2025playground} provides an open-source suite of GPU-accelerated environments with demonstrated zero-shot sim-to-real transfer across quadrupeds, humanoids, and dexterous hands. 
What's more, \textbf{DiffMJX}~\cite{paulus2025diffmjx} further improves differentiable simulation under hard contacts, making gradient-based optimization more reliable in contact-rich control settings.
On the content front, \textbf{RoboCasa}~\cite{nasiriany2024robocasa} integrates generative AI tools (LLMs and text-to-image models) to synthesize 120 kitchen scenes with over 2,500 objects across 100 tasks, narrowing the visual realism gap. 
Its successor, \textbf{RoboCasa365}~\cite{nasiriany2025robocasa}, scales this further to 365 everyday tasks across 2,500 diverse kitchen environments, accompanied by over 600 hours of human demonstration data and 1,600 hours of synthetically generated data, providing a large-scale training and benchmarking framework for generalist robots. 
Additionally, \textbf{BiGym}~\cite{chernyadev2024bigym} establishes a benchmark of 40 tasks for mobile bi-manual manipulation within MuJoCo. 
Together, the combination of the \textbf{MJX parallel computation backend} with \textbf{generative content pipelines} positions MuJoCo as a capable platform for producing foundation-model training data that is both physically grounded and semantically diverse.

\textbf{SAPIEN\@: Articulated Physical Foundations and Cross-Domain Generalization}

As a foundational physics kernel in embodied AI research, SAPIEN~\cite{xiang2020sapien} establishes a critical bridge between visual perception and physically grounded interaction. 
Built upon the PhysX backend, it natively supports analytical joint constraints and part-level mobility annotations, directly addressing the physical complexities of articulated manipulation. 
This architecture is sustained by the PartNet-Mobility dataset~\cite{xiang2020sapien}, providing over 2,300 assets with hierarchical kinematic trees and motion-limit metadata to establish a standardized corpus for learning visual-semantic priors within a unified kinematic framework.

Building upon this kernel, the ManiSkill series evolved computational paradigms to scale skill generalization and simulation efficiency. ManiSkill 1~\cite{mu2021maniskill} established the baseline for 3D point cloud-based policies, while ManiSkill 2~\cite{gu2023maniskill2} integrated a Position-Based Dynamics (PBD) solver to encompass fluid and soft-body tasks, yielding over 4 million expert frames. Recognizing the data generation bottleneck in policy learning, ManiSkill 3~\cite{tao2025maniskill3} transitioned to a GPU-parallelized architecture. By achieving throughputs exceeding 30,000 FPS on consumer hardware, it drastically accelerates reinforcement learning and enables the massive data synthesis required for scalable whole-body and dexterous manipulation. The task complexity within this ecosystem is further enriched by ManiSkill-HAB~\cite{shukla2025maniskillhab}, which enforces realistic physical contact for long-horizon rearrangement by restricting teleportation, and GAPartNet~\cite{geng2023gapartnet}, which introduces a part-centric perception framework to enable zero-shot interaction with novel functional components. Extending the cognitive demands of this ecosystem, \textbf{MIKASA-Robo}~\cite{cherepanov2026mikasarobo} leverages ManiSkill3 to benchmark memory-intensive tabletop manipulation, requiring agents to retain and recall task-relevant context across extended interaction horizons.

The maturity of this ecosystem naturally shifts the research focus from qualitative simulation to rigorous sim-to-real evaluation and automated data synthesis. 
The SIMPLER~\cite{li2025simpler} framework establishes an empirical bridge, demonstrating a high Pearson correlation between simulated success rates and real-world deployment. 
This predictable fidelity enables the paradigm advances toward synthesizing complex, high-dimensional coordination data. RoboTwin~\cite{mu2024robotwin} pioneered bimanual digital twin benchmarks to mitigate the visual-dynamic gap for dual-arm coordination. This was significantly scaled by RoboTwin 2.0~\cite{chen2025robotwin2}, which utilizes multimodal large language models and simulation-in-the-loop feedback to automatically synthesize over 100,000 domain-randomized expert trajectories across diverse dual-arm tasks. Ultimately, pushing beyond the constraints of a single physical embodiment, RoboFactory~\cite{qin2025robofactory} scales this automated synthesis to multi-agent manipulation. By enforcing strict logical, spatial, and temporal compositional constraints, RoboFactory guarantees collision-free collaboration, completing the SAPIEN ecosystem's trajectory from single-arm physics validation to scalable, multi-agent behavioral generation.

\textbf{Isaac Ecosystem: Massive GPU Parallelism and Foundation-Scale Data Synthesis}

Within the embodied AI ecosystem, the NVIDIA Isaac Gym introduces a GPU-resident simulation architecture that significantly accelerates robot learning pipelines. Built upon NVIDIA PhysX with GPU-accelerated rigid-body dynamics, Isaac enables physics stepping, observation extraction, and reward computation to remain within GPU memory, eliminating the frequent CPU–GPU data transfers that characterize traditional simulators. This design allows thousands to tens of thousands of parallel environments to be simulated simultaneously, yielding orders-of-magnitude improvements in simulation throughput and enabling large-scale reinforcement learning for robotic control.

The ecosystem's expansion is anchored by the integration of NVIDIA Omniverse and the adoption of OpenUSD as the universal schema for physical AI\@. This architecture enables Isaac to ingest heterogeneous industrial assets from CAD and CG pipelines, converting them into semantically rich digital twins. The paradigm is exemplified by the \textbf{BEHAVIOR-1K} dataset~\cite{li2024behavior}, which leverages Omniverse to provide 1,000 household tasks with physically grounded common-sense priors, and \textbf{ARNOLD}~\cite{gong2023arnold}, which establishes a benchmark for long-horizon language-grounded interactions. To address embodied data scarcity, massively parallel physics stepping is harnessed to explore high-dimensional configuration spaces at scale, as demonstrated by the \textbf{DexGraspNet} dataset~\cite{wang2023dexgraspnetlargescaleroboticdexterous}, which performs 1.32 million grasp attempts across 5,355 object categories directly within the GPU physics kernel. Leveraging the same infrastructure, \textbf{FurnitureSim}~\cite{heo2023furniturebench}---built upon Isaac Gym and Factory---provides a high-fidelity simulated counterpart to the FurnitureBench real-world assembly benchmark, enabling accurate simulation of screw-fastening dynamics for long-horizon precision tasks. At the framework level, \textbf{Isaac Lab}~\cite{mittal2023orbit} serves as the unified, modular entry point for robot learning within this ecosystem, supporting reinforcement learning, imitation learning, and motion planning across diverse embodiments.

Beyond raw throughput, Isaac has advanced toward high fidelity and differentiable intelligence. The Newton Solver ensures the physical fidelity required for precision-critical and quasi-static tasks by providing sub-millimeter contact stability---a prerequisite for benchmarks such as \textbf{Factory}~\cite{narang2022factory} and \textbf{IndustReal}~\cite{schoonbeek2024industreal}. 
Complementing this, NVIDIA Warp provides the differentiable backbone necessary for gradient-based optimization, which is critical for emerging paradigms that utilizes Warp's differentiable kernel to align high-frequency tactile feedback with linguistic instructions in contact-rich environments.
Beyond differentiable physics, scaling training corpora demands procedural generation and high-throughput architectures. 
\textbf{MimicGen}~\cite{mandlekar2023mimicgen} addresses this by autonomously multiplying limited human demonstrations through spatial randomization. 
On the other hand, \textbf{BiDexHands}~\cite{chen2022bidexhands} tackles the computational bottleneck of high-dimensional continuous control by leveraging fully tensorized, GPU-resident simulation to execute thousands of bimanual manipulation tasks in parallel. 
Ultimately, these complementary strategies shift the embodied data pipeline from labor-intensive manual collection to industrialized, compute-driven synthesis.

The ecosystem's dominance in high-DoF morphology is further solidified by its capacity to resolve the intricate mapping between human kinesthetics and robotic execution. 
By operationalizing massively parallel motion retargeting, Isaac functions as a morphological bridge: frameworks such as ExBody~\cite{cheng2024exbody} and OmniH2O~\cite{he2024omnih2o} leverage its accelerated dynamics to reconcile structural discrepancies between biological and mechanical skeletons in real-time. 
This capability facilitates a ``human-to-humanoid'' data flywheel, where raw teleoperated expertise is distilled into robust whole-body controllers via the HumanoidBench~\cite{sferrazza2024humanoidbenchsimulatedhumanoidbenchmark} evaluation suite. 
By unifying massively parallel sampling, semantic asset ingestion, and differentiable intelligence, Isaac establishes itself as a self-evolving engine for the large-scale pre-training and validation of embodied foundation models.

\textbf{Specialized Task Platforms: From Logic Primitives to Multi-Physics Autonomy}

Beyond the massive-scale frameworks, specialized platforms provide critical validation grounds for specific sub-fields of embodied AI, offering heterogeneous kernels tailored for logic planning, semantic navigation, and multi-physics synthesis.

\textbf{Lightweight Logic and Action Primitives.} \textbf{PyBullet}\footnote{\url{http://pybullet.org}} and \textbf{CoppeliaSim}~\cite{rohmer2013vrep} constitute the long-standing lightweight physics engines for manipulation research, valued for their reliable contact dynamics and streamlined computational overhead. Upon these engines, two foundational benchmarks have anchored most subsequent studies: \textbf{RLBench}~\cite{james2020rlbench}, built on CoppeliaSim, provides a diverse library of 100+ precision-critical tasks spanning pick-and-place, assembly, and tool use; \textbf{Ravens}~\cite{zhang2019raven}, built on PyBullet, pioneered the learning of spatial pick-and-place primitives through Transporter Networks. Extending the CoppeliaSim ecosystem to long-horizon evaluation, \textbf{CALVIN}~\cite{mees2022calvin} introduced a benchmark for multi-step, language-conditioned task execution, shifting the focus from single-step precision to sequential compositional reasoning. Along a complementary axis, \textbf{PerAct2}~\cite{shridhar2024peract2} extends the RLBench ecosystem to bimanual manipulation, benchmarking coordinated dual-arm tasks that demand richer action spaces beyond single-arm settings. However, these classic benchmarks primarily evaluate in-distribution performance, leaving robustness and out-of-distribution generalization underexplored. To close this gap, \textbf{Colosseum}~\cite{pumacay2024colosseum} systematically augments RLBench tasks with multi-axis physical perturbations (lighting, texture, object pose), while \textbf{GEM-Bench}~\cite{garcia2025gembench} establishes hierarchical generalization metrics to quantify zero-shot transferability across novel objects, instructions, and task compositions.

\textbf{Navigation and Semantic Reasoning.} \textbf{Habitat}~\cite{szot2022habitat20} and \textbf{AI2-THOR}~\cite{kolve2022ai2thor} serve as the two dominant simulation platforms for indoor navigation and semantic reasoning research. Their capabilities are grounded in large-scale 3D scene assets: Habitat draws on photorealistic reconstructions such as Matterport3D~\cite{chang2017matterport3d} and HM3D~\cite{ramakrishnan2021hm3d}, while AI2-THOR provides procedurally generated interactive rooms with fine-grained object affordances. Upon these platforms, early benchmarks defined the core evaluation axes: \textbf{ALFRED}~\cite{shridhar2020alfred} established language-guided multi-step task completion in AI2-THOR, and \textbf{TEACh}~\cite{padmakumar2022teach} extended this to dialogic human--agent collaboration. Within the Habitat ecosystem, OpenEQA~\cite{majumdar2024openeqa} introduced embodied question answering over real-world 3D scans, while \textbf{HomeRobot}~\cite{yenamandra2023homerobot} proposed the Open-Vocabulary Mobile Manipulation (OVMM) benchmark, requiring agents to locate arbitrary objects in unseen environments and place them on specified receptacles---jointly testing open-vocabulary recognition, navigation, and grasping. However, each of these benchmarks targets an isolated capability slice---instruction following, question answering, or mobile manipulation---making cross-capability comparison difficult. To unify this fragmented landscape, \textbf{EmbodiedBench}~\cite{yang2025embodiedbench} consolidates heterogeneous tasks across AI2-THOR, Habitat, and Minecraft, spanning high-level semantic planning and low-level atomic control, and quantifies 15 distinct embodied capabilities under consistent metrics.

\textbf{Generative Multi-Physics and Coupled Interactions.} The frontier of simulation for embodied AI is advancing from isolated rigid-body dynamics toward unified multi-physics and generative task synthesis. Simulating highly articulated and deformable objects introduces extreme degrees of freedom, challenging traditional contact dynamics. The evolution of this domain can be traced through specialized, material-specific simulators. For 1D and 2D deformables, such as ropes and garments, early environments like \textbf{SoftGym}~\cite{lin2021softgym} and \textbf{DeformableRavens}~\cite{seita2021deformableravens} established essential baselines. As the demand for complex topological interactions and varied material properties grew, comprehensive simulation benchmarks like \textbf{GarmentLab}~\cite{lu2024garmentlab} emerged. Concurrently, manipulation frameworks advanced from task-specific pipelines like \textbf{ClothFunnels}~\cite{canberk2022clothfunnels} to unified, category-level approaches such as \textbf{UniGarmentManip}~\cite{wu2024unigarmentmanip}.
For 3D elastoplastic materials and fluids, environments like \textbf{FluidLab}~\cite{xian2023fluidlab} pioneered complex state representations, which were further supported by differentiable skill frameworks such as \textbf{DiffSkill}~\cite{lin2022diffskill}. Researchers also extended these capabilities to the dexterous manipulation of soft bodies via \textbf{DexDeform}~\cite{li2023dexdeform}, and began generating deformable manipulation datasets by retrofitting soft-body support into traditional engines.
However, these approaches predominantly relied on fragmented, standalone solvers and early differentiable physics optimization~\cite{hu2020difftaichi}, struggling to natively handle cross-material interactions. Recently, unified platforms such as \textbf{Genesis}~\cite{genesis2024} have fundamentally restructured simulation pipelines by integrating diverse numerical methods---including the Material Point Method (MPM), FEM, and Position Based Dynamics (PBD)---into a single cohesive framework. By natively supporting the simulation of ``Coupled Interaction Data,'' these platforms capture the complex interplay among rigid, soft, and fluid bodies. This comprehensive ecosystem allows models to learn rich, complex physical interactions through massively scaled synthetic data, enabling robust, zero-shot generalizability in real-world deployments.

\subsubsection{Automatic Data Generation}

To address the challenge of data scarcity in embodied AI, simulation and synthetic data generation have emerged as critical solutions. This section reviews recent advancements across three key dimensions: the automated synthesis of \textbf{tasks and demonstrations} to scale behavioral data, the construction of diverse \textbf{environmental scenes} with physical fidelity, and the modeling of \textbf{dynamic evolution} via generative world models.

\textbf{Scalable Environment and Scene Synthesis}

Traditional simulation environments have long been constrained by the prohibitive costs of manual modeling, creating a severe bottleneck in scene diversity for training data. The evolution of this field originated with the brute-force augmentation of environment scale via \textbf{Procedural Content Generation (PCG)} and has gradually converged towards semantic controllability. \textbf{ProcTHOR}~\cite{deitke2022procthor} established the benchmark for large-scale generation by algorithmically sampling over 10,000 3D indoor environments with randomized layouts and lighting stored in memory. This strategy successfully addressed the out-of-distribution (OOD) generalization problem in navigation tasks, demonstrating that environmental scale is a critical dimension for enhancing the robustness of embodied agents. However, pure stochastic generation often lacks semantic logic (e.g., placing a bed in a kitchen), making it inadequate for complex long-tail tasks. To bridge this gap, \textbf{Holodeck}~\cite{yang2024holodeck} introduced Large Language Models (LLMs) as a semantic intermediary atop PCG, enabling a precise mapping from natural language instructions to 3D environments. Through this ``semantic interface,'' researchers can tailor specific scenarios via high-level prompts (e.g., ``generate a crowded music room''), thereby achieving a critical leap from randomized scaling to semantic customization.

With the issues of layout scale and semantics addressed, the research focus shifted towards physical interactivity and visual perceptual fidelity to mitigate the Sim-to-Real gap. \textbf{PhyScene}~\cite{yang2024physcene} targeted common physical artifacts in synthetic scenes—such as object interpenetration and levitation—by utilizing physics engine constraint solving to optimize the kinematic joints and collision volumes of generated objects, thus ensuring scene interactability. Concurrently, \textbf{BEHAVIOR Vision Suite}~\cite{ge2024bvs} focused on the robustness of visual perception. By allowing high-degree customization of sensor parameters, lighting conditions, and object textures within the simulation, BVS provided a systematic testbed that closely mimics the noise distribution of the real world. This correction from mere geometric appearance to dual alignment of physics and perception lays a credible foundation for simulation-based training of fine-grained manipulation tasks. Complementing these from-scratch approaches, \textbf{GAIA}~\cite{ko2026gaia} grounds simulation generation directly in the real world: given a single RGB image and a task instruction, a pre-trained VLM jointly infers the visual context and task requirements, then places task-relevant objects---including those not visible in the original image---to produce an interactive, task-ready environment without additional training or manual setup. Policies learned within GAIA-generated environments transfer to the corresponding target scenes, demonstrating that real-context anchoring can complement physics and perceptual alignment in closing the sim-to-real gap.

Currently, scene synthesis is advancing towards fully automated closed-loops and macroscopic scales, marking the evolution of simulation environments from static training grounds to self-generating embodied data factories. \textbf{RoboGen}~\cite{wang2023robogen} represents a form of automation, constructing an LLM-driven self-evolution loop that automatically retrieves assets, generates scenes, and trains policies based on task requirements, completely eliminating human intervention. On the application side, \textbf{RoboCasa}~\cite{nasiriany2024robocasa} leveraged Generative AI (Text-to-3D/Image) to create high-fidelity kitchen environments covering hundreds of everyday household tasks, serving as a core benchmark for generalist manipulation capabilities. Furthermore, \textbf{GRUtopia}~\cite{wang2024grutopia} expanded the horizon from indoor settings to city-level scales, supporting large-scale agent collaboration and evolution within complex environments involving social interactions. Collectively, these works construct a synthetic data ecosystem capable of self-generation, self-verification, and infinite scalability.

\textbf{Generative Task and Demonstration Synthesis}

The methods for automatically generating robot training data have evolved along three axes: from spatial augmentation of existing demonstrations to LLM-driven task synthesis, and most recently toward physically grounded generation that handles complex morphologies and failure recovery. Early work established the core primitives: \textbf{MimicGen}~\cite{mandlekar2023mimicgen} multiplies a single human demonstration into thousands of valid trajectories via rigid-body spatial transformations, while \textbf{ROSIE}~\cite{yu2023scalingrobotlearningsemantically} uses text-to-image diffusion to randomize visual appearances. With the rise of LLMs, \textbf{GenSim}~\cite{wang2024gensim} and \textbf{RoboGen}~\cite{wang2023robogen} began automating the creation of tasks themselves by generating simulation code and reward functions, and \textbf{SkillMimicGen}~\cite{garrett2024skillmimicgen} decomposed long demonstrations into reusable skill primitives to prevent compounding errors in multi-step tasks.

A critical challenge is teaching robots to recover from errors and handle edge cases. \textbf{IntervenGen}~\cite{hoque2024intervengen} algorithmically generates recovery trajectories from a small set of human corrections, providing examples of how to return to a successful path after failure. Scaling this to fully automated boundary-case mining, \textbf{Fail2Progress}~\cite{huang2025fail2progress} uses Stein Variational Inference (SVI) in parallel simulations to synthesize massive, high-variance difficult-scenario data from real robot failure cases, efficiently patching policy vulnerabilities without additional human intervention. Closing the loop from data generation to policy refinement, \textbf{Self-Improving VLA (PLD)}~\cite{xiao2025selfimprovingvla} lets a residual RL expert actively intervene during base policy rollouts, automatically generating high-quality mixed trajectories containing both corrective and recovery behaviors, which are then used to fine-tune the downstream foundation model in a self-improving cycle.

Generating demonstrations for complex and diverse robot morphologies presents a distinct set of challenges. \textbf{MoMaGen}~\cite{li2026momagengeneratingdemonstrationssoft} formulates demonstration generation for bimanual mobile manipulators as a constrained optimization problem, jointly satisfying hard constraints (e.g., arm reachability) and soft constraints (e.g., camera visibility during navigation) to synthesize diverse datasets from a single source demonstration. \textbf{HumanoidGen}~\cite{jing2025humanoidgen} targets bimanual humanoid robots with dexterous hands, using LLMs to coordinate macro arm movements with fine-grained finger manipulations. For dexterous grasping at scale, \textbf{UltraDexGrasp}~\cite{yang2026ultradexgrasp} fuses optimization-driven grasp pose synthesis with planning-based demonstration generation, automatically constructing a 20-million-frame dataset of multi-strategy bimanual grasps spanning thousands of object categories. To bridge the embodiment gap between human demonstrations and robot execution, \textbf{OmniRetarget}~\cite{yang2025omniretarget} employs an interaction-mesh Laplacian deformation engine to retarget human videos into robot training data while rigorously maintaining human-object-environment contact relationships without interpenetration artifacts. Scaling this autonomous synthesis paradigm to foundation-model pre-training, \textbf{InternData-A1}~\cite{cai2026internvla} constructs a fully decoupled pipeline atop the Isaac ecosystem that separates asset specification, skill policies, task composition, and photorealistic rendering, producing 630k trajectories across 4 embodiments, 70 tasks, and 227 scenes---covering rigid, articulated, deformable, and fluid objects---at a cost below \$0.003 per episode on 8 RTX 4090 GPUs.

A parallel line of work bypasses physics simulators entirely, leveraging 3D representations to proliferate demonstrations with high visual fidelity. \textbf{DemoGen}~\cite{xue2025demogen} generates new training data directly by editing 3D point clouds or synthesizing new spatial configurations from real-world data. \textbf{Tether}~\cite{liang2026tether} advances this to dynamic physical interaction, using semantic correspondence to warp demonstrations to new configurations and employing VLMs to verify and self-curate successful trials in the real world. \textbf{RoboSplat}~\cite{robosplat2025} operates within a unified 3D Gaussian Splatting representation to perform scene-level object replacement, SE(3)-equivariant manipulation, and novel viewpoint synthesis, enabling a single real-world demonstration to be proliferated into massive multi-modal training datasets. \textbf{Real2Render2Real (R2R2R)}~\cite{yu2025r2r2r} pushes this further by completely bypassing rigid-body dynamics simulation, relying purely on 3D Gaussian reconstruction and spherical linear interpolation for kinematic rendering to multiply a single real human video into thousands of high-quality cross-configuration robot training trajectories.

\textbf{Generative World Models Synthesis}

Traditional physics simulators struggle to accurately model soft objects, fluids, and complex visual details, limiting their usefulness for training real-world robots. 
To address these limitations, recent research has leveraged world models to synthesize future states and long-horizon trajectories within learned latent spaces; hereafter, we refer to such models as \textbf{generative world models}.
By approximating transitions through latent dynamics predictors, these frameworks support data-driven rollouts that facilitate planning and policy learning in complex environments.

Recent work focuses on making these video generation models obey basic physical laws. For example, \textbf{NVIDIA Cosmos}~\cite{nvidia2025cosmos} trains on large-scale video datasets to predict future frames that implicitly follow Newtonian mechanics and contact dynamics, without relying on a traditional physics engine. Its successor, \textbf{Cosmos-Predict2.5}~\cite{nvidia2025cosmospredict25}, further unifies Text2World, Image2World, and Video2World generation in a single flow-based architecture, achieving substantial improvements in video quality and instruction alignment. 
\textbf{iVideoGPT}~\cite{wu2024ivideogpt} treats physical simulation as a sequence prediction problem, taking the robot's current visual observation and planned action to predict the next visual state. 
To specifically simulate robot-environment interactions, \textbf{IRASim}~\cite{zhu2025irasim} uses a Diffusion Transformer trained on real robot data to generate realistic video responses to specific robot commands, capturing the actual kinematics and physical feedback of the hardware. 
To explicitly prevent the visual hallucinations common in these generative processes, \textbf{AnchorDream}~\cite{ye2025anchordream} repurposes video diffusion models by injecting exact robot kinematic trajectories as conditional anchors, ensuring the synthesized interactions remain strictly embodiment-aware and physically consistent.

Beyond simply predicting the next frame, these models allow robots to imagine and plan for new situations. 
\textbf{RoboDreamer}~\cite{zhou2024robodreamer} uses the latent space of generative models to combine different objects and actions, allowing the robot to practice tasks it has never seen before. 
\textbf{DreMa}~\cite{barcellona2025dream} uses video generation to simulate different future trajectories, helping the robot choose the best sequence of actions through self-supervised learning.
Extending this imagination capability to explicit data synthesis, \textbf{DreamGen}~\cite{jang2025dreamge} employs action-conditioned video foundation models to simulate environment state transitions within continuous latent spaces, generating pre-rehearsal datasets with rich long-tail attributes for training cross-scene robot policies. Expanding beyond targeted demonstrations,, \textbf{DreamDojo}~\cite{gao2026dreamdojo} extracts generalized interaction dynamics features directly from massive first- and third-person internet videos, synthesizing physically plausible world model pre-training priors in latent space without requiring any robot-specific data.
Despite their capabilities, standard video models frequently suffer from physically impossible artifacts. To address this, EnerVerse~\cite{huang2025enerverse} enforces strict 3D geometric constraints by predicting how a scene's 3D volume changes over time, whereas EnerVerse-AC~\cite{jiang2025enerverseac} integrates explicit action conditions to guarantee that the generated 3D future precisely aligns with the robot's control intentions. As these neural simulators become more accurate, they are increasingly used to train and test robot policies directly: \textbf{WorldGym}~\cite{quevedo2025worldgym} uses generative world models as full RL environments, allowing robots to learn by interacting with generated video streams rather than a physics engine.

\subsection{Data Engine}\label{sec:data_augmentation}

Because collecting real-world robot data is expensive and time-consuming, researchers increasingly rely on data augmentation and refinement to maximize the value of existing demonstrations. We categorize these efforts into three areas: \textbf{visual augmentation} to improve perception, \textbf{physical and temporal constraints} to ensure augmented data remains realistic, and \textbf{data filtering} to extract high-quality signals from noisy datasets.

\subsubsection{Visual Augmentation and Perceptual Synthesis}\label{sec:trajectory_redrawing}

To help robots recognize objects in new environments, researchers use generative AI to alter the visual appearance of training data without changing the underlying robot actions. \textbf{GenAug}~\cite{chen2023genaug} takes this further by automatically pasting new 3D objects into the scene to act as visual distractors. To handle different camera setups, \textbf{RoVi-Aug}~\cite{chen2024roviaug} leverages image-to-image diffusion models to perform robot and viewpoint augmentation, transforming demonstrations into diverse robot appearances and camera perspectives. This synthetic diversification improves policy generalization across robot platforms and camera configurations. Advancing this viewpoint augmentation with explicit pose control, \textbf{ROPA}~\cite{chen2025ropa} fine-tunes a skeleton-pose-conditioned diffusion model to offline synthesize third-person RGB-D data and action labels from entirely novel viewpoints, while preserving physical contact constraints and geometric consistency. Taking a complementary approach, \textbf{Action-View Augmentation}~\cite{pan2025one} combines novel view synthesis with manipulation-level trajectory warping to automatically derive one thousand high-variance training trajectories from a single expert demonstration while maintaining kinematic coherence. Going beyond static appearance and viewpoints to address covariate shift, \textbf{DMD}~\cite{zhang2024dmd} leverages diffusion models to hallucinate corrective eye-in-hand observations from perturbed states, effectively enabling DAgger-style recovery learning without the need for costly physical data collection.

However, applying these image-based edits frame-by-frame often causes flickering, which disrupts the robot's ability to track motion. To solve this, \textbf{VISTA}~\cite{wang2024vista} uses video diffusion models to ensure the generated visual streams remain smooth over time. To prevent physical impossibilities—like generated objects floating in mid-air—\textbf{DABI}~\cite{li2024dabi} uses depth maps to guide the generation process, ensuring the new visual elements respect the 3D geometry of the original scene. Finally, \textbf{LucidSim}~\cite{mishra2024lucidsim} combines these generative techniques with visual distillation, creating a robust pipeline that effectively bridges the visual gap between synthetic simulations and the physical world.

\subsubsection{Physical Constraints and Temporal Logic}\label{sec:dynamics_constraints}

When augmenting the physical trajectories of a robot (e.g., shifting a grasping motion 5cm to the left), the new data must still obey the laws of physics. \textbf{EIL}~\cite{we2024equivariant} applies SE(3) equivariance, a mathematical property ensuring that if an object is rotated or translated, the robot's predicted action rotates or translates by the exact same amount. To prevent augmented trajectories from passing through solid objects, \textbf{EAD}~\cite{geng2024environment} and \textbf{DAAG}~\cite{jia2024dynamics} incorporate collision checking and friction models, ensuring that every newly generated trajectory is actually executable in the real world. Extending this to one-shot trajectory augmentation, \textbf{CP-Gen}~\cite{lin2025cpgen} explicitly integrates kinematic boundaries and contact dynamics constraints during the generation and augmentation of single-sample supervised trajectories, ensuring that the synthesized visual-motor policy data remains both safe and physically feasible.

Beyond spatial constraints, researchers are augmenting the temporal logic of tasks to teach robots how to recover from mistakes and handle sparse rewards. \textbf{Diff-DAgger}~\cite{lo2024diffdagger} uses diffusion models to generate ``recovery trajectories''—simulating what the robot should do if it accidentally drops an object or misses a grasp. \textbf{Generative Predecessor Models}~\cite{yang2024generative} work backward, taking a successful final state and generating plausible past states that could have led to it. To address environments where success signals are rare, \textbf{GoalGAIL}~\cite{ding2024goalgail} uses adversarial mechanisms to generate virtual goal trajectories. Finally, \textbf{Self-Imitation by Planning}~\cite{sun2024self} takes failed robot attempts and uses motion planning to correct them, turning useless failure data into valuable training signals.

\subsubsection{Data Filtering and Systemic Optimization}\label{sec:data_refinement}

As datasets grow larger and incorporate data from various sources (teleoperation, scripted policies, different robot types), they inevitably contain noise and suboptimal behaviors. The focus here shifts from generating more data to filtering out the bad data. \textbf{L2D}~\cite{yue2024learning} uses a preference-based learning mechanism to automatically identify and extract the smoothest, most efficient segments from messy human demonstrations. For offline learning, \textbf{ILID}~\cite{belkhale2023data} and \textbf{DC-IL}~\cite{kim2024distribution} apply distribution correction techniques to prevent the robot from imitating the worst behaviors in a dataset, ensuring it only learns from the highest-quality signals.

When mixing datasets from entirely different robots and labs, simply combining them can actually hurt performance (negative transfer). \textbf{Re-Mix}~\cite{jason2024remix} solves this by dynamically adjusting the training weights of different datasets during training, ensuring the model doesn't overfit to one specific robot type. To handle massive, redundant datasets, \textbf{JUICER}~\cite{haldar2024juicer} uses distillation to throw away repetitive data and keep only the most informative samples, while \textbf{RoboPearls}~\cite{li2024robopearls} specifically searches for and upweights rare, long-tail scenarios that the robot struggles with. Pushing self-supervised curation to the scale of foundation models, \textbf{SCIZOR}~\cite{zhang2026scizor} operates at the fine-grained transition level to filter out both suboptimal and redundant state-action pairs without requiring external reward signals. It employs a task progress predictor to discard segments lacking meaningful progression (e.g., collisions or jittering) and clusters joint visual-action embeddings to deduplicate overrepresented patterns while preserving rare variations. Approaching curation from a gradient-based perspective, \textbf{CUPID}~\cite{agia2025cupidcuratingdatarobot} pioneers the use of influence functions at the end of automatic data generation pipelines, evaluating each synthetic trajectory's direct gradient contribution to policy performance to intelligently curate and discard harmful noise samples. Extending this influence-based paradigm to generative data distillation, \textbf{FT-NCFM}~\cite{wang2025NCFM} introduces a Fact-Tracing engine that combines influence functions with neural characteristic function matching to distill high-value synthetic coresets from large-scale embodied manipulation datasets, maximizing downstream policy performance per training sample. Complementing this with semantic reasoning, \textbf{RoboCurate}~\cite{kim2026robocurate} introduces a VLM-based physical attribute alignment mechanism to intelligently filter and curate noisy synthetic actions produced by video generation models, ensuring that only physically plausible data survives the curation process. Pushing this further to the algorithmic level, \textbf{OptimSyn}~\cite{fan2026optimsyn} employs influence-function-guided reward models to train a synthetic data evaluation engine (Rubrics), automatically selecting and generating the training sets that most improve downstream task performance. Finally, to ensure efficient data flow within distributed architectures like \textbf{CACTI}~\cite{cacti2024} and \textbf{Scalable Multi-Task IL}~\cite{mtil2024}, \textbf{OXE-AUGE}~\cite{oxe2024auge} provides a standardized way to align data across different robot embodiments, which is further complemented by \textbf{OILCA}~\cite{oilca2024} to solve environment generalization issues through contextual augmentation across petabyte-scale heterogeneous data streams.

\subsection{Large-Scale Data Pipelines}\label{sec:data_factories_robotic_arms}

The evolution of robotic manipulation increasingly relies on scalable infrastructures capable of generating, managing, and utilizing large datasets. Unlike general navigation, manipulation requires processing high-frequency contact dynamics and fine-grained hand-eye coordination. This necessitates \textbf{Integrated Data Pipelines} that combine procedural generation, physical simulation, and real-world data collection. This section reviews eight key frameworks that support large-scale data-centric manipulation research.

\textbf{RoboTwin}~\cite{mu2024robotwin} provides a comprehensive data generation framework for industrial dual-arm manipulation. Its initial iteration established a high-fidelity simulation environment for contact-rich tasks (e.g., belt drive assembly) by capturing the kinematics and sensor noise of physical setups to bridge the Sim-to-Real gap. Building on this, \textbf{RoboTwin v2}~\cite{chen2025robotwin2} expands the system to support automated task generation. By leveraging Multimodal Large Language Models (MLLMs), the platform automates the design of manipulation tasks, parsing abstract assembly instructions to synthesize executable task codes and asset configurations. This allows the system to scale from a fixed set of tasks to a wide variation of industrial scenarios, reducing the need for manual engineering in complex manufacturing environments.

\textbf{GRUTOPIA}~\cite{wang2024grutopia} serves as a simulation environment for multi-agent social interactions and service tasks. While the initial architecture provided infrastructure for large-scale simulation, \textbf{GRUTOPIA v2} expands its data generation capabilities through two core modules: GRScenes and GRResidents. GRScenes introduces a dataset of 100,000 interactive, annotated 3D scenes covering 89 functional categories. Complementing this, GRResidents deploys an LLM-driven NPC system to simulate human behaviors with long-term memory and intent. This combination enables the generation of social-norm-aware manipulation data—such as delivering items without interrupting conversations—training agents to operate in populated dynamic environments.

\textbf{MimicLabs}~\cite{mimicLabs2024} provides an analytical framework for data generation, investigating the scaling laws specific to manipulation tasks. It quantifies the trade-offs between data quantity, diversity (visual and geometric), and quality. By synthesizing heterogeneous demonstrations ranging from teleoperation to scripted policies, it establishes protocols for data composition. This provides guidelines on how data generation pipelines like RoboTwin and GRUTOPIA should balance their computational budgets to improve policy generalization across unseen objects and layouts.

\textbf{AIRSPEED}~\cite{xia2024airspeed} provides a universal data production platform that connects data collection hardware with simulators. It decouples the data generation logic from hardware backends, offering an API to integrate VR collection rigs, mocap systems, and various simulation environments. While initially demonstrated on high-speed dynamics, its modular design allows research groups to set up custom data collection pipelines for novel robotic arm configurations. This architecture lowers the barrier to entry, enabling the construction of simulation environments without rebuilding the underlying software stack.

\textbf{RoboVerse}~\cite{roboverse2024} provides a unified platform, dataset, and benchmark designed to support scalable robot learning. To address data fragmentation, it utilizes a standardized infrastructure, MetaSim, which integrates heterogeneous data sources and simulation interfaces. Based on this infrastructure, RoboVerse supplies a large-scale synthetic dataset comprising over 1,000 manipulation tasks, 5,500 object assets, and approximately 10 million interaction transitions (equivalent to about 500,000 trajectories). Beyond data generation, the platform establishes a systematic evaluation benchmark with multiple generalization levels. This allows for the structured assessment of learned policies across various dimensions, including environmental variations, unseen objects, and cross-domain transferability.

\textbf{LeRobot}~\cite{lerobot2024} acts as an open-source repository and framework for sharing robotic datasets. It simplifies the collection and sharing of real-world arm trajectories by providing a standardized data format. This ecosystem encourages the aggregation of fragmented datasets from isolated labs into larger, unified collections. By hosting these resources on platforms like Hugging Face, LeRobot facilitates the training of foundation models that can operate across distinct robotic arm morphologies, helping to address the data fragmentation problem.

\textbf{AgiBot World Colosseo}~\cite{agibot2025} provides a large-scale data collection and simulation Platform for humanoid robots, integrating physical data gathering with digital synthesis to address data scarcity for bi-manual manipulation. It employs a dual-cycle pipeline: physically, it deploys a fleet of over 100 AgiBot G1 humanoids to capture real-world interaction data; digitally, it utilizes the Genie Sim engine to augment this data via physics-compliant domain randomization. This platform has produced a dataset of over 1 million trajectories across 217 tasks, which is used to train the GO-1 policy through imitation and reinforcement learning, which is proven to be highly effective for training advanced manipulation policies via imitation and reinforcement learning.

\textbf{NVIDIA Isaac GR00T}\footnote{\url{https://developer.nvidia.com/isaac/gr00t}} provides an end-to-end platform for humanoid robot development, spanning synthetic data generation, foundation model training, and sim-to-real deployment. At its core, the \textbf{GR00T N1} foundation model adopts a dual-system architecture inspired by human cognition: a ``fast-thinking'' diffusion transformer generates continuous motor actions at high frequency, while a ``slow-thinking'' vision-language module handles deliberate reasoning and multi-step planning. The platform's data engine, \textbf{GR00T-Dreams}, leverages Cosmos world foundation models to synthesize large-scale neural trajectories from a single image input, extracting action tokens that can train manipulation and locomotion policies without real-world collection. Built atop the GPU-accelerated Isaac Lab simulation framework with RTX-based photorealistic rendering and Newton rigid-body dynamics, the ecosystem integrates high-fidelity physics simulation, domain randomization, and heterogeneous sensor modeling into a unified pipeline that bridges the sim-to-real gap for training generalist humanoid policies.

\textbf{DexFlyWheel}~\cite{zhu2025dexflywheel} constructs an autonomous flywheel by combining imitation learning with residual reinforcement learning in simulation. It actively explores and harvests high-success-rate contact trajectories for dexterous manipulation through trial-and-error, eliminating the need for manual intervention. Shifting the data source from simulation to real-world human videos, \textbf{RoboWheel}~\cite{zhang2025robowheel} builds a data engine that converts monocular hand-object interaction (HOI) footage into cross-embodiment robot supervision. It reconstructs contact-rich trajectories with RL-based physical plausibility refinement, then retargets them to diverse morphologies---arms, dexterous hands, and humanoids---providing the first quantitative evidence that HOI video can match teleoperation as an effective data source for robotic learning.

With the data infrastructure in place---from curated datasets through acquisition pipelines to industrial-scale data engines---the next question becomes how to convert this data into manipulation capability. The following section examines the learning paradigms that consume these data streams, spanning high-level cognitive planning, low-level policy learning, and the integration architectures that bridge them.

\section{Learning Paradigms of Robotic Arms}\label{sec:embodiment_method}

The data ecosystem of \S\ref{sec:embodiment_dataset} provides the raw material; the learning paradigms determine how that material is transformed into manipulation intelligence. Borrowing a neuroscience analogy, we decompose this transformation along two complementary pathways that mirror the division between the cerebral cortex and the cerebellum: \emph{cognitive planning} (\S\ref{sec:embodied_brain})---the ``embodied brain''---converts semantic intent into structured plans, progressing from code generation through task decomposition and foundation reasoning to geometric grounding and physical imagination; \emph{policy learning} (\S\ref{sec:policy_learning_low_level})---the ``embodied cerebellum''---maps perceptual observations to motor commands via explicit regression, discrete tokenization, or iterative generation. Neither pathway suffices alone: planning produces goals that lack motor precision, while policies execute reflexes that lack semantic awareness. Section~4.3 therefore examines integration architectures (Type~I/II/III) that fuse both pathways, and \S4.4 traces the deployment lifecycle from offline pre-training through online adaptation to runtime execution.

\subsection{Cognitive Planning: The Embodied Brain}\label{sec:embodied_brain}

The core function of high-level cognition is to bridge the semantic gap between abstract human intent and low-level robotic control. As the field matures, this interface has evolved from simple language-conditioned scoring to physics-informed generative simulation. We organize this evolution into five paradigms: \emph{Semantic Logic and Code Generation} (from probabilistic scoring to Turing-complete program synthesis), \emph{Long-Horizon Task Decomposition} (formal solvers and hierarchical structures to mitigate planning hallucinations), \emph{Foundation Reasoning Models} (multimodal understanding and physical common-sense grounding), \emph{Geometric Grounding and Spatial Intelligence} (translating semantic intent into precise 3D geometric actions), and \emph{Generative Dynamics and Physical Imagination} (predicting physical consequences before execution via learned world models). Table~\ref{tab:brain_paradigms} summarizes the paradigm-level properties and representative methods; Figure~\ref{fig:brain_timeline} traces their chronological evolution, and Figure~\ref{fig:brain_structure} illustrates their architectural relationships.
%% TODO: Insert \ref{fig:brain_timeline} — timeline figure here
%% TODO: Insert \ref{fig:brain_structure} — structure/taxonomy figure here

\begin{table*}[t!]
\centering
\caption{Taxonomy of Cognitive Planning Paradigms. \textbf{B1}\,=\,Semantic Logic \& Code Generation; \textbf{B2}\,=\,Long-Horizon Task Decomposition; \textbf{B3}\,=\,Grounded Multimodal Reasoning; \textbf{B4}\,=\,Geometric Grounding \& Spatial Intelligence; \textbf{B5}\,=\,Generative Dynamics \& Physical Imagination. Reasoning Level indicates the cognitive abstraction; Output Abstraction denotes the generated artifact; Execution Gap specifies what additional computation converts the output into motor commands; Data Tier indicates the dataset level (\S\ref{sec:eai_datasets}) that primarily supports each paradigm ($\mathcal{K}_1$\,=\,Skill, $\mathcal{K}_2$\,=\,Task, $\mathcal{K}_3$\,=\,Foundation).}\label{tab:brain_paradigms}
\renewcommand{\arraystretch}{1.15}
\setlength{\tabcolsep}{2.5pt}
\scriptsize
\begin{tabular}{
  c
  c
  c
  c
  c
  c
  p{0.22\linewidth}
}
\toprule
\textbf{ID} &
\textbf{Data Tier} &
\textbf{Sub-direction} &
\textbf{Reasoning Level} &
\textbf{Output Abstraction} &
\textbf{Execution Gap} &
\textbf{Representative Methods} \\
\midrule
\multirow{4}{*}{\textbf{B1}}
 & \multirow{4}{*}{$\mathcal{K}_2$} & Probabilistic Alignment & Semantic & Skill ranking & Skill API & SayCan~\cite{brohan2023can} \\
 & & Program Synthesis & Semantic & Executable code & Runtime interp. & CaP~\cite{liang2022code}, ProgPrompt~\cite{singh2022progprompt}, Demo2Code~\cite{wang2023demo2code} \\
 & & Evolutionary Sim-to-Real & Semantic & Reward functions & RL optimizer & DrEureka~\cite{dreureka2024}, Lang.\ to Reward~\cite{ma2023reward} \\
 & & Structured Control Flow & Logical & BT / XML & Reactive exec. & LLM-as-BT~\cite{ao2025llmbt} \\
\midrule
\multirow{4}{*}{\textbf{B2}}
 & \multirow{4}{*}{$\mathcal{K}_2$} & Formal Solvers & Logical & PDDL plan & Classical planner & LLM+P~\cite{liu2023llmp}, L3M+P~\cite{agarwal2025l3mp} \\
 & & Semantic Search & Logical & Sub-task sequence & Motion planner & SayPlan~\cite{sayplan2023} \\
 & & Closed-Loop Feedback & Logical & Re-plan & Skill API & Inner Monologue~\cite{huang2023inner} \\
 & & State Maintenance & Logical & State-dep.\ code & Runtime interp. & Statler~\cite{yoneda2024statler} \\
\midrule
\multirow{2}{*}{\textbf{B3}}
 & \multirow{2}{*}{$\mathcal{K}_2/\mathcal{K}_3$} & Multimodal Grounding & Multimodal & Text / action & Skill API & PaLM-E~\cite{driess2023palme}, RoboBrain~\cite{ji2025robobrain} \\
 & & Embodied CoT & Multimodal & CoT / ReAct trace & Skill API & EmbodiedGPT~\cite{mu2023embodiedgpt}, ReAct~\cite{yao2023react} \\
\midrule
\multirow{5}{*}{\textbf{B4}}
 & \multirow{5}{*}{$\mathcal{K}_1$} & Physical Constraint & Geometric & Contact / SE(3) pose & Motion planner & ManipLLM~\cite{manipllm2024}, LLM-GROP~\cite{llmgrop2024}, A3VLM~\cite{huang2024a3vlm} \\
 & & 2D Affordance & Geometric & Heatmap / mask & Grasp planner & CLIPort~\cite{shridhar2021cliport}, PASG~\cite{zhu2025pasg}, Moka~\cite{moka2024}, RAM~\cite{ram2025}, UAD~\cite{tang2025uad} \\
 & & 3D Neural Repr. & Geometric & Neural field & Motion planner & NDF~\cite{simeonov2021ndf}, GNFactor~\cite{ze2023gnfactor}, F3RM~\cite{f3rm2023}, D3Fields~\cite{wang2024d3fields}, GraspSplats~\cite{graspsplats2025}, CoPa~\cite{copa2024} \\
 & & Constraint Reasoning & Geometric & Value map / keypts & Trajectory opt. & VoxPoser~\cite{voxposer2023}, ReKep~\cite{rekep2024}, FoundationPose~\cite{wen2024foundationpose} \\
 & & Dexterous Manip. & Geometric & Grasp pose & Finger controller & DexGraspNet~\cite{wang2023dexgraspnetlargescaleroboticdexterous}, UniDexGrasp++~\cite{unidexgrasp2024}, GraspGPT~\cite{tang2023graspgpt}, OmniManip~\cite{pan2025omnimanip} \\
\midrule
\multirow{3}{*}{\textbf{B5}}
 & \multirow{3}{*}{$\mathcal{K}_1/\mathcal{K}_3$} & Latent Dynamics & Predictive & Latent rollout & MPC / actor-critic & Daydreamer~\cite{wu2023daydreamer}, TD-MPC2~\cite{hansen2024tdmpc2}, V-JEPA2~\cite{assran2025vjepa} \\
 & & Generative Video & Predictive & Future frames & Inverse dynamics & UniPi~\cite{du2023unipi}, VILA~\cite{du2024vila}, UWM~\cite{zhu2025uwm}, Genie~\cite{bruce2024genie}, Cosmos~\cite{nvidia2025cosmos}, DreamDojo~\cite{gao2026dreamdojo} \\
 & & Physics-Informed Plan. & Predictive & Trajectory / MCTS & Direct exec. & Deep Lagrangian~\cite{schulze2025lagrangian}, Particle-Grid~\cite{zhang2025particlegrid}, STORM~\cite{lin2025storm}, LUMOS~\cite{nematollahi2025lumos} \\
\bottomrule
\end{tabular}
\end{table*}

\subsubsection{Semantic Logic and Code Generation}\label{sec:code_gen}

The interface between high-level semantic reasoning and low-level control is the central challenge of cognitive planning. Early approaches framed this as a \textbf{Probabilistic Alignment} problem. Epitomized by \textbf{SayCan}~\cite{brohan2023can}, these methods ground abstract plans in physical affordances by multiplying the semantic probability of an instruction (derived from Large Language Models) with the execution probability of a primitive skill (derived from Value Functions). While effective for sequential tasks, this ``scoring-and-ranking'' paradigm is inherently limited to atomic actions and lacks the capacity to represent complex logic such as loops, conditionals, or variable passing.

To transcend these bottlenecks, the paradigm shifted towards \textbf{Program Synthesis}, where LLMs function as policy interpreters. 
\textbf{Code as Policies}~\cite{liang2022code} pioneered the generation of executable Python code, synthesizing behaviors from natural language. 
Enhancing this constraint adherence, \textbf{ProgPrompt}~\cite{singh2022progprompt} introduced a native programming environment prompting strategy that asserts preconditions within the prompt structure, effectively grounding intent into rigorous logical assertions.

Recent advancements have elevated this paradigm from one-shot generation to \textbf{Evolutionary Sim-to-Real Optimization}. \textbf{DrEureka}~\cite{dreureka2024} represents the state-of-the-art in this direction. Unlike static code generators, DrEureka employs an LLM as an evolutionary search operator to iteratively write and refine reward functions. By analyzing simulation statistics and rewriting the code in a closed loop, it enables robots to acquire highly dynamic physical skills—such as quadrupedal locomotion on a yoga ball—and seamlessly transfer these policies to the real world without manual reward engineering.

Beyond imperative scripts, researchers have explored \textbf{Structured Control Flow} to enhance interpretability and robustness. \textbf{LLM-as-BT-Planner}~\cite{ao2025llmbt} integrates LLMs with Behavior Trees (BTs), generating modular XML-based tree structures. This approach facilitates reactive switching and error recovery, offering a verifiable alternative to unstructured linear code, allowing agents to dynamically adapt to execution failures without replanning the entire sequence.

In the realm of optimization, a distinct branch focuses on \textbf{Language-to-Reward} generation. Instead of synthesizing direct actions, systems like \textbf{Language to Reward}~\cite{ma2023reward} decouple semantic task specification from motor execution by generating programmatic objective functions, thereby delegating the continuous trajectory optimization to low-level reinforcement learning\@. Expanding the input modalities for programmatic reasoning, \textbf{Demo2Code}~\cite{wang2023demo2code} translates visual and kinematic expert demonstrations directly into executable symbolic code. 
This approach explicitly grounds logical script generation in physical reality, merging the sample efficiency of imitation learning with the rigorous interpretability of code generation.

\subsubsection{Long-Horizon Task Decomposition}\label{sec:task_decomp}

Although foundation models provide strong semantic priors and broad common-sense knowledge, their performance often degrades in long-horizon manipulation tasks, where compounding errors in task decomposition and state tracking can lead to inconsistent plans. To mitigate this, a prominent research stream integrates LLMs with logic slovers. \textbf{LLM+P}~\cite{liu2023llmp} addresses logical hallucinations by translating natural language instructions into Planning Domain Definition Language(PDDL), delegating the reasoning process to classical solvers to guarantee feasibility. This symbolic grounding was further enhanced by \textbf{L3M+P}~\cite{agarwal2025l3mp}, which enables the solver to adaptively learn domain transition models from experience rather than relying solely on pre-defined rules.

However, classical solvers face combinatorial explosions in large-scale environments. To address scalability, \textbf{SayPlan}~\cite{sayplan2023} integrates a \textbf{Semantic Search Mechanism} anchored in 3D Scene Graphs. By hierarchically pruning the search space—filtering out irrelevant rooms or containers before planning—SayPlan enables efficient task decomposition in expansive, multi-room environments where flat PDDL planning would be computationally intractable. Complementing this, \textbf{TaPA}~\cite{wu2023tapa} grounds LLM-generated plans in physical scene constraints by conditioning on the set of objects actually present in the environment, ensuring that each planned action is executable given the current world state.

Beyond static planning, establishing a Closed-Loop Feedback mechanism also crucial for robustness. \textbf{Inner Monologue}~\cite{huang2023inner} pioneered this paradigm by injecting environmental feedback back into the LLM's context window. This allows the model to perform feedback-conditioned in-context replanning, enabling immediate plan revision after execution failure.

A distinct challenge from feedback is world state maintenance which prevents the agent from forgetting the consequences of past actions. \textbf{Statler}~\cite{yoneda2024statler} addresses long-horizon state consistency by treating planning as state-dependent code generation with explicit world-state memory, thereby reducing open-loop errors caused by forgotten object states.

\subsubsection{Grounded Multimodal Reasoning}\label{sec:multimodal_reasoning}

While B1 and B2 operate in purely symbolic or logical spaces---generating text plans, code, or PDDL specifications---they rely on downstream modules to interpret these outputs in the physical world. Grounded Multimodal Reasoning marks the transition layer where foundation models first ingest raw sensory signals (vision, proprioception) alongside language, producing representations that are semantically rich yet physically situated. The key distinction from B1/B2 is the input modality: rather than reasoning over text alone, B3 models jointly process visual observations and language within a single pretrained backbone, enabling zero-shot transfer of internet-scale knowledge to embodied contexts.

A foundational approach, \textbf{PaLM-E}~\cite{driess2023palme}, directly injects continuous state vectors and visual patches into the frozen embedding space of an LLM, treating sensor inputs as ``multimodal sentences'' to enable zero-shot referential reasoning. Complementing this implicit inference, \textbf{RoboBrain}~\cite{ji2025robobrain} focuses on structured knowledge retention, mining affordance-centric relationships from unstructured data to build a queryable semantic memory bank that mitigates the hallucination risks inherent in purely generative models.

To extend this grounded reasoning into the temporal domain, \textbf{EmbodiedGPT}~\cite{mu2023embodiedgpt} pioneers multimodal chain-of-thought by fine-tuning models on ``Observation-Thought-Action'' triplets, forcing the agent to explicitly decompose high-level instructions into hierarchical sub-goals via an Embodied CoT that conditions on visual observations at each step. Building on this, \textbf{ReAct}~\cite{yao2023react} introduces a closed-loop reasoning-acting framework where the agent alternates between generating reasoning traces and executing actions, dynamically adjusting its plan based on environmental feedback---transforming static chain-of-thought into an interactive, perception-conditioned decision process. Unlike B2's purely logical re-planning (e.g., Inner Monologue operating over text state), these approaches ground each reasoning step in the current visual observation, making them robust to perceptual ambiguity and partial observability.

As these multimodal reasoning capabilities mature, they naturally point toward the need for precise geometric grounding (B4) and predictive physical simulation (B5). The progression from semantic understanding to spatial precision is further unified in Type~I integration architectures (\S4.3), which absorb both reasoning and action generation into a single end-to-end backbone.

\subsubsection{Geometric Grounding and Spatial Intelligence}\label{sec:affordance_mapping}

While reasoning models determine \textbf{what} to do, the geometric grounding module determines exactly \textbf{where} and \textbf{how} to interact. This layer bridges the gap between high-level semantic intent and low-level geometric execution, evolving from physical constraint satisfaction through 2D affordance heatmaps to continuous 3D representations and dexterous manipulation.

A first class of methods leverages foundation models to ground semantic plans in precise physical constraints. \textbf{ManipLLM}~\cite{manipllm2024} fine-tunes VLMs to predict pixel-wise contact masks rather than generic text, ensuring stable grasping on arbitrary objects. \textbf{LLM-GROP}~\cite{llmgrop2024} formulates spatial relationship reasoning (e.g., ``left of'') as a geometric satisfaction problem to determine collision-free SE(3) poses. \textbf{A3VLM}~\cite{huang2024a3vlm} deepens this grounding by explicitly reasoning about kinematic structure---predicting joint axes and motion types (e.g., distinguishing revolute from prismatic joints), enabling manipulation of unseen articulated objects by understanding their underlying mechanisms.

% \paragraph{From 2D Affordance Heatmaps to Zero-Shot Geometric Grounding.}
Beyond these constraint-based approaches, the foundational paradigm for spatial affordance is established by \textbf{CLIPort}~\cite{shridhar2021cliport}, which introduces a \textbf{Two-Stream Architecture} disentangling semantic understanding (via CLIP) from spatial precision (via Transporter Networks). By cross-attending semantic features with rotation-invariant spatial crops, it allows agents to localize language-conditioned affordances with pixel-perfect accuracy. However, heatmap-based approaches often lack structural understanding. To address this, \textbf{PASG (Primitive-Aware Semantic Grounding)}~\cite{zhu2025pasg} introduces a closed-loop framework which automatically extracts geometric primitives such as functional axes and keypoints, employing VLMs to perform \textbf{Semantic Anchoring} that dynamically couples these primitives with task descriptions. This self-corrective mechanism ensures that the predicted affordance is not just visually plausible but geometrically aligned with the object's function (e.g., aligning with the pouring axis). Parallel to this, \textbf{Moka}~\cite{moka2024} employs \textbf{Mark-Based Visual Prompting}, overlaying actionable points directly on the image to guide the VLM, limiting the search space for open-vocabulary manipulation. To achieve zero-shot generalization, recent works pivot to \textbf{Retrieval and Distillation}. \textbf{RAM (Retrieval-Based Affordance Transfer)}~\cite{ram2025} proposes a memory-based paradigm, retrieving expert affordance masks from external databases to adapt to novel objects instantaneously. Similarly, \textbf{UAD (Unsupervised Affordance Distillation)}~\cite{tang2025uad} and \textbf{UniAff}~\cite{uniaff2025} bypass manual annotation by distilling dense affordance knowledge from foundation models, enabling the inference of actionable regions on unseen categories via internet-scale priors. These advances enable affordance prediction on novel objects without manual annotation, but remain confined to 2D pixel space.
% \paragraph{Lifting to 3D: Neural Fields, Semantic Geometry, and Efficient Representations.}
The deviation from discrete 2D affordance was marked by \textbf{NDF}~\cite{simeonov2021ndf}, which formulates object representation as continuous SE(3)-equivariant coordinate fields, enabling manipulation based on shape similarity rather than memorizing global templates. \textbf{GNFactor}~\cite{ze2023gnfactor} advances this by jointly optimizing a generalizable neural feature field and a Perceiver Transformer over a shared deep 3D voxel representation, unifying scene reconstruction with decision-making for multi-task real-robot learning. To further augment this geometric substrate with open-vocabulary semantics, \textbf{F3RM}~\cite{f3rm2023} embeds CLIP features into the neural field, allowing functional affordances (e.g., ``the handle'') to be localized via language prompts, and \textbf{D3Fields}~\cite{wang2024d3fields} extends this to dynamic scenes using DINO and CLIP features for zero-shot correspondence across unseen instances. Manipulation in 3D space requires explicit \textbf{Geometric Awareness}: \textbf{RoboSpatial}~\cite{robospatial2025} and \textbf{GeoManip}~\cite{geomanip2025} lift affordance to 3D point clouds, modeling geometric curvature and collision landscapes for kinematic feasibility. Addressing the latency of neural fields, \textbf{GraspSplats}~\cite{graspsplats2025} embeds grasp affordances directly into 3D Gaussian Splatting primitives for real-time differentiable interaction, while \textbf{CoPa}~\cite{copa2024} decomposes objects into task-relevant geometric parts to ground physics at the fine-grained component level.
% \paragraph{Constraint-Based Geometric Reasoning and Robust Tracking.}
To translate these representations into actionable plans, the focus shifts from observation to constraint satisfaction. \textbf{VoxPoser}~\cite{voxposer2023} utilizes LLMs to synthesize 3D value maps, effectively grounding abstract language instructions (e.g., ``avoid the vase'') into dense potential fields that guide motion planners through safe trajectories. For complex dynamic interactions, \textbf{ReKep}~\cite{rekep2024} advances this reasoning by defining manipulation tasks not as fixed trajectories, but as systems of time-varying geometric constraints between keypoints. This formulation allows agents to solve for actions that satisfy physical relations (e.g., keeping a cup upright) robustly against disturbances. Complementing this, \textbf{FoundationPose}~\cite{wen2024foundationpose} provides the robust 6D pose tracking essential for contact-rich tasks.

Beyond simple parallel-jaw grasping, \textbf{DexGraspNet}~\cite{wang2023dexgraspnetlargescaleroboticdexterous} provides massive-scale priors for multi-finger contact. Building on this, \textbf{UniDexGrasp++}~\cite{unidexgrasp2024} advances the policy learning itself, employing a geometry-aware curriculum to master universal grasping across thousands of object categories. Furthermore, \textbf{GraspGPT}~\cite{tang2023graspgpt} introduces \textbf{Language-Conditioned Dexterity}, leveraging LLMs to interpret abstract instructions (e.g., ``grasp the pen for writing'') and generate task-specific functional grasp poses, effectively bridging the gap between semantic reasoning and high-dimensional finger control. Ultimately, these capabilities converge in \textbf{OmniManip}~\cite{pan2025omnimanip}, which synthesizes field-based and constraint-based modalities into an \textbf{Omni-Generalizable} architecture, enabling agents to autonomously reason about and manipulate diverse objects in open-world environments.

\subsubsection{Generative Dynamics and Physical Imagination}\label{sec:generative_dynamics}

The ability to predict and simulate physical consequences before acting represents the frontier of embodied cognition. Unlike reactive systems, generative dynamics models enable robots to \textbf{imagine} future outcomes and plan through mental simulation---marking the transition from trial-and-error learning to physics-informed foresight.

The foundational approach is \textbf{latent dynamics world models}, which compress observations into compact state spaces for autoregressive rollouts. \textbf{Daydreamer}~\cite{wu2023daydreamer} first demonstrates that real robots can learn visuomotor skills from scratch within hours via online RSSM-based latent imagination, and moving beyond pixel reconstruction, \textbf{TD-MPC2}~\cite{hansen2024tdmpc2} jointly optimizes value learning and MPC trajectory search in a purely task-driven latent space, while \textbf{DynaMo}~\cite{cui2024dynamo} introduces dynamics pre-training to generate control-oriented representations. At the frontier, \textbf{V-JEPA2}~\cite{assran2025vjepa} abandons reconstruction altogether, predicting future states in an abstract embedding space---representing the leap from \emph{reconstructive} to \emph{predictive} world modeling.

While latent models offer efficiency, they sacrifice visual fidelity. A parallel stream leverages large-scale \textbf{video generation models} as explicit physical forward simulators, treating visual sequences as the richest carriers of contact and dynamics priors. \textbf{UniPi}~\cite{du2023unipi} pioneers this by reformulating policy as video generation---synthesizing future trajectories via text-conditioned diffusion, from which inverse dynamics extracts actions. \textbf{VILA}~\cite{du2024vila} extends this to planning verification by imagining pixel-level consequences before execution, and \textbf{GVM}~\cite{gvm2024} closes the loop by using synthesized frames as visual sub-goals for conditioned exploration. More recent works tighten the prediction-action coupling: \textbf{VPP}~\cite{hu2025vpp} learns inverse dynamics \emph{within} the diffusion model itself, and \textbf{UWM}~\cite{zhu2025uwm} unifies video and action diffusion in a single transformer with independent diffusion timesteps. For compositional reasoning, \textbf{RoboDreamer}~\cite{zhou2024robodreamer} supports pre-interaction visual imagination to evaluate manipulation feasibility. Scaling toward foundation-level simulation, \textbf{UniSim}~\cite{yang2024unisim} learns interactive real-world simulators from heterogeneous video, \textbf{Genie}~\cite{bruce2024genie} extracts controllable environments from unlabeled internet video in a fully unsupervised manner, and \textbf{IRASim}~\cite{zhu2025irasim} generates fine-grained robot-object interaction videos. To inject 3D geometric rigidity, \textbf{GWM}~\cite{lu2025gwm} combines Diffusion Transformers with 3D Gaussian-aware VAEs for multi-view consistency under complex contact. At the industrial frontier, \textbf{Cosmos Policy}~\cite{nvidia2025cosmos} projects actions into latent frames co-participating in diffusion generation, \textbf{DreamDojo}~\cite{gao2026dreamdojo} pre-trains on 44k hours of egocentric human video for cross-embodiment generalization, and \textbf{Ctrl-World}~\cite{guo2026ctrlworld} enables policy evaluation and improvement entirely within imagination via controllable policy-in-the-loop rollouts.

However, purely data-driven world models lack intrinsic adherence to physical laws, causing cumulative implausibility over long horizons. \cite{schulze2025lagrangian} embed Lagrangian mechanics into dynamics models for 7-DoF manipulators, enforcing energy conservation, while \textbf{Particle-Grid Neural Dynamics}~\cite{zhang2025particlegrid} constructs particle-mesh GNN dynamics for deformable objects under large deformations. \textbf{STORM}~\cite{lin2025storm} introduces generative MCTS where a foundation model serves as action prior and a video world model as simulation engine with UCB pruning to discipline hallucinations, and \textbf{LUMOS}~\cite{nematollahi2025lumos} embeds language instructions into the prediction loop for multi-stage task compliance.

\subsection{Policy Learning: The Embodied Cerebellum}\label{sec:policy_learning_low_level}

Low-level motion mapping is the core mechanism by which robots transform perceptual representations into physical control commands. The fundamental objective is to learn a policy $\pi$ that maps high-dimensional observations $\mathcal{S}$---spanning visual, proprioceptive, and tactile modalities---to actuator action spaces $\mathcal{A}$. As the field evolves from single-task behavior cloning to generalist robot learning, the central challenge has shifted from trajectory fitting to modeling complex multimodal action distributions where a single state may admit multiple valid responses. We organize recent advances into four paradigms based on how the action space is formulated: \emph{Explicit Continuous Policies} that treat control as deterministic regression, \emph{Discrete Action Policies} that reframe control as sequence classification, \emph{Iterative Generative Policies} that model complex distributions via diffusion and flow matching, and \emph{Efficient Policy Architectures} that balance computational efficiency with multi-task inference speed. Figure~\ref{fig:policy_timeline} traces the chronological evolution of these paradigms, and Figure~\ref{fig:policy_structure} illustrates their architectural relationships.
%% TODO: Insert \ref{fig:policy_timeline} — timeline figure here
%% TODO: Insert \ref{fig:policy_structure} — structure/taxonomy figure here

\begin{table*}[t!]
\centering
\caption{Taxonomy of Policy Learning Paradigms for Embodied Manipulation. \textbf{P1}\,=\,Explicit Continuous (1-step forward pass, MSE/L1 regression); \textbf{P2}\,=\,Discrete Action ($T$-step autoregressive, cross-entropy); \textbf{P3}\,=\,Iterative Generative ($K$-step denoising/flow). Data Tier indicates the dataset level (\S\ref{sec:eai_datasets}) that primarily supports each paradigm ($\mathcal{K}_1$\,=\,Skill, $\mathcal{K}_2$\,=\,Task, $\mathcal{K}_3$\,=\,Foundation). Abbrev.: Percept.\,=\,Perception; Regr.\,=\,Regression; Resid.\,=\,Residual; Discr.\,=\,Discretization; Enhanc.\,=\,Enhancement; Consist.\,=\,Consistency; Diff.\,=\,Diffusion; refine.\,=\,refinement; Cont.\,=\,Continuous.}\label{tab:policy_paradigms}
\renewcommand{\arraystretch}{1.15}
\setlength{\tabcolsep}{2.5pt}
\scriptsize
\begin{tabular}{
  c
  c
  c
  c
  c
  c
  c
  p{0.18\linewidth}
}
\toprule
\textbf{ID} &
\textbf{Data Tier} &
\textbf{Sub-direction} &
\textbf{Action Space} &
\textbf{Scale} &
\textbf{Input} &
\textbf{Multimodality} &
\textbf{Representative Methods} \\
\midrule
\multirow{3}{*}{\textbf{P1}}
 & \multirow{3}{*}{$\mathcal{K}_1/\mathcal{K}_2$} & Percept.-Action Regr. & Cont. & {$<$100M} & 2D / 3D & Mode-averaging & BC-Z~\cite{jang2022bc}, MVP~\cite{xiao2022mvp}, PointPolicy~\cite{haldar2025pointpolicy} \\
 & & Temporal Chunking & Cont.\ chunk & {100M--1B} & 2D / multi-view & Mode-averaging & DT~\cite{chen2021decisiontransformer}, ACT~\cite{zhao2023learning}, Mobile ALOHA~\cite{fu2024mobile}, BAKU~\cite{haldar2024baku} \\
 & & Recovery \& Transfer & Cont.\ / resid. & {100M--1B} & 2D / video & Goal-conditioned & MimicPlay~\cite{wang2023mimicplay}, Iter.\ Resid.~\cite{chi2022iterativeresidualpolicy}, Motion Tracks~\cite{ren2025motion_tracks}, UniSkill~\cite{kim2025uniskill} \\
\midrule
\multirow{3}{*}{\textbf{P2}}
 & \multirow{3}{*}{$\mathcal{K}_2/\mathcal{K}_3$} & Uniform \& Spatial Discr. & Fixed bins & {1B--55B} & 2D / voxel & Categorical modes & RT-1~\cite{brohan2023rt1}, Gato~\cite{reed2022generalistagent}, PerAct~\cite{shridhar2022peract}, Act3D~\cite{gervet2023act3d}, RVT~\cite{goyal2023rvt} \\
 & & Learned Quantization & Codebook & {100M--1B} & 2D & Codebook modes & BeT~\cite{shafiullah2022bet}, VQ-BeT~\cite{lee2024vqbet}, CARP~\cite{gong2025carp} \\
 & & Action-as-Language & LM tokens & {1B--10B+} & 2D / multi-modal & LM categorical & VIMA~\cite{jiang2023vima}, GR-1~\cite{wu2023gr1}, Seer~\cite{tian2024seer}, ARP~\cite{zhang2025arp}, ICIL~\cite{fu2025icil} \\
\midrule
\multirow{4}{*}{\textbf{P3}}
 & \multirow{4}{*}{$\mathcal{K}_1/\mathcal{K}_3$} & Diffusion Foundations & Denoised cont. & {$<$500M} & 2D & Score-based refine. & Diff.\ Policy~\cite{bekris2023diffusion}, Diffuser~\cite{diffuser_2022}, Beso~\cite{beso_2023}, StructDiff.~\cite{structdiffusion_2023} \\
 & & 3D Geometric Enhanc. & Denoised cont. & {$<$500M} & 3D / point cloud & Equivariant refine. & DP3~\cite{dp3_2024}, iDP3~\cite{idp3_2024}, EquiBot~\cite{equibot_2025}, ManiGaussian~\cite{lu2024manigaussian} \\
 & & Flow \& Acceleration & Flow cont. & {$<$500M} & 2D / 3D & OT flow & Consistency Policy~\cite{prasad2024consistency}, RecFlow~\cite{xue2025recflow}, SeFA~\cite{sefa_policy_2024} \\
 & & Generalist Scaling & Denoised cont. & {1B--8B} & 2D / multi-view & Score-based refine. & Octo~\cite{octomodelteam2024}, RDT-1B~\cite{rdt1b_2024}, FLOWER~\cite{reuss2025flower} \\
\bottomrule
\end{tabular}
\end{table*}

\subsubsection{Explicit Continuous Policies}\label{sec:explicit_continuous_policies}

The most direct path from perception to motor command is \emph{explicit continuous regression}: a single forward pass maps observations to actions without discrete tokenization or iterative denoising, trading distributional expressiveness for inference simplicity and deterministic reproducibility. The foundational design separates \emph{perception} from \emph{action regression}: a pre-trained visual encoder extracts task-relevant features, and a lightweight action head maps them to motor commands. \textbf{BC-Z}~\cite{jang2022bc} demonstrates that this design scales to multi-task zero-shot generalization, and \textbf{MVP}~\cite{xiao2022mvp} shows that large-scale masked visual pre-training yields representations so robust that the policy reduces to a simple MLP~\cite{Radosavovic2022}. While these methods primarily operate on 2D images, \textbf{PointPolicy}~\cite{haldar2025pointpolicy} extends explicit regression to 3D point clouds, leveraging geometric information to handle occlusion and spatial ambiguity.

Single-step prediction, however, sacrifices the temporal coherence required for smooth physical execution---each action is generated independently, ignoring the trajectory context that precedes it. \textbf{Decision Transformer}~\cite{chen2021decisiontransformer} addresses this by reframing control as autoregressive sequence modeling with regression objectives. The \textbf{ACT} line of work~\cite{zhao2023learning} operationalizes this for real robots via \emph{action chunking} and CVAE-based modeling, predicting short-horizon sequences to mitigate jitter and latency---a formulation that extends naturally to bimanual coordination (\textbf{Bi-ACT}) and whole-body control (\textbf{Mobile ALOHA})~\cite{fu2024mobile}. For multi-task efficiency, \textbf{BAKU}~\cite{haldar2024baku} combines CNN-based granular feature extraction with a transformer policy, avoiding the overhead of full-scale sequence transformers.

Yet even with temporal structure, explicit regression remains inherently open-loop within each prediction step, compounding errors under distribution shift. Two complementary strategies mitigate this fragility. For \emph{long-horizon reusability}, \textbf{MimicPlay}~\cite{wang2023mimicplay} learns goal-conditioned behaviors from play-style data that compose at test time, while \textbf{Iterative Residual Policy}~\cite{chi2022iterativeresidualpolicy} repeatedly predicts residual corrections to recover from execution drift. For \emph{cross-embodiment transfer}, \textbf{Motion Tracks}~\cite{ren2025motion_tracks} learns a unified motion representation supporting rapid adaptation from limited demonstrations, and \textbf{UniSkill}~\cite{kim2025uniskill} distills embodiment-agnostic skill representations from large-scale unlabeled cross-embodiment video, enabling human video prompts to guide robot policies without aligned datasets.

\subsubsection{Discrete Action Policies}\label{sec:discrete_action_policies}

In contrast to explicit continuous policies, \textbf{Discrete Action Policies} reframe robot control as a \textbf{sequence modeling problem}---quantizing continuous actions into discrete tokens and optimizing via categorical cross-entropy, which naturally captures multi-modal action distributions that unimodal regression cannot represent. Foundational approaches employ \textbf{uniform discretization}: RT-1 maps continuous dimensions to 256 uniform bins for large-scale manipulation~\cite{brohan2023rt1}, while Gato uses $\mu$-law encoding to tokenize diverse modalities into a unified sequence~\cite{reed2022generalistagent}. Extending this to 3D, PerAct classifies actions over voxelized grids~\cite{shridhar2022peract}, Act3D selects keypoints on 3D feature fields~\cite{gervet2023act3d}, and RVT predicts actions via multi-view heatmaps~\cite{goyal2023rvt}. However, fixed binning forces a fundamental trade-off between precision and vocabulary size. \textbf{Learned quantization} methods resolve this: BeT defines action modes via k-means clustering~\cite{shafiullah2022bet}, VQ-BeT learns compact codebooks through VQ-VAE~\cite{lee2024vqbet}, and CARP introduces coarse-to-fine hierarchical prediction for fine-grained control without vocabulary explosion~\cite{gong2025carp}.
% Action-as-Language
With robust tokenization in place, the paradigm culminates in \textbf{Action-as-Language} models that treat control as semantic token generation. \textbf{VIMA}~\cite{jiang2023vima} pioneers this by interleaving visual objects with text instructions for multimodal prompt conditioning. Unified-IO 2 and GR-1 push toward ``Any-to-Any'' generation, unifying vision, language, and action in a single autoregressive stream~\cite{lu2023unifiedio2, wu2023gr1}, while \textbf{Seer}~\cite{tian2024seer} validates that the discrete architecture itself---not massive language pre-training---drives the performance gains. \textbf{ARP}~\cite{zhang2025arp} further demonstrates that chunk-level autoregressive action generation yields a universal policy architecture generalizing across diverse robots and task configurations.
% complex control demands
Scaling this formulation to \textbf{complex control demands} reveals both its versatility and its unique affordances. Dense Policy introduces bidirectional autoregressive learning for denser supervision signals~\cite{su2025densepolicy}, and Humanoid Policy extends tokenization to high-dimensional whole-body control~\cite{qiu2025humanoidpolicyhuman}. Crucially, the alignment with next-token prediction uniquely enables \textbf{In-Context Imitation Learning}: agents adapt to unseen tasks solely by conditioning on demonstration prompts, treating imitation as sequence completion without gradient updates~\cite{fu2025icil}---an emergent capability unavailable to continuous or diffusion-based policies.

\subsubsection{Iterative Generative Policies}\label{sec:iterative_generative_policies}

Bridging the precision of continuous regression and the multi-modal expressiveness of discrete approaches, \textbf{Iterative Generative Policies} formulate robot control as a conditional denoising process. The foundational work, \textbf{Diffusion Policy}~\cite{bekris2023diffusion}, adapts DDPMs to visuomotor control by learning the score function of the action distribution, avoiding mode-averaging pitfalls~\cite{diffusion_bc_2023}. \textbf{Beso}~\cite{beso_2023} extends this with goal-conditioned score-based density estimation. Beyond immediate control, diffusion naturally extends to long-horizon planning: \textbf{Diffuser}~\cite{diffuser_2022} treats planning as probabilistic inpainting, \textbf{Decision Diffuser}~\cite{decision_diffuser_2023} decouples generation from constraint satisfaction, and \textbf{StructDiffusion}~\cite{structdiffusion_2023} constructs physically valid multi-object structures from language prompts.

To address the geometric ambiguity of 2D image-based diffusion, a body of work incorporates explicit 3D representations. \textbf{DP3} and its humanoid extension \textbf{iDP3} integrate point cloud processing with diffusion heads, enhancing generalization to viewpoint changes~\cite{dp3_2024, idp3_2024}. \textbf{ManiGaussian}~\cite{lu2024manigaussian} utilizes dynamic Gaussian Splatting for high-fidelity manipulation, while \textbf{RDP}~\cite{xue2025reactive} introduces reactive diffusion for visual-tactile modalities. Pushing further, \textbf{EquiBot}~\cite{equibot_2025} and \textbf{Diffusion-EDFs}~\cite{diffusion_edfs_2024} enforce SE(3)/SIM(3) equivariance for zero-shot generalization to rigid transformations, refined by spherical Fourier representations~\cite{spherical_fourier_2025} and geometric algebra for dexterous grasping~\cite{sun2025hybrid_pga,zhong2025gagrasp}.

However, iterative denoising incurs prohibitive inference latency, motivating a paradigm shift toward \textbf{Flow Matching}, which learns velocity fields transporting noise to data along straight trajectories with far fewer ODE steps. \textbf{AdaFlow}~\cite{hu2024adaflow} pioneered this direction with a variance-adaptive ODE solver that automatically reduces to single-step generation for unimodal distributions. \textbf{RFMP} and \textbf{SRFMP}~\cite{braun2024riemannian, ding2025fast} extend flow matching to Riemannian manifolds with LaSalle-based stability guarantees for end-effector pose generation. Two strategies then emerged to push toward single-step inference: \textbf{Consistency Policy}~\cite{prasad2024consistency} distills a teacher Diffusion Policy's ODE trajectory into few-step generation, while \textbf{RecFlow}~\cite{xue2025recflow} and \textbf{SeFA}~\cite{sefa_policy_2024} rectify the flow itself---the latter reducing latency by over 98\% through selective alignment against expert demonstrations. Subsequent work explores complementary axes: \textbf{FlowPolicy}~\cite{zhang2025flowpolicy} achieves single-step 3D point-cloud-conditioned inference with 7$\times$ speedup; \textbf{Streaming Flow}~\cite{jiang2025streaming} streams actions on-the-fly during ODE integration for receding-horizon control; \textbf{ReSeFlow}~\cite{wang2025reseflow} injects rectified-flow efficiency into SE(3)-equivariant models; \textbf{ManiFlow}~\cite{yan2025maniflow} scales to 1--2 step dexterous generation across single-arm, bimanual, and humanoid platforms via DiT-X; and \textbf{Guided Flow Policy}~\cite{tiofack2026guided} extends flow matching to offline RL with value-aware selective regularization.

Iterative generative policies have become the backbone for cross-embodiment foundation models. \textbf{Octo}~\cite{octomodelteam2024} establishes the open-source baseline: a transformer-based diffusion policy pretrained on 800k trajectories across 25 robot configurations, supporting language and goal-image conditioning with few-hour finetuning. \textbf{RDT-1B}~\cite{rdt1b_2024} scales this to 1.2B parameters with a Physically Interpretable Unified Action Space across 46 datasets, producing emergent zero-shot generalization and state-of-the-art bimanual performance. Extensions span multimodal sensing (\textbf{DP-CA}~\cite{kang2025dpca}, force-feedback for contact-rich tasks), efficient foundation-model scaling (\textbf{FLOWER}~\cite{reuss2025flower}, 950M parameters competitive with billion-scale models at 200 GPU hours), cross-morphology transfer (\textbf{Grasp2Grasp}~\cite{zhong2025grasp2grasp}, Schr\"{o}dinger Bridge for dexterous grasp translation), adaptive inference compute (\textbf{DA-SIP}~\cite{chun2025dynamictesttimecomputescaling}, 2.6--4.4$\times$ reduction via difficulty-aware solver selection), and training-free closed-loop correction (\textbf{DCDP}~\cite{wu2026dcdp}, injecting real-time dynamics into pre-committed action chunks).\@

\subsection{Integration Models}

RT-2\cite{brohan2023rt2} marks the practical onset of modern robotic integation models by showing that language-conditioned action tokenization can unify semantic understanding and low-level control in an end-to-end policy.
(TypeI) systems then scaled this paradigm through larger backbones and broader datasets (e.g., OpenVLA, RoboVLMs, UniVLA, NORA-1.5, ReconVLA), improving task breadth and cross-domain transfer\cite{kim2024openvla,li2026matters,bu2025univla,hung2025nora15,song2025reconvla}.
As autoregressive latency became a deployment bottleneck in contact-rich settings, the field moved toward two structural directions—time-scale decoupling (TypeII) and spatially grounded state modeling (TypeIII)\cite{belkhale2024rth,chen2025fis,zhen20243dvla,qu2025spatialvla}; Table~\ref{tab:embodied_models_2026_v4} summarizes this transition from 2023 to 2026. We note that these three types are not mutually exclusive: a given model may exhibit characteristics of multiple categories (e.g., a world model with a hierarchical reasoning loop), and we classify each by its primary architectural contribution.

\subsubsection{Type I\@: Unified End-to-End Foundation Models}

The core philosophy of Type I models is to leverage the immense capacity of large-scale architectures---primarily Transformers and, more recently, State Space Models---to absorb vast amounts of diverse robotic data and generalize across tasks and embodiments. By treating actions as just another modality (often tokenized alongside text and image patches), these models inherit the scaling properties of their vision-language predecessors. The seminal work in this category is \textbf{RT-2}~\cite{brohan2023rt2}, which demonstrated that a Vision-Language Model (VLM) could be co-fine-tuned on robotic trajectory data to output discrete action tokens. Building upon this foundation, the \textbf{Open X-Embodiment} initiative introduced \textbf{RT-X}~\cite{padalkar2023open}, scaling the architecture across multiple robot morphologies. To democratize access, \textbf{OpenVLA}~\cite{kim2024openvla} was introduced as a fully open-source 7B parameter vision-language-action model, while \textbf{RoboFlamingo}~\cite{li2023roboflamingo} adapted the Flamingo architecture for robotic control. Recent efforts continue to push scaling boundaries; \textbf{InternVLA-A1}~\cite{cai2026internvlaa1} unifies understanding and generation for manipulation, and \textbf{GR00T}~\cite{nvidia2025grootn1} extends the foundation model paradigm to generalist humanoid robots.

As the field progressed, researchers sought to address the limitations of standard autoregressive generation, particularly regarding inference speed and the handling of continuous action spaces. To bridge the gap between discrete tokens and smooth physical execution, \textbf{$\pi_0$}~\cite{black2024pi0} introduced a flow-matching architecture combining VLM semantics with continuous diffusion control. Similarly, \textbf{HybridVLA}~\cite{liu2025hybridvla} explores a collaborative paradigm, utilizing autoregression for high-level semantic steps and diffusion for low-level continuous execution, while \textbf{F1-VLA}~\cite{lv2025f1vla} bridges understanding and generation to actions. To improve discrete action modeling, \textbf{VQ-VLA}~\cite{wang2025vqvla} demonstrates that scaling vector-quantized action tokenizers significantly reduces quantization errors. Beyond standard transformers, \textbf{RoboMamba}~\cite{liu2024robomamba} replaces self-attention with State Space Models (SSMs), achieving linear time complexity that is crucial for processing long-horizon multimodal sequences. 

While scaling up model capacity unlocks emergent reasoning capabilities, the prohibitive latency of massive architectures has catalyzed a parallel push towards affordable, edge-deployable models. To overcome this computational bottleneck, inference acceleration techniques have gained significant traction: \textbf{DeeR-VLA}~\cite{yue2024deervla} introduces a dynamic early-exit mechanism that halts transformer computation once action confidence is sufficient, \textbf{MoLe-VLA}~\cite{zhang2025molevla} achieves similar acceleration through dynamic layer-skipping via a Mixture-of-Layers design, and \textbf{VLA-Cache}~\cite{xu2025vlacache} optimizes autoregressive caching for real-time execution. At the extreme end of compression, \textbf{BitVLA}~\cite{wang2025bitvla} pushes quantization to 1-bit weights. To make the adaptation of these models more accessible, \textbf{VLA-Adapter}~\cite{wang2025vlaadapter} introduces parameter-efficient tuning paradigms specifically tailored for tiny-scale models like \textbf{TinyVLA}~\cite{wen2025tinyvla}, \textbf{SmolVLA}~\cite{shukor2025smolvla}, \textbf{Edge VLA}~\cite{budzianowski2025edgevla}, and \textbf{Evo-1}~\cite{lin2025evo1}, while \textbf{OpenVLA-OFT}~\cite{kim2025openvlaoft} optimizes the speed and success rate of fine-tuning full-scale models for downstream tasks. Notably, while inference-side compression has received extensive attention, the training-side computational asymmetry remains underexplored: pre-training a 7B model on OXE-scale data typically demands thousands of GPU-hours, yet the cost-performance frontier between compact models (e.g., TinyVLA, SmolVLA) and full-scale architectures lacks standardized benchmarking, hindering principled resource allocation for both academic and industrial practitioners.

To better align model outputs with complex environments and enhance perceptual grounding, \textbf{ReconVLA}~\cite{song2025reconvla} introduces auxiliary reconstructive objectives, forcing the model to explicitly encode 3D spatial structures. \textbf{BridgeVLA}~\cite{li2025bridgevla} aligns inputs and outputs specifically for 3D manipulation, while \textbf{NORA-1.5}~\cite{hung2025nora15} incorporates world-model preference rewards. Furthermore, the modality spectrum is expanding: \textbf{Interleave-VLA}~\cite{fan2026interleave} processes interleaved image-text sequences for complex multi-step reasoning, and \textbf{VLAS}~\cite{zhao2025vlas} natively integrates speech instructions to bypass text-transcription latency. While Type I models offer a compelling vision of end-to-end embodied intelligence, they often struggle with tasks requiring precise 3D spatial reasoning or complex, long-horizon planning. Complementing these architectural advances, \textbf{VLANeXt}~\cite{wu2026vlanext} provides a systematic ablation of the Type~I design space, distilling practical recipes across foundational components, perception, and action modeling. These limitations motivate the development of Type II (Hierarchical Temporal Decoupling) and Type III (Spatial Grounding \& Predictive World Models) architectures, which explicitly incorporate structural priors to handle the nuances of the physical world.


<!- llava-vla Rethinking the Practicality of Vision-language-action Model: A Comprehensive Benchmark and An Improved Baseline>

\begin{table*}[t!]
\centering
\caption{Comprehensive Taxonomy of Embodied Foundation Models (2023--2026).
\textbf{Brain (B):} B1: Semantic Logic, B2: L-H Decomp, B3: Grounded MM Reas., B4: Geo.\ Grounding, B5: Spatial Intel.
\textbf{Policy (P):} P1: Expl.\ Continuous, P2: Discrete Action, P3: Iter.\ Generative.
\textbf{Ctrl.~Hz}: Low~$<$10\,Hz; High~20--100\,Hz; RT~$\geq$100\,Hz; (Step)~=~reasoning-loop; (Hier.)~=~hierarchical slow/fast.
\textbf{Hardware}: Div.\ = Diverse multi-robot; Prop.\ = Proprietary.
\textbf{Data Form}: Real Traj.\ = real robot trajectories; Web Vid.\ = internet video; Syn.\ = synthetic.}\label{tab:embodied_models_2026_v4}
\renewcommand{\arraystretch}{1.2}
\setlength{\tabcolsep}{3pt}
\scriptsize

\begin{tabular}{
  p{0.14\linewidth}
  p{0.08\linewidth}
  p{0.17\linewidth}
  p{0.08\linewidth}
  p{0.13\linewidth}
  p{0.168\linewidth}
  p{0.205\linewidth}
}
\toprule
\textbf{Model} &
\textbf{Taxonomy} &
\textbf{Backbone \& Params} &
\textbf{Ctrl.\ Hz} &
\textbf{Hardware} &
\textbf{Primary Data Scale} &
\textbf{Eval: Sim \& Real} \\
\midrule

% --- TYPE I ---
\multicolumn{7}{c}{\textbf{Type I\@: Unified End-to-End Foundation Models}} \\
\midrule

\textbf{RT-2}~\cite{brohan2023rt2} & B3, P2 & PaLI-X (55B) & Low & Google Robot & 130K Real Traj.\ + Web Vid. & Real: 700+ manipulation tasks \\

\textbf{$\pi$0}~\cite{black2024pi0} & B3, P3 & PaliGemma (3B) & High & Div.\ (6 types) & $\sim$10K Real Traj.\ (OXE+Prop.) & Real: 7 dexterous tasks \\

\textbf{RoboVLMs}~\cite{li2026matters} & B3, P1 & KosMos / LLaVA (2B+) & Low & Div. & Multi-task Real Traj.\ (CALVIN) & Sim: CALVIN (ABC$\rightarrow$D) \\

\textbf{OpenVLA}~\cite{kim2024openvla} & B3, P2 & Llama-2 + SigLIP (7B) & Low & Franka+WidowX & 970K Real Traj.\ (OXE) & Real: 29 manipulation tasks \\

\textbf{UniVLA}~\cite{bu2025univla} & B3, P2 & Emu3 (8.5B) & Low & Div. & OXE Real Traj.\ + Web Vid. & Sim: LIBERO (95.5\%) | Real: 10 cross-embodiment tasks \\

\textbf{NORA-1.5}~\cite{hung2025nora15} & B3, P3 & NORA (7B) & Low & Div. & OXE Real Traj.\ + Pref.\ data & Sim: SimplerEnv, LIBERO \\

\textbf{Interleave-VLA}~\cite{fan2026interleave} & B3, P2 & VILA-U (7B) & Low & Div. & Interleaved OXE Real Traj. & Real: 10+ manipulation tasks \\

\textbf{RDT-2}~\cite{liu2026rdt2} & B3, P3 & Qwen2.5-VL (7B) & High & ALOHA & 10K h Real Traj.\ (UMI) & Real: 10+ bimanual tasks \\

\textbf{ReconVLA}~\cite{song2025reconvla} & B3, P2 & LLaVA-7B (Qwen2-7B, $\sim$7B) & Low & AgileX PiPer & 2M samples (100K+ traj., BridgeV2/LIBERO) & Sim: CALVIN (ABCD$\rightarrow$D) | Real: 5 manipulation tasks \\

\midrule

% --- TYPE II ---
\multicolumn{7}{c}{\textbf{Type II\@: Hierarchical \& Modular Reasoning Models}} \\
\midrule

\textbf{RT-H}~\cite{belkhale2024rth} & B2, P2 & PaLI-X (7B+) & Low & Google Robot & Kitchen Real Traj. & Real: 8 manipulation tasks \\

\textbf{OK-Robot}~\cite{liu2024okrobot} & B1, P1 & OWL-ViT + CLIP & Low (Step) & Stretch & Pre-trained (Zero-shot) & Real: 171 mobile tasks \\

\textbf{CoT-VLA}~\cite{zhao2025cotvla} & B2, P2 & VILA-U (7B) & Low (Step) & Div. & OXE Real Traj.\ + EPIC & Real: 10 long-horizon tasks \\

\textbf{ThinkAct}~\cite{huang2025thinkact} & B1, P2 & Qwen2.5-VL (7B) & Low (Step) & Franka & OXE + RoboVQA Real Traj. & Sim: CALVIN | Real: 10 long-horizon tasks \\

\textbf{FiS-VLA}~\cite{chen2025fis} & B2, P1 & Llama-2 (7B) & RT (Hier.) & Franka & OXE Real Traj. & Real: 6 reactive tasks \\

\textbf{HAMSTER}~\cite{li2025hamster} & B5, P1 & VLM + 3D Policy & Low & Div. & Off-domain Web Vid. & Real: 8 cross-embodiment tasks \\

\textbf{Hi Robot}~\cite{shi2025hirobot} & B2, P3 & VLM + $\pi$0 & Low (Hier.) & Div. & Syn.\ data & Real: 10 manipulation tasks \\

\textbf{GO-1}~\cite{agibot2025} & B2, P3 & ViLLA (7B+) & Low & Prop.\ (AgiBot) & 1M Real Traj.\ (AgiBot) & Real: 217 bimanual tasks \\

\textbf{GR00T~N1}~\cite{nvidia2025grootn1} & B2, P3 & Eagle2-VLM + Flow Policy ($\sim$3B) & Low (Hier.) & Humanoid (Div.) & Syn.\ + Real Traj.\ (Isaac Sim) & Sim: Isaac Sim | Real: 10+ humanoid tasks \\

\midrule

% --- TYPE III ---
\multicolumn{7}{c}{\textbf{Type III\@: Spatial Grounding \& Predictive World Models}} \\
\midrule

\textbf{GR-2}~\cite{cheang2024gr2} & B5, P3 & GPT-Style CVAE & Low & Franka & 38M Web Vid. & Sim: CALVIN | Real: 100+ manipulation tasks \\

\textbf{UWM}~\cite{zhu2025uwm} & B5, P3 & DiT (12L, 768d) & High & Franka & DROID + Vid.\ Co-train & Sim: LIBERO | Real: 5 tasks (92\% SR) \\

\textbf{V-JEPA 2}~\cite{assran2025vjepa} & B5, P1 & ViT-g (1B) & Low & Div. & 22M Web Vid.\ (VideoMix) & Sim: Habitat (Zero-shot plan) \\

\textbf{DreamVLA}~\cite{zhang2025dreamvla} & B5, P3 & Diffusion WM (7B) & Low & Franka & CALVIN Real Traj. & Sim: CALVIN | Real: 10 manipulation tasks \\

\textbf{3D-VLA}~\cite{zhen20243dvla} & B5, P3 & 3D-LLM / Flan-T5-XL (3B) & Low & Franka & 3D-Lang-Action Real Traj. & Sim: RLBench, CALVIN | Real: 10 3D tasks \\

\textbf{SpatialVLA}~\cite{qu2025spatialvla} & B5, P2 & PaLiGemma 2 (4B) & Low & Franka+WidowX & OXE + RH20T Real Traj. & Sim: RLBench | Real: 10+ spatial tasks \\

\textbf{BridgeVLA}~\cite{li2025bridgevla} & B4, P2 & VLM Backbone & Low & Franka & Bridge / OXE Real Traj. & Real: 10 3D tasks \\

\textbf{RoBridge}~\cite{zhang2025robridge} & B4, P1 & VLM Planner & Low & WidowX & Syn.\ + Real Traj. & Real: 12 manipulation tasks \\

\textbf{PointVLA}~\cite{li2026pointvla} & B5, P1 & Qwen2-VL (2B) & Low & Franka & OXE Real Traj. & Real: 4 3D tasks \\

\bottomrule
\end{tabular}
\end{table*}

\subsubsection{Type II\@: Hierarchical \& Modular Reasoning Models}

Type II models emerge from a fundamental systems-level observation: as end-to-end Type~I models scale in parameter count to achieve broad semantic generalization, the computational latency of autoregressive decoding becomes fundamentally misaligned with the high-frequency control requirements of contact-rich manipulation~\cite{belkhale2024rth,chen2025fis,xu2025a0}. This paradigm retains the rich priors of foundation models while restructuring the policy stack, explicitly decoupling \emph{slow semantic deliberation} from \emph{fast sensorimotor execution}. In this hierarchical formulation, high-level modules are dedicated to task decomposition, intention tracking, and failure recovery, while low-level policies maintain reactive, high-bandwidth actuation under strict real-time deadlines~\cite{zhao2025cotvla,huang2025thinkact,zhang2025robridge}. This architectural shift reframes the primary bottleneck from pure model capacity to the \emph{time-scale coordination} across heterogeneous reasoning and control loops.

A prominent trajectory within this paradigm focuses on explicit reasoning and Chain-of-Thought (CoT) processes. Early hierarchical action decomposition, as seen in RT-H~\cite{belkhale2024rth}, has rapidly evolved into sophisticated visual latent planning and explicit reasoning supervision. Models such as CoT-VLA~\cite{zhao2025cotvla}, ThinkAct~\cite{huang2025thinkact}, and ECoT~\cite{zawalski2025ecot} demonstrate that intermediate reasoning steps---whether expressed in natural language or visual tokens---significantly enhance long-horizon robustness. Recognizing that verbose CoT traces incur prohibitive inference latency, \textbf{Fast-ThinkAct}~\cite{huang2026fastthinkact} distills explicit reasoning into compact \emph{verbalizable latent plans} via a preference-guided objective, retaining long-horizon planning and failure-recovery capabilities while reducing inference latency by up to 89\%. Beyond compressing language-based reasoning, the field is diversifying the \emph{medium} of intermediate deliberation: \textbf{CoA-VLA}~\cite{li2025coavla} replaces sub-task language with visual-textual \emph{chains of affordance}, grounding each reasoning step in actionable contact regions, while \textbf{ACoT-VLA}~\cite{zhong2026acotvla} argues that the most effective reasoning deliberates directly in the action space, proposing coarse reference trajectories (Explicit Action Reasoner) combined with latent action priors (Implicit Action Reasoner) to condition the downstream policy. This is further augmented by graph-based reasoning (GraphCoT-VLA~\cite{jin2024graphcot}), and specialized reasoning agents (e.g., Cosmos-Reason1~\cite{cosmosreason2024}), which externalize intermediate decisions to make failure recovery and instruction revision more tractable.

In parallel, the structural design of dual-system architectures has matured to formalize the interface between cognition and execution. FiS-VLA~\cite{chen2025fis} operationalizes dual-speed control by coupling a slower semantic branch with a real-time execution branch, directly alleviating control-side latency pressure. This modular pattern is reinforced by systems like A0~\cite{xu2025a0}, RoBridge~\cite{zhang2025robridge}, and Maestro~\cite{shi2025maestro}, which establish explicit intermediate intent representations. 

Recent extensions broaden this modular trend along the axes of adaptive reasoning, safety, and cross-embodiment deployment. On the front of dynamic execution and reasoning, systems like Hi Robot~\cite{shi2025hirobot} employs a System 1/System 2 hierarchical architecture to process complex open-ended instructions and real-time human feedback. Safety and stability are explicitly addressed by integrating Lyapunov-stable neural control~\cite{yang2024lyapunov} and safe reinforcement learning frameworks. Notably, SafeVLA~\cite{zhang2025safevla} exemplifies this by formulating policy alignment as a Constrained Markov Decision Process (CMDP), explicitly trading off task performance to ensure that policy outputs strictly adhere to real-world safety bounds. On the deployment front, frameworks like HAMSTER~\cite{li2025hamster} emphasize cross-morphology transfer under mixed semantic-control pipelines, and methods like ControlVLA~\cite{li2025controlvla} enable highly efficient object-centric adaptation without catastrophic forgetting. Ultimately, industrial-scale efforts like GO-1~\cite{agibot2025} and GR00T N1~\cite{nvidia2025grootn1} scale this hierarchical design to bimanual and humanoid settings. Collectively, these works establish hierarchy as a fundamental design principle for preserving semantic competence under real-time physical constraints. Complementary to reasoning chains, \textbf{MemoryVLA}~\cite{shi2026memoryvla} introduces a dual-memory architecture fusing short-term working memory with long-term episodic-semantic memory, enabling temporally coherent decision-making across extended manipulation horizons.

\subsubsection{Type III\@: Spatial Grounding \& Predictive World Models}

While Type II architectures successfully resolve the temporal mismatch between reasoning and control, they fundamentally inherit the spatial ambiguities of their 2D vision-language backbones. When tasks demand precise geometric consistency under occlusion, articulated contacts, or viewpoint shifts, time-scale decoupling alone is insufficient. This motivates the Type III transition, where the core objective shifts from \emph{how to schedule reasoning} to \emph{how to maintain and predict an executable world state}. The progression within this category follows a natural hierarchy: explicit 3D spatial representations first establish an accurate model of the current world state, which then serves as the foundation for predictive world models that simulate future states to guide planning and control.

The foundational shift in this paradigm involves lifting representations from 2D pixels to 3D point clouds and voxels. Pioneering frameworks like 3D-VLA~\cite{zhen20243dvla} and SpatialVLA~\cite{qu2025spatialvla} natively integrate 3D perception into the LLM embedding space, enabling spatial reasoning and 3D goal generation. This architecture is further optimized by PointVLA~\cite{li2026pointvla} and CubicVLA~\cite{ji2026cubicvla}, which streamline 3D tokenization to achieve efficient, high-frequency manipulation without sacrificing spatial granularity.

Beyond raw 3D inputs, injecting geometric priors significantly enhances sample efficiency and generalization. 3D Diffuser Actor~\cite{3ddiffuseractor2024} formulates manipulation as a 3D conditional generation problem, leveraging spatial attention. To handle SE(3) transformations robustly, EquiGraspFlow~\cite{lim2024equigraspflow} enforce strict equivariance and symmetry constraints, ensuring that learned policies generalize zero-shot to novel object poses and spatial configurations.

To capture fine-grained physical interactions, recent works leverage Neural Fields and 3D Gaussian Splatting (3DGS). 3DAffordSplat~\cite{wei20253daffordsplat} embeds actionable affordances directly into 3DGS representations, enabling real-time geometric queries. For dynamic and articulated objects, PhysTwin~\cite{jiang2025phystwin} explicitly models kinematic constraints and elastodynamics. Furthermore, maintaining spatial consistency across video frames is critical for closed-loop execution; thus, models like TraceVLA~\cite{zheng2025tracevla} and TrackVLA~\cite{wang2025trackvla} integrate temporal tracking mechanisms to handle dynamic environments and visual occlusions. Beyond temporal continuity, precise 3D spatial reasoning is essential for executing complex physical tasks. To this end, RoboRefer~\cite{zhou2025roborefer} advances spatial referring by precisely localizing specific 3D targets based on complex language instructions. Building upon this, RoboTracer~\cite{zhou2025robotracer} extends spatial reasoning to continuous, metric-grounded spatial tracing, generating ordered 3D keypoint sequences for multi-step manipulation trajectories.MolmoAct~\cite{lee2025molmoact} elevates this paradigm by explicitly factorizing the decision-making process into depth-aware perception tokens and human-editable visual reasoning traces. By visualizing its intended 2.5D spatial trajectory before execution, MolmoAct not only grounds semantic intent in structural reality but also ensures highly steerable and interpretable low-level control.

Spatial awareness naturally extends into the temporal domain via predictive World Models. By simulating future states, models like RoboDreamer~\cite{zhou2024robodreamer}, and GWM~\cite{lu2025gwm} enable agents to ``imagine'' the consequences of their actions before execution. This capability is elevated to 4D by Learning 4D Embodied World Models in TesserAct~\cite{zhen2025tesseract} and SplatSim~\cite{qureshi2024splatsim}, which predict dynamic 3DGS environments. Concurrently, the ``Video-as-Action'' paradigm leverages video generation as a proxy for planning. Frameworks such as VideoVLA~\cite{shen2025videovla}, Gen2Act~\cite{bharadhwaj2024gen2act}, Dreamitate~\cite{liang2024dreamitate}, Dream2Real~\cite{kapelyukh2024dream2real}, and Unified Video Action Model~\cite{li2025uva} synthesize future video frames to guide low-level control, effectively bypassing explicit state estimation by predicting the visual future. Pushing this unification further, \textbf{VITA}~\cite{gao2025vita} establishes a cross-modal shared codebook whose latent tokens are simultaneously decoded into future frame predictions and robot actions, internalizing visual dynamics as an implicit inductive bias. Eliminating the need for action labels during pretraining, \textbf{ViPRA}~\cite{routray2025vipra} learns motion-centric latent action priors from unlabeled videos and refines them via flow matching into continuous control policies. Scaling this unsupervised paradigm to foundation-model level, \textbf{LAPA}~\cite{ye2025lapa} trains a VQ-VAE to discretize latent actions from internet-scale video, then pretrains a 7B model via next-token prediction on these pseudo-action tokens---achieving competitive manipulation performance with orders-of-magnitude less robot data.

Crucially, recent advancements have pushed world models beyond simple video generation towards robust, action-conditioned latent representations. \textbf{GR-2}~\cite{cheang2024gr2} exemplifies this by employing a GPT-style CVAE to learn a highly scalable, action-conditioned video prediction model, demonstrating that scaling generative world models directly translates to improved multi-task manipulation performance. Whereas GR-2 conditions video generation on actions but separates the two objectives, Unified World Models~\cite{zhu2025uwm} demonstrates that a single Diffusion Transformer can jointly model both modalities by assigning \emph{independent} diffusion timesteps to action and observation channels. This decoupling enables the same network to function as a policy ($t_{o'}\!=\!T$), forward dynamics model ($t_a\!=\!0$), inverse dynamics model ($t_{o'}\!=\!0$), or video predictor ($t_a\!=\!T$) at inference time, depending solely on which timestep is clamped. By treating missing action labels as maximally noised inputs, UWM naturally absorbs action-free video during pretraining, yielding up to 20\% success-rate gains over Diffusion Policy on real Franka tasks and establishing a principled bridge between world modeling and policy learning within a single architecture. Similarly, \textbf{V-JEPA 2}~\cite{assran2025vjepa} shifts the paradigm from pixel reconstruction to latent predictive reasoning, utilizing a Joint-Embedding Predictive Architecture to infer physics-compliant features without the noise of texture generation. To bridge the gap between these predictive models and physical execution, \textbf{DreamVLA}~\cite{zhang2025dreamvla} integrates a diffusion-based world model directly into the Type~I pipeline, allowing the agent to plan via imagined trajectories. Furthermore, \textbf{BridgeVLA}~\cite{li2025bridgevla} and \textbf{RoBridge}~\cite{zhang2025robridge} focus on aligning these spatial and predictive representations specifically for 3D manipulation, ensuring that the imagined states are kinematically feasible and geometrically grounded for the downstream policy. Along the axis of latent motion representations, \textbf{Moto}~\cite{chen2025moto} learns a bridging ``language'' of motion by unsupervised tokenization of video into latent motion token sequences, pre-training via next motion token prediction and co-fine-tuning for downstream robot control. Extending this direction, \textbf{CoWVLA}~\cite{yang2026chainofworld} introduces world model thinking in latent motion space, enabling temporally continuous dynamic modeling before action generation.

Finally, true spatial intelligence requires multi-modal physical grounding beyond vision. Contact-rich manipulation necessitates explicit force and tactile feedback. Tactile-VLA~\cite{huang2025tactilevla} integrate force-torque and tactile sensor streams directly into the foundation model backbone. 3D-ViTac~\cite{huang20253dvitac}, AnySkin~\cite{bhirangi2024anyskin} advance this by mapping high-resolution tactile signals to 3D spatial representations, enabling dexterous manipulation of deformable or visually occluded objects. Furthermore, audio provides a critical proxy for physical contact and material properties; ManiWAV~\cite{liu2024maniwav} demonstrates that integrating audio-visual signals significantly enhances the robustness of dynamic manipulation tasks, completing the multi-modal spatial awareness required for generalist embodied agents.

\subsection{Deployment Lifecycle}

While the preceding sections classify embodied models by \emph{architectural design} (Type~I--III), a complementary engineering question remains: how does a multi-billion-parameter model---whether a unified autoregressive model, a diffusion policy, a hierarchical reasoning system, or a predictive world model---survive the transition from cloud-scale training to continuous 10\,Hz control on resource-constrained hardware? This section organizes recent work into a 3-Stage Deployment Lifecycle based on when the core computation occurs: Stage~1 (offline---in simulation, on static datasets, or in the cloud), Stage~2 (online---requiring direct physical interaction), and Stage~3 (runtime---during continuous inference on the edge). Table~\ref{tab:deployment_lifecycle} summarizes the bottlenecks and representative techniques at each stage.

\begin{table*}[t!]
\centering
\caption{Deployment bottlenecks and representative techniques across the 3-Stage pipeline.
\textbf{S1}\,=\,Offline Pre-deployment; \textbf{S2}\,=\,Online Physical Alignment; \textbf{S3}\,=\,Inference-Time Execution.}\label{tab:deployment_lifecycle}
\renewcommand{\arraystretch}{1.15}
\setlength{\tabcolsep}{3pt}
\scriptsize
\begin{tabular}{
  c
  p{0.20\linewidth}
  p{0.33\linewidth}
  p{0.34\linewidth}
}
\toprule
\textbf{Stage} & \textbf{Deployment Bottleneck} & \textbf{Technique Category} & \textbf{Representative Methods} \\
\midrule
\multirow{4}{*}{\textbf{S1}}
 & Sim-to-Real Gap & Physics alignment, domain adaptation & Transic~\cite{jiang2024transic}, ASAP~\cite{he2025asap} \\
 & Cross-Embodiment & Latent action spaces, R2S2R pipelines & X-Sim~\cite{dan2025xsim}, Lat.\ Act.\ Diff.~\cite{bauer2025latentaction}, CycleVAE~\cite{dastider2025cyclevae} \\
 & Model Compression & Knowledge distillation, action tokenization & Shallow-$\pi$~\cite{jeon2026shallowpi}, FAST~\cite{pertsch2025fast} \\
 & Offline RL at Scale & Autoregressive Q-fn, dense return pred. & Q-Transformer~\cite{chebotar2023qtransformer}, ReinboT~\cite{zhang2025reinbot} \\
\midrule
\multirow{4}{*}{\textbf{S2}}
 & RL Fine-Tuning & Online PPO/GRPO, RL data generation & RLDG~\cite{xu2024rldg}, SimpleVLA-RL~\cite{li2025simplevla}, FLaRe~\cite{hu2024flare} \\
 & Contact-Rich Adapt. & Reinforced FT, HITL, visuo-tactile loop & ConRFT~\cite{chen2025conrft}, HITL~\cite{luo2025humaninloop}, VT-Refine~\cite{huang2025vtrefine} \\
 & Forgetting \& Stability & Consistency reg., latent-space RL & DASP~\cite{liu2025dasp}, DSRL~\cite{wagenmaker2025steering} \\
 & Residual Learning & Generalist--specialist decomposition & ManipTrans~\cite{li2025maniptrans}, PLD~\cite{xiao2025pld}, OFT~\cite{kim2025oft} \\
\midrule
\multirow{4}{*}{\textbf{S3}}
 & Quantization \& Pruning & Mixed-precision, 1-bit PTQ, token pruning & QVLA~\cite{xu2026qvla}, HBVLA~\cite{yan2026hbvla}, LightVLA~\cite{jiang2025lightvla} \\
 & Decoding Acceleration & Parallel decoding, distillation, async pipe. & PD-VLA~\cite{pdvla2025}, TinyVLA~\cite{wen2025tinyvla}, SwiftVLA~\cite{ni2025swiftvla} \\
 & Safety \& Robustness & Contrastive decoding, certifiable planning & Policy CD~\cite{wu2026policycontrastive}, Neural MP~\cite{dalal2024neuralmp} \\
 & Failure Recovery & Uncertainty detection, language-guided repair & FAIL-Detect~\cite{xu2025detectfailures}, RACER~\cite{dai2024racer}, AHA~\cite{duan2024aha} \\
\bottomrule
\end{tabular}
\end{table*}

\subsubsection{Stage 1: Offline Pre-deployment (The Preparation Phase)}

Pure simulation training frequently degrades upon physical deployment because contact-rich dynamics---surface friction, cable compliance, sensor noise---remain difficult to model faithfully. Transic~\cite{jiang2024transic} uses offline corrections augmented by human-in-the-loop interventions to transfer short-horizon manipulation skills, while ASAP~\cite{he2025asap} focuses on automated physics alignment via high-fidelity simulation and system identification for whole-body humanoid locomotion. Both reduce the sim-to-real gap through complementary mechanisms---task-specific human corrections and automated system identification, respectively---yet residual discrepancies typically persist and propagate as adaptation challenges for Stage~2. Beyond dynamics alignment, the absence of paired multi-robot datasets further restricts foundation pre-training; recent solutions converge on \emph{shared latent action spaces} as the interoperability layer. X-Sim~\cite{dan2025xsim} deploys a Real-to-Sim-to-Real pipeline routing through simulation, Latent Action Diffusion~\cite{bauer2025latentaction} constructs embodiment-agnostic representations in a shared latent action space, and CycleVAE~\cite{dastider2025cyclevae} aligns latent motion sequences from guided demonstrations---all decoupling what the robot does from how it is actuated, but differing in whether alignment occurs in observation, action, or motion space. At the representation pretraining level, \textbf{DeFI}~\cite{zhang2026disentangled} decouples forward and inverse dynamics pretraining, exploiting large-scale unlabeled video for physical prediction and labeled robot data for action alignment, enabling visual encoders to acquire both capabilities simultaneously.

Multi-billion-parameter models generally exceed the compute budget of onboard hardware. Two complementary strategies address this: \emph{structural distillation} (Shallow-$\pi$~\cite{jeon2026shallowpi} distills backbone and action head by up to 70\% in parameter count) and \emph{representation compression} (FAST~\cite{pertsch2025fast} replaces per-dimension binning with Frequency-space Action Sequence Tokenization via the Discrete Cosine Transform). These strategies can potentially compound in practice, but their joint design and empirical interaction remain largely unexplored. Meanwhile, training RL policies on static datasets avoids costly online interaction but introduces optimization challenges at billion-parameter scale. Q-Transformer~\cite{chebotar2023qtransformer} demonstrates that autoregressive Q-functions can scale offline actor-critic learning to massive transformers; ReinboT~\cite{zhang2025reinbot} relaxes the need for optimal demonstrations by predicting dense returns; and further architectural modifications~\cite{springenberg2024offlineactorcritic} stabilize offline actor-critic methods at extreme scale. Together, these establish that value-based pre-training is viable across model sizes, though its interaction with downstream online fine-tuning remains underexplored.

The decisions made in Stage~1 directly constrain Stage~2: aggressive compression reduces the model's capacity to absorb online fine-tuning gradients, and the choice of sim-to-real strategy determines whether residual gaps manifest as systematic dynamics errors or coverage gaps---each propagating as a distinct adaptation challenge.

\subsubsection{Stage 2: Online Physical Alignment (The Adaptation Phase)}

Once deployed on the physical robot, the model encounters real-world friction, out-of-distribution objects, and contact forces absent from any simulation---conditions that offline pipelines rarely anticipate in full.

Reinforcement learning offers an autonomous alternative to costly supervised fine-tuning, and recent work shows rapid iteration in both algorithms and engineering recipes. RLDG~\cite{xu2024rldg} isolates RL instability by first generating high-quality trajectory data via sample-efficient real-world RL, then distilling it into the generalist policy. SimpleVLA-RL~\cite{li2025simplevla} eliminates this two-phase separation, applying group relative policy optimization (GRPO) directly to 7B+ models, while an empirical study~\cite{liu2026rlvla} reports that online PPO can improve both semantic understanding and execution robustness under their evaluation setting. OFT~\cite{kim2025oft} identifies practical recipes (parallel decoding, continuous actions, L1 regression loss) that maximize success rate per fine-tuning hour, and FLaRe~\cite{hu2024flare} addresses the tension between generalist priors and specialist performance through adaptive large-scale RL fine-tuning. To standardize these diverse fine-tuning pipelines, Policy-Agnostic RL~\cite{mark2024policyagnosticrl} introduces a unified framework capable of conducting both offline and online RL fine-tuning across arbitrary policy classes and backbones, effectively decoupling the RL algorithm from the underlying model architecture. DayDreamer~\cite{wu2023daydreamer} improves sample efficiency through learned world models: a latent dynamics model trained on real robot data generates imagined rollouts for policy optimization, enabling a quadruped to learn locomotion within one hour and a robot arm to acquire pick-and-place behaviors within several hours of physical interaction, WMPO~\cite{zhu2025wmpo} directly extends world model-based policy optimization to multi-modal foundation model architectures, enabling rapid online alignment without heavily relying on costly physical rollouts. 
Parallel to gradient-based and model-based methods, \textbf{EvoRL}~\cite{zheng2025evorl} introduces a GPU-accelerated evolutionary reinforcement learning framework that overcomes the computational bottleneck of population-based search, enabling massive parallel exploration for complex contact-rich tasks on a single consumer-grade GPU. 
Overall, RL fine-tuning is often competitive with supervised alternatives, but the optimal recipe (direct vs.\ staged, PPO vs.\ GRPO) remains task- and constraint-dependent. For tasks involving tight tolerances or multi-finger contact, a harsher adaptation regime emerges: ConRFT~\cite{chen2025conrft} integrates offline BC/Q-learning with an online consistency policy and human safety checkpoints; for sub-millimeter dexterous tasks where autonomous RL alone often proves insufficient,~\cite{luo2025humaninloop} deploys interactive RL with periodic human interventions---underscoring the current limits of purely autonomous adaptation; and VT-Refine~\cite{huang2025vtrefine} demonstrates that continuous visuo-tactile feedback is essential for physically aligning dual robotic arms at microscopic coordination scales.

A critical concern for all Stage~2 methods is catastrophic forgetting. Residual approaches address this by freezing the base policy and training lightweight correction modules: ManipTrans~\cite{li2025maniptrans} fine-tunes a specialist residual module under strict physics constraints for bimanual tasks, while PLD~\cite{xiao2025pld} trains lightweight residual actors via off-policy RL to autonomously probe failure states and distill recoveries back. To stabilize full-model fine-tuning directly, DASP~\cite{liu2025dasp} uses a diffusion autodecoder encoding skill primitives as stable attractors, and DSRL~\cite{wagenmaker2025steering} applies RL strictly in latent space for diffusion policies. These mechanisms reveal a fundamental trade-off: stronger regularization preserves pre-trained knowledge but limits adaptation to novel physics, while weaker regularization enables rapid adaptation at the risk of skill collapse.

The adaptation strategies chosen here directly shape Stage~3: full-parameter RL fine-tuning yields higher task performance but harder-to-quantize weight distributions, while residual architectures maintain the original weight structure---simplifying Stage~3 compression---at the cost of additional inference-time computation. Distribution shifts from online RL also propagate as a safety concern, as decision boundaries that moved during adaptation may produce novel failure modes absent from both simulation and fine-tuning.

\subsubsection{Stage 3: Inference-Time Execution (The Runtime Phase)}

During continuous operation, the model need sustain 10\,Hz+ control loops on resource-constrained edge hardware while providing safety guarantees---the bottleneck has shifted from \emph{learning} to \emph{serving}.

Applying uniform-bit quantization to large robotic models typically causes significant degradation in manipulation benchmarks because action-critical channels require higher precision than language-generation channels. QVLA~\cite{xu2026qvla} addresses this via action-centric, channel-wise bit allocation guided by action-space sensitivity, while HBVLA~\cite{yan2026hbvla} pushes to 1-bit post-training quantization via policy-grounded saliency and Haar wavelet transforms. Orthogonally, token pruning attacks the input dimension: LightVLA~\cite{jiang2025lightvla} adaptively drops irrelevant visual background tokens, and EfficientVLA~\cite{yang2025efficientvla} introduces a training-free algorithm optimizing the KV cache on the fly. Beyond weight and token compression, action chunking in autoregressive models introduces temporal latency that compounds across the control loop. PD-VLA~\cite{pdvla2025} reformulates autoregressive decoding into parallel fixed-point iterations generating entire action chunks simultaneously; TinyVLA~\cite{wen2025tinyvla} distills massive models into fast, compact variants; SwiftVLA~\cite{ni2025swiftvla} intelligently routes spatiotemporal features to minimize compute overhead; SARA-RT~\cite{shi2024sarart} reduces the quadratic attention complexity of large robotics transformers to linear cost by up-training softmax attention into self-adaptive robust attention, achieving approximately 14\% wall-clock speedup without degrading manipulation accuracy; and at the system level, Xiaomi-Robotics-0~\cite{cai2026xiaomirobotics0} achieves 80\,ms end-to-end latency via asynchronous execution with adaptive loss re-weighting. SP-VLA~\cite{li2025spvla} dynamically switches between the heavy foundation model and a lightweight action generator based on task complexity. These methods reveal a design spectrum from \emph{algorithmic} acceleration (PD-VLA, SwiftVLA) to \emph{system-level} pipelining (Xiaomi-Robotics-0), with the optimal combination likely being hardware-dependent.

Large vision-conditioned policies often over-rely on spurious visual correlations, producing catastrophic failures in uncontrolled environments. Policy Contrastive Decoding~\cite{wu2026policycontrastive} addresses this by contrasting action distributions from original versus object-masked inputs. However, even a highly robust neural policy does not by itself provide \emph{formal} collision-avoidance guarantees; Neural MP~\cite{dalal2024neuralmp} bridges this gap by combining neural intent generation with certifiable motion planning. Detecting and recovering from execution failures is equally critical. FAIL-Detect~\cite{xu2025detectfailures} uses flow-based density estimators to flag out-of-distribution states as candidate failures using only positive demonstrations; RACER~\cite{dai2024racer} deploys language-guided failure recovery policies; and AHA~\cite{duan2024aha} extends this to open-ended diagnosis by fine-tuning an open-source VLM to produce free-form language rationales for physical failures---charting a progression from binary detection to structured recovery to open-ended self-diagnosis. Reflective Test-Time Planning~\cite{hong2026rttp} pushes this progression further by converting diagnosed failures into self-supervised learning signals: dual reflection loops drive LoRA policy updates at runtime, reporting a $\sim$3$\times$ improvement on long-horizon household tasks in BEHAVIOR-1K (33.7\% vs.\ 11.2\%) under their evaluation setting and demonstrating that online self-improvement is a viable Stage~3 capability when the overhead fits the control-loop budget.

Table~\ref{tab:data_learning_map} synthesizes the mapping between the data ecosystem (\S\ref{sec:eai_datasets}) and the learning paradigms of this section, making explicit how each data regime constrains the choice of cognitive planning, policy learning, and integration architecture.

\begin{table*}[t]
\centering
\caption{Mapping from data regimes to learning paradigms. Each row traces how a data tier (\S\ref{sec:eai_datasets}) flows through cognitive planning (Brain, \S\ref{sec:embodied_brain}), policy learning (Cerebellum, \S\ref{sec:policy_learning_low_level}), and integration architecture (\S4.3), together with the resulting system-level implication. Paradigm codes follow Tables~\ref{tab:brain_paradigms}--\ref{tab:policy_paradigms}. Note that $\mathcal{K}_3 \supset \mathcal{K}_2 \supset \mathcal{K}_1$ by the recursive knowledge-state definition (Eq.~\ref{eq:knowledge_state}); each row highlights the \emph{marginal} information contribution of the corresponding tier.}\label{tab:data_learning_map}
\renewcommand{\arraystretch}{1.3}
\setlength{\tabcolsep}{3pt}
\scriptsize
\begin{tabular}{
  c
  >{\centering\arraybackslash}p{0.06\linewidth}
  c
  p{0.25\linewidth}
  >{\centering\arraybackslash}p{0.10\linewidth}
  p{0.25\linewidth}
}
\toprule
\textbf{Data Regime} &
\textbf{Brain} &
\textbf{Cerebellum} &
\textbf{Planning--Policy Coordination} &
\textbf{Integration} &
\textbf{System Implication} \\
\midrule
Skill-level ($\mathcal{K}_1$) &
B4 &
P1; P3 &
Loose---geometric targets $\to$ local force execution &
Type~III (spatial) &
Precise single-step primitives; no multi-step composition \\
\midrule
Task-level ($\mathcal{K}_2$) &
B1/B2; B3 &
P1; P2 &
Tight---slow decomposition $\to$ fast tracking &
Type~II (hierarchical) &
Long-horizon tasks; bounded by scene diversity and format alignment \\
\midrule
\multirow{2}{*}[-1.2em]{Foundation-level ($\mathcal{K}_3$)} &
B3; B5 &
P2; P3 &
Implicit---shared backbone absorbs both &
Type~I (unified) &
Broad transfer; 3D precision remains a bottleneck \\
\cmidrule{2-6}
 &
B4/B5 &
P3 &
Predictive---world model imagines, policy executes &
Type~III (world models) &
Pre-execution verification; fidelity degrades over long horizons \\
\bottomrule
\end{tabular}
\end{table*}

The learning paradigms and deployment strategies outlined above establish \emph{how} manipulation intelligence is constructed and operationalized. Yet a critical question remains: how do we systematically measure whether these systems work? The next section addresses this gap by synthesizing the emerging evaluation landscape.

\section{Evaluation metrics of Embodied Manipulation}\label{sec:evaluation}

Sections~\ref{sec:embodiment_hardware}--\ref{sec:embodiment_method} have traced the organizational chain from hardware through data to learned policies and their deployment. Yet without rigorous evaluation, this chain remains open-loop: we cannot quantify progress, diagnose failure modes, or compare systems across platforms. Traditional evaluation relies predominantly on binary success rates, which obscure \emph{why} a policy fails and hinder reproducible comparison across diverse embodiments. As task-level datasets scale in diversity (\S\ref{sec:task_datasets_narrative}) and foundation models grow to billions of parameters (\S4.3), the evaluation bottleneck has become acute~\cite{zhou2025autoeval, atreya2025roboarena}. This section synthesizes the emerging landscape of evaluation methodologies, organizing them into four complementary dimensions: \textbf{task execution quality} (\S\ref{sec:eval_task_execution}), which diagnoses completion rates, trajectory smoothness, and grasp stability; \textbf{data and sample evaluation} (\S\ref{sec:eval_data_sample}), which quantifies dataset quality and learning efficiency; \textbf{human preference and generalization} (\S\ref{sec:eval_human_generalization}), which captures subjective quality and out-of-distribution robustness; and \textbf{energy and efficiency} (\S\ref{sec:eval_energy}), which addresses the physical constraints of continuous deployment.

\subsection{Task Execution Evaluation}\label{sec:eval_task_execution}

\begin{table*}[t]
\centering
\caption{Evaluation metrics and representative works for embodied manipulation, organized by four complementary dimensions.}\label{tab:eval_metrics}
\renewcommand{\arraystretch}{1.15}
\setlength{\tabcolsep}{3pt}
\scriptsize
\begin{tabular}{
  c
  p{0.21\linewidth}
  p{0.33\linewidth}
  p{0.29\linewidth}
}
\toprule
\textbf{Dimension} & \textbf{Metric Category} & \textbf{Representative Metrics} & \textbf{Representative Works} \\
\midrule
\multirow{4}{*}{\shortstack{\textbf{Task}\\\textbf{Execution}\\(\S\ref{sec:eval_task_execution})}} 
 & Task Completion & Stage-wise Success, Task Progression, Normalized Efficiency & RoboEval~\cite{wang2025roboeval}, REALM~\cite{sedlacek2025realm}, Eval-Actions~\cite{liu2026evalactions} \\
 & Trajectory Smoothness & SPARC, LDLJ, Path Length, Jerk & SPARC~\cite{balasubramanian2015sparc}, LDLJ~\cite{antonelli2023movement}, RoboEval~\cite{wang2025roboeval} \\
 & Trajectory Similarity \& Grasp & DTW, Fr\'{e}chet Distance, $\epsilon$-metric, $v$-metric & Fr\'{e}chet~\cite{agarwal2012computingdiscretefrechetdistance}, Ferrari \& Canny~\cite{ferrari1992grasp} \\
 & Dual-Arm Coordination & Velocity Divergence, Height Discrepancy & RoboEval~\cite{wang2025roboeval} \\
\midrule
\multirow{3}{*}{\shortstack{\textbf{Data \&}\\\textbf{Sample}\\(\S\ref{sec:eval_data_sample})}} 
 & Info-theoretic & Mutual Information, Demo Contribution & DemInf~\cite{hejna2025robot} \\
 & Gradient-based & Influence Scores, Hessian-based Valuation & CUPID~\cite{agia2025cupidcuratingdatarobot}, QoQ~\cite{lee2026qoq} \\
 & Sample Efficiency & AULC & General IL benchmarks \\
\midrule
\multirow{3}{*}{\shortstack{\textbf{Human Pref.}\\\textbf{\& Gen.}\\(\S\ref{sec:eval_human_generalization})}} 
 & Human Preference & Bradley-Terry, VLM-based Eval & RoboArena~\cite{atreya2025roboarena}, AutoEval~\cite{zhou2025autoeval} \\
 & OOD Generalization & $\Delta SR_{\text{OOD}}$, RMSD & REALM~\cite{sedlacek2025realm} \\
 & Sim-to-Real & MMRV, Pearson $\rho$, VLM-as-Judge, $\Delta SR$ & SIMPLER~\cite{li2025simpler}, WorldArena~\cite{shang2026worldarena} \\
\midrule
\multirow{4}{*}{\shortstack{\textbf{Energy \&}\\\textbf{Efficiency}\\(\S\ref{sec:eval_energy})}} 
 & Power Modeling & Joint Power $P_i(t)$, Torque-squared $E_{\text{proxy}}$ & Classical electrodynamics \\
 & Energy Decomposition & $P_{\text{elec}}$, $P_{\text{friction}}$, $P_{\text{copper}}$, $P_{\text{mech}}$ & ECDP~\cite{heredia2023ecdp} \\
 & Normalized Efficiency & $\eta_E(\mathbf{q})$, $f_U/f_C/f_R$, CoT & Vidussi et al.~\cite{vidussi2021energy}, \cite{heredia2025evaluating}, Tucker~\cite{tucker1975energetic} \\
 & Multi-objective Optimization & Pareto front via LSTM surrogate & Rezali et al.~\cite{rezali2025lstm} \\
\bottomrule
\end{tabular}
\end{table*}

Task execution evaluation serves as the primary validation layer for embodied manipulation policies, progressing from coarse binary success rates to fine-grained diagnostic metrics that isolate failure modes across perception, planning, and control. This subsection examines the evolution from stage-wise completion metrics to trajectory quality analysis, grasp stability quantification, and dual-arm coordination assessment.

\paragraph{Task Completion and Stage-wise Progress.}
Binary success rates, while interpretable, fail to distinguish between policies that fail immediately versus those that complete 90\% of a task before encountering a final obstacle. \textbf{RoboEval}~\cite{wang2025roboeval} addresses this limitation by introducing \textbf{stage-wise success metrics}, decomposing long-horizon tasks into $K$ sequential stages and computing progress as $\text{Progress} = \frac{1}{K}\sum_{k=1}^K S_k$, where $S_k = \mathbb{1}[\text{stage } k \text{ complete}]$. This granular decomposition enables precise failure diagnosis: a policy achieving 0.8 progress has successfully navigated 80\% of the task structure, providing actionable feedback for targeted improvement. Complementing this discrete formulation, \textbf{REALM}~\cite{sedlacek2025realm} introduces \textbf{Task Progression} (TP), a continuous scoring mechanism that assigns importance weights to each stage: $TP = \sum_{k=1}^K w_k \cdot s_k$, where $s_k \in [0,1]$ represents partial completion and $\sum_{k=1}^K w_k = 1$. To quantify robustness under environmental perturbations, REALM further defines the \textbf{Root Mean Square Deviation} (RMSD) across perturbation types $p$:
\begin{equation}
  \text{RMSD}(p) = \sqrt{\frac{1}{|M||T_{\text{tasks}}|} \sum_{m,\tau} \left(TP_{m,\tau,p} - TP_{m,\tau,\text{default}}\right)^2}
\end{equation}
where $m$ indexes models and $\tau$ indexes tasks. An ideal generalist model exhibits RMSD $\to 0$, indicating invariance to visual, semantic, and physical perturbations. Building on automated evaluation infrastructure, \textbf{Eval-Actions}~\cite{liu2026evalactions} leverages vision-language models to autonomously assess task completion, computing \textbf{normalized efficiency} as $E = T_{\text{expert}}/T_{\text{agent}} \in (0,1)$, where values approaching 1 indicate expert-level temporal efficiency.

\paragraph{Trajectory Smoothness and Execution Quality.}
Beyond task completion, execution quality reflects the policy's internal coherence and physical plausibility. Jerky, oscillatory motions indicate poor generalization or unstable control, even when tasks ultimately succeed. \textbf{SPARC}~\cite{balasubramanian2015sparc} quantifies smoothness in the frequency domain by measuring the arc length of the normalized velocity spectrum:
\begin{equation}
  \text{SPARC} = -\int_0^{\hat\omega} \sqrt{\left(\frac{1}{\hat\omega}\right)^2 + \left(\frac{d\hat{V}(\omega)}{d\omega}\right)^2}\, d\omega
\end{equation}
where $\hat{V}(\omega) = V(\omega)/V(0)$ is the normalized velocity magnitude spectrum and $\hat\omega$ is the cutoff frequency (typically 10\,Hz). SPARC values are negative, with values closer to zero indicating smoother motion. Complementing this frequency-domain perspective, \textbf{LDLJ} (Log Dimensionless Jerk)~\cite{antonelli2023movement} provides a time-domain alternative via dimensionless normalization:
\begin{equation}
  \text{LDLJ} = -\log\left(\frac{D^5}{A^2 T^3} \int_0^T \left(\frac{d^3 v}{dt^3}\right)^2 dt\right)
\end{equation}
where $D$ is movement amplitude, $A$ is peak velocity, and $T$ is movement duration. Higher LDLJ values (closer to zero) indicate smoother trajectories. RoboEval extends these classical metrics with robot-specific diagnostics: \textbf{path length} in joint space ($L_{\text{joint}} = \sum_{t=1}^{T-1} \|\mathbf{q}_{t+1} - \mathbf{q}_t\|_2$) and Cartesian space ($L_{\text{cart}} = \sum_{t=1}^{T-1} \|\mathbf{p}_{t+1} - \mathbf{p}_t\|_2$), \textbf{jerk} computed via third-order finite differences ($J = \frac{1}{T-3}\sum_{t=1}^{T-3} \|\Delta^3 \mathbf{q}_t / \Delta t^3\|_2$), and discrete event counters for \textbf{collisions} and \textbf{slip events}. To assess trajectory fidelity relative to expert demonstrations, \textbf{Dynamic Time Warping} computes the minimum-cost alignment between two sequences via dynamic programming: $\text{DTW}(P, Q) = \min_{\text{path}} \sum_{(i,j)} d(p_i, q_j)$, allowing temporal warping to accommodate execution speed variations. In contrast, the \textbf{Fr\'{e}chet Distance}~\cite{agarwal2012computingdiscretefrechetdistance} measures the worst-case deviation under monotonic reparameterization: $d_F(P, Q) = \inf_{\alpha,\beta} \max_{t\in[0,1]} d(P(\alpha(t)), Q(\beta(t)))$, making it particularly sensitive to critical waypoint deviations in contact-rich tasks. Beyond trajectory-level metrics, manipulation success ultimately hinges on stable grasps. The \textbf{$\epsilon$-metric}~\cite{ferrari1992grasp} quantifies force closure robustness as the minimum singular value of the grasp matrix: $\epsilon = \lambda_{\min}(\mathbf{G}\mathbf{G}^\top)$, where $\mathbf{G}$ maps contact forces to object wrenches; larger $\epsilon$ values indicate greater resistance to perturbations. The complementary \textbf{$v$-metric} measures global grasp versatility as the volume of the feasible wrench space: $v = \text{Vol}(\mathcal{W}_c)$, where $\mathcal{W}_c$ is the convex hull of achievable 6D force-torque combinations. These geometric metrics integrate naturally with modern dexterous manipulation benchmarks such as DexGraspNet~\cite{wang2023dexgraspnetlargescaleroboticdexterous}.

\paragraph{Dual-Arm Coordination Metrics.}
Bimanual manipulation introduces coordination challenges absent in single-arm systems. RoboEval defines \textbf{velocity divergence} as $V_{\text{div}} = \frac{1}{T}\sum_{t=1}^T \|\dot{\mathbf{p}}_t^L - \dot{\mathbf{p}}_t^R\|_2$, quantifying asynchrony between left and right end-effectors, and \textbf{height discrepancy} as $H_{\text{disc}} = \frac{1}{T}\sum_{t=1}^T |z_t^L - z_t^R|$, measuring vertical misalignment. These metrics are particularly relevant for evaluating dual-arm systems where tight coordination is essential for tasks like cloth folding and bimanual assembly.

\subsection{Data and Sample Evaluation}\label{sec:eval_data_sample}

As foundation models scale to billions of parameters, the bottleneck increasingly lies not in model capacity but in data quality and sample efficiency. This subsection examines information-theoretic quality metrics, gradient-based data curation, and learning curve analysis.

\paragraph{Dataset Quality via Information Theory.}
\textbf{DemInf}~\cite{hejna2025robot} frames dataset quality as a mutual information maximization problem: $I(S; A) = H(A) - H(A|S)$, where $S$ represents states and $A$ represents actions. High action entropy $H(A)$ indicates broad state-space coverage, while low conditional entropy $H(A|S)$ indicates deterministic, unambiguous actions given states. To quantify the marginal contribution of individual demonstrations, DemInf computes:
\begin{equation}
  \Delta I_i = I(S; A \mid \mathcal{D}) - I(S; A \mid \mathcal{D} \setminus \{z_i\})
\end{equation}
where $z_i$ is the $i$-th demonstration. Demonstrations with high $\Delta I_i$ provide unique information and should be prioritized, while those with near-zero $\Delta I_i$ are redundant. Practical implementation involves training a VAE to compress high-dimensional observations into low-dimensional latents, then applying the k-NN-based KSG estimator to compute mutual information.

\paragraph{Influence Functions for Data Curation.}
Complementing information-theoretic approaches, \textbf{CUPID}~\cite{agia2025cupidcuratingdatarobot} and \textbf{QoQ}~\cite{lee2026qoq} employ influence functions to identify demonstrations that most impact validation performance:
\begin{equation}
  \text{Influence}(z_i) = -\nabla_\theta \mathcal{L}(z_{\text{val}})^\top H_\theta^{-1} \nabla_\theta \mathcal{L}(z_i)
\end{equation}
where $H_\theta = \nabla^2_\theta \mathcal{L}_{\text{train}}$ is the Hessian matrix, approximated via TRAK or EK-FAC for computational tractability. Positive influence scores indicate beneficial demonstrations, while negative scores suggest harmful outliers that degrade generalization. This gradient-based perspective complements the self-supervised filtering discussed in \S\ref{sec:data_refinement}, forming a comprehensive data curation stack that spans information-theoretic, gradient-based, and self-supervised paradigms.

\paragraph{Sample Efficiency.}
To quantify how efficiently policies utilize training data, the \textbf{Area Under the Learning Curve} (AULC) integrates success rate over the number of demonstrations:
\begin{equation}
  \text{AULC} = \int_0^{N_{\max}} SR(n)\, dn \approx \sum_{k} SR(n_k) \cdot \Delta n_k
\end{equation}
where $SR(n)$ is the success rate after training on $n$ demonstrations. Higher AULC indicates superior data utilization, a critical metric for few-shot imitation learning and data augmentation strategies (\S\ref{sec:data_augmentation}).

\subsection{Human Preference and Generalization Evaluation}\label{sec:eval_human_generalization}

While automated metrics provide scalability, human judgment remains the ultimate arbiter of manipulation quality in open-world settings. This subsection examines pairwise preference ranking, out-of-distribution generalization, and sim-to-real validation.

\paragraph{Human Preference via Pairwise Comparison.}
\textbf{RoboArena}~\cite{atreya2025roboarena} establishes a distributed real-world testbed for policy benchmarking, employing the \textbf{Bradley-Terry model} to infer global skill ratings from pairwise comparisons:
\begin{equation}
  P(\text{policy } i \succ j) = \frac{e^{\beta_i}}{e^{\beta_i} + e^{\beta_j}}
\end{equation}
where $\beta_i$ represents the latent skill of policy $i$. Maximum likelihood estimation recovers skill parameters via:
\begin{equation}
  \hat{\boldsymbol{\beta}} = \arg\max_{\boldsymbol{\beta}} \sum_{(i,j)\in\mathcal{C}} \left[y_{ij}\log P(i \succ j) + (1-y_{ij})\log P(j \succ i)\right]
\end{equation}
where $\mathcal{C}$ is the set of all pairwise comparisons and $y_{ij}=1$ if evaluators prefer policy $i$. Empirically, rankings converge after approximately 100 comparisons. Complementing the distributed physical infrastructure of RoboArena, \textbf{AutoEval}~\cite{zhou2025autoeval} automates success detection via vision-language models, enabling large-scale evaluation without manual trial-by-trial annotation---addressing the scalability bottleneck where manual annotation becomes prohibitively expensive as dataset scale and task diversity expand.

\paragraph{Out-of-Distribution Generalization.}
Generalization to novel environments is quantified via the \textbf{OOD success rate gap}: $\Delta SR_{\text{OOD}} = SR_{\text{in-dist}} - SR_{\text{out-dist}}$. REALM categorizes perturbations into three types: \textbf{visual} (background, lighting, color variations), \textbf{semantic} (object category substitution), and \textbf{physical} (initial pose offsets). The RMSD metric (defined in \S\ref{sec:eval_task_execution}) aggregates robustness across these perturbation types, with ideal generalist models exhibiting RMSD $\to 0$. These metrics align with existing benchmarks such as LIBERO~\cite{liu2023libero} and SimplerEnv~\cite{li2025simpler}, which systematically vary environmental factors to stress-test policy robustness.

\paragraph{Sim-to-Real Validation.}
\textbf{SIMPLER}~\cite{li2025simpler} establishes empirical bridges between simulated benchmarks and real-world deployment by quantifying \textbf{simulation ranking validity}. The \textbf{Mismatched Ranking Violations} (MMRV) metric measures the fraction of policy pairs whose relative ranking differs between simulation and reality:
\begin{equation}
  \text{MMRV} = \frac{1}{\binom{N}{2}} \sum_{i < j} \mathbb{1}\left[\text{rank}_{\text{sim}}(i,j) \neq \text{rank}_{\text{real}}(i,j)\right]
\end{equation}
where $N$ is the number of policies evaluated. Lower MMRV indicates higher sim-to-real correlation, validating simulation as a reliable proxy for real-world performance. Building on the simulation infrastructure discussed in \S\ref{sec:simulation_infrastructure}, SIMPLER demonstrates that carefully calibrated simulators can achieve low MMRV, enabling cost-effective policy screening before expensive real-world deployment. Extending this validation hierarchy one layer further, \textbf{WorldArena}~\cite{shang2026worldarena} asks whether \emph{learned world models} can themselves serve as reliable evaluation environments---replacing not only real-world trials but also hand-crafted simulators. WorldArena evaluates embodied world models (EWMs) along three functional roles: as \emph{data synthesis engines} (measuring the policy success rate gain $\Delta SR$ from model-generated training data), as \emph{policy evaluators} (measuring the Pearson correlation $\rho$ between world-model-predicted and simulator-measured success rates), and as \emph{action planners} (measuring closed-loop task success rate $SR_{\text{closed-loop}}$ when the world model drives autoregressive rollouts decoded into actions via an inverse dynamics model). Perceptual fidelity is assessed through standard video quality metrics (FID, FVD, SSIM, LPIPS) alongside a VLM-as-Judge protocol that scores physical interaction realism on a normalized Likert scale. Together, SIMPLER and WorldArena establish a three-tier evaluation trust chain---real world $\to$ simulator $\to$ world model---providing increasingly scalable but progressively approximate proxies for physical deployment.

\subsection{Energy and Efficiency Evaluation}\label{sec:eval_energy}

As embodied systems transition from laboratory prototypes to continuous industrial deployment, energy consumption emerges as a critical constraint alongside task performance. This subsection examines electrodynamic power models, component-wise energy decomposition, and normalized efficiency metrics that enable cross-platform comparison.

\paragraph{Electrodynamic Power Models.}
The instantaneous power consumption of joint $i$ is modeled as:
\begin{equation}
  P_i(t) = \frac{\tau_i^2(t)}{K_{T,i}^2} R_i + b_i \dot{q}_i^2(t) + P_{\text{idle},i}
\end{equation}
where the first term represents copper loss (dominant, proportional to torque squared), the second term captures viscous friction, and the third term accounts for idle electronics power. Total energy consumption integrates over the trajectory: $E = \int_0^T \sum_{i=1}^{n} P_i(t)\, dt$. For simulation environments lacking detailed motor parameters, a \textbf{torque-squared proxy} provides a parameter-free alternative: $E_{\text{proxy}} = \sum_{i=1}^{n} \int_0^T \tau_i^2(t)\, dt$, which correlates strongly with true energy consumption in practice.

\paragraph{Component-wise Energy Decomposition.}
The \textbf{ECDP}~\cite{heredia2023ecdp} provides a principled three-stage methodology---excitation trajectory design, system identification, and energy disaggregation---to decompose total instantaneous power into four physically interpretable components:
\begin{equation}
  P_{\text{total}}(t) = \underbrace{P_{\text{elec}}}_{\text{electronics}} + \underbrace{P_{\text{friction}}(t)}_{\text{mechanical friction}} + \underbrace{P_{\text{copper}}(t)}_{\text{copper loss}} + \underbrace{P_{\text{mech}}(t)}_{\text{useful mechanical work}}
\end{equation}
Specifically, the electronics component $P_{\text{elec}} = \sum_{i=1}^{n} P_{\text{idle},i}$ is a motion-independent constant representing control-box standby power; the friction component $P_{\text{friction}}(t) = \sum_{i=1}^{n} (b_i \dot{q}_i^2(t) + f_{c,i} |\dot{q}_i(t)|)$ captures both viscous ($b_i$) and Coulomb ($f_{c,i}$) friction losses; the copper loss $P_{\text{copper}}(t) = \sum_{i=1}^{n} \tau_i^2(t) R_i / K_{T,i}^2$ is the dominant motion-dependent term, where $K_{T,i}$ and $R_i$ are the motor torque constant and winding resistance respectively; and the mechanical power $P_{\text{mech}}(t) = \sum_{i=1}^{n} \tau_i(t)\dot{q}_i(t)$ represents useful work delivered to the load. Each component integrates to yield its energy share $E_k = \int_0^T P_k(t)\,dt$, with the ratio $\rho_k = E_k / E_{\text{total}}$ diagnosing where optimization effort should be directed. Validated across UR5e, UR10e, Franka FR3, and Kinova Gen3 platforms, ECDP reveals a consistent hierarchy $\rho_{\text{elec}} \gg \rho_{\text{copper}} \approx \rho_{\text{friction}} \gg \rho_{\text{mech}}$, with 60--90\% of total energy consumed by control electronics rather than mechanical motion. This finding has profound implications for energy-aware policy design: trajectory optimization alone has a low ceiling for energy savings, and substantial gains require co-optimization of hardware electronics alongside learned control policies.

\paragraph{Normalized Efficiency Metrics.}
To enable cross-platform comparison, several dimensionless metrics have been proposed. 
The \textbf{local energy indicator}~\cite{vidussi2021energy} quantifies configuration-dependent efficiency by combining the robot's inertia matrix $\mathbf{M}(\mathbf{q})$ with motor electrical parameters. Defining the electrical weight matrix $\mathbf{W} = \text{diag}(R_1/K_{T,1}^2, \ldots, R_n/K_{T,n}^2)$, the energy-weighted inertia matrix is constructed as $\mathbf{B}(\mathbf{q}) = \mathbf{W}^{1/2}\mathbf{M}(\mathbf{q})\mathbf{W}^{1/2}$, yielding the local index $\eta_E(\mathbf{q}) = 1/\|\mathbf{B}(\mathbf{q})\|_F$, where $\|\cdot\|_F$ denotes the Frobenius norm. This formulation---analogous to Yoshikawa's manipulability ellipsoid but targeting energy rather than velocity---enables workspace-wide \emph{energy maps} without running dynamics simulations. Integrating along a trajectory yields the \textbf{Trajectory Energy Index} (TEI): $\text{TEI} = \int_0^T \eta_E(\mathbf{q}(t))\,dt$, which exhibits a monotonic inverse relationship with true energy consumption ($E \propto 1/\text{TEI}$), allowing rapid comparison of candidate task placements from configuration sequences alone. This enables TEI-guided workspace optimization for both classical and learning-based manipulation pipelines without running full dynamics simulations. Recent work~\cite{heredia2025evaluating} introduces three normalized metrics: \textbf{energy utilization} $f_U = E_{\text{mech}}/E_{\text{total}} \in [0,1]$ (higher is better), \textbf{conversion efficiency} $f_C = E_{\text{out}}/E_{\text{in}} \in [0,1]$, and \textbf{reliability} $f_R = 1 - \sigma_P/\bar{P} \in [0,1]$, where $\sigma_P$ and $\bar{P}$ are power standard deviation and mean respectively. These dimensionless metrics enable fair comparison across robots of different scales. Finally, the \textbf{Cost of Transport}~\cite{tucker1975energetic}, originally developed for legged locomotion, normalizes energy by mass and distance: $\text{CoT} = E/(mg \cdot d)$, where $m$ is robot mass, $g$ is gravitational acceleration, and $d$ is end-effector displacement. Rezali et al.~\cite{rezali2025lstm} further demonstrate that LSTM-based surrogate models can map trajectory parameters to energy consumption, enabling multi-objective Pareto optimization over execution time, energy, and jerk simultaneously via NSGA-II\@.

While energy metrics remain underexplored in current embodied manipulation benchmarks, the deployment constraints outlined in \S4.4 and \S\ref{sec:deployment_safety} necessitate their integration into future evaluation protocols. As embodied systems scale to continuous operation in industrial and domestic settings, energy efficiency will become as critical as task success rate.

The evaluation dimensions above expose recurring gaps that no single section can resolve in isolation---limitations that cut across hardware, data, learning, and deployment. The final section distills these into actionable research directions.

\section{Open Challenges and Future Directions}\label{sec:challenges}

The organizational chain traced through this survey---from embodiment constraints (\S\ref{sec:embodiment_hardware}) through data design (\S\ref{sec:embodiment_dataset}) to learning paradigms (\S\ref{sec:embodiment_method}) and evaluation (\S\ref{sec:evaluation})---reveals that the most persistent bottlenecks are not confined to any single stage but arise at the \emph{interfaces} between stages: representation gaps between perception and planning, safety gaps between training and deployment, and standardization gaps across the ecosystem. This section diagnoses these cross-cutting limitations through three complementary perspectives: representation and reasoning bottlenecks (\S\ref{sec:representation_reasoning}) address what models can perceive and how they plan; deployment and safety challenges (\S\ref{sec:deployment_safety}) examine real-world operation constraints; and ecosystem infrastructure gaps (\S\ref{sec:ecosystem_infrastructure}) identify systemic barriers in data, evaluation, and community practices.

\subsection{Representation and Reasoning Bottlenecks}\label{sec:representation_reasoning}

Despite advancements in vision-language models, fundamental gaps persist in how embodied systems perceive physical dynamics and execute complex, multi-step actions. To address these representation bottlenecks, 
key research directions can focus on: (i) physics-grounded multimodal integration that moves beyond static visual priors by deeply fusing proprioceptive and tactile data (e.g., leveraging dual-path tactile encoders in OmniVTLA~\cite{cheng2025omnivtla}, and localized force feedback in ForceVLA~\cite{yu2025forcevla}) and enforcing physical reasoning through dynamic pre-training objectives; 
(ii) morphology-agnostic generalization that decouples task-level objectives from embodiment-specific kinematics to overcome the reliance on rigid serial manipulators---starkly highlighted by the zero-shot failure of standard Type~I models on soft arms~\cite{su2025bridgingembodiment}---and enable fluid transfer to novel morphologies via universal representation spaces (e.g., 3D point flows in PointWorld~\cite{huang2026pointworld}) or by extracting cross-embodiment physical priors directly from massive-scale human egocentric data to scale high-DoF dexterous manipulation (e.g., EgoScale~\cite{zheng2026egoscale});
(iii) mitigating catastrophic forgetting and the reasoning-acting trade-off: as embodied foundation models are heavily fine-tuned on narrow, embodiment-specific manipulation trajectories, they frequently suffer from catastrophic forgetting, degrading the broad semantic reasoning and visual understanding inherited from their foundation VLM backbones. Bridging the gap between high-level generalist reasoning and low-level specialist control requires novel continuous learning paradigms, modular architectures (such as Mixture-of-Experts or decoupled action heads, e.g., InstructVLA~\cite{yang2025instructvla}), and balanced instruction-tuning strategies to ensure the acquisition of physical dexterity does not cannibalize universal semantic knowledge.

\subsection{Deployment and Safety Challenges}\label{sec:deployment_safety}

Deploying large-scale neural policies into real-world environments introduces strict constraints on latency, safety, and energy consumption that isolated algorithmic improvements cannot resolve. 
Key deployment research directions include: (i) latency-aware edge inference and sustainable deployment via architecture-specific acceleration, model compression, and hardware-software co-design to consistently meet the strict control frequencies (e.g., $>10$~Hz) required for closed-loop manipulation while optimizing standardized energy consumption metrics for continuous adaptation;
(ii) verifiable safety and robustness through the integration of neural policies with hard-constraint safety shields (e.g., control barrier functions) and out-of-distribution detection to defend against severe adversarial perturbations (e.g., multimodal backdoor injection attacks exposed by BadVLA~\cite{zhou2025badvla}); 
and (iii) runtime anomaly detection and execution resilience to actively mitigate policy collapses during critical state transitions. Since end-to-end models frequently experience sudden control degradation during complex motion-transition phases (e.g., the vision-proprioception failures highlighted by~\cite{lu2026whenwould}).

\subsection{Ecosystem Infrastructure}\label{sec:ecosystem_infrastructure}

The embodied manipulation ecosystem is severely bottlenecked by proprietary, fragmented data pipelines and inconsistent evaluation protocols that hinder reproducible progress. 
Establishing a mature infrastructure requires community-wide alignment on: (i) unified multimodal data standards that create universal formats for action spaces and sensory data (e.g., OHRA~\cite{wei2026one} establishing a parameterized canonical URDF representation to unify diverse dexterous hand architectures), 
systematically expand embodiment diversity, and crucially integrate failure trajectories and robust linguistic-action grounding (e.g., LingBot-VA~\cite{li2026lingrobotva}) back into training corpora; 
(ii) diagnostic evaluation frameworks that shift from binary success rates to standardized, reconfigurable real-world benchmarks rigorously categorizing failure modes across perception, planning, and execution (e.g., RoboChallenge~\cite{yakefu2025robochallenge} building a real-world robotics testing and evaluation platform supporting diverse tasks and multiple robots, and WorldArena~\cite{shang2026worldarena} establishing the first unified standard for evaluating embodied world models); 
and (iii) closed-loop sim-to-real validation leveraging quantifiable predictive validity frameworks to bridge the reality gap by correlating simulation benchmarks with dynamic, physical-world performance metrics.


\subsection{Open Challenges Along the Constraint Chain}\label{sec:open_challenges}

Building on the discussion above, the main open problems of embodied manipulation can be read as a linked chain of constraints. Limits in embodiment shape the data that can be collected and aligned. Those data limitations then narrow the architectural choices available to learning systems. When architectural interfaces remain unstable, evaluation becomes harder to interpret and less useful as feedback. Once these systems leave controlled settings, trustworthiness becomes an additional challenge that cannot be inferred from task success alone. The first four challenges follow this chain directly. The fifth becomes central at deployment time.
\paragraph{Embodiment foundations.}\label{sec:challenge_embodiment}
A first challenge concerns the lack of shared abstractions across robot embodiments. 
Most manipulation systems are still organized around local sensing layouts, end-effector geometry, kinematics, and control interfaces. 
That makes progress hard to transfer across platforms. 
The same task may require different contact strategies, motion scales, and feedback loops on different robots, while signals that look similar can carry very different functional meaning. 
This difficulty appears early, and its effects carry forward. If embodiments cannot be related on common terms, downstream data collection, model design, and evaluation all inherit that mismatch. 
The field therefore needs an abstraction layer that preserves task-relevant meaning across heterogeneous platforms without ignoring their physical constraints. 
In practice, this means more consistent ways to describe state, action, and contact, better mechanisms for retargeting trajectories and policies across morphologies, and broader support for modalities such as force, touch, bimanual coordination, and mobile manipulation. 
Recent work has begun to move in this direction, yet a stable cross-embodiment foundation is still missing.
\paragraph{Data ecosystem.}\label{sec:challenge_data}
A second challenge concerns the data ecosystem that supports manipulation learning. Real-world teleoperation still offers the most faithful record of contact, failure, recovery, and state transition, yet it is slow, costly, and difficult to scale across tasks, sites, and robot platforms. Simulation and synthetic data improve throughput and coverage, though they remain imperfect substitutes when success depends on friction, compliance, deformation, timing, and embodiment-specific control effects. The difficulty does not end with collection. Even large datasets are often hard to combine because their observation schemas, action conventions, time bases, calibration assumptions, and task annotations do not line up cleanly. For manipulation, more data alone will not solve the problem. The field needs data that are faithful enough to train on, structured enough to align, and reusable enough to accumulate across projects. That points to a broader infrastructure agenda: scalable collection, shared formats and metadata, stronger curation and versioning, and deliberate inclusion of failure cases and long-tail interactions. Recent datasets and toolchains have expanded what is possible, yet the data substrate remains fragmented.
\paragraph{Learning architectures.}\label{sec:challenge_learning}
A third challenge concerns architectural design. End-to-end policies, hierarchical systems, vision-language-action models, and planning-augmented pipelines each capture part of the problem, yet they still differ on several basic questions: what counts as an executable action unit, what kind of intermediate structure is genuinely useful, and where embodiment-specific adaptation should live. Because these choices are often made locally, comparisons across methods are easily confounded by interface design. A reported gain may come from a more convenient decomposition or action representation, so it is often hard to tell what has really improved at the architectural level. The key issue, then, is to stabilize the interfaces that connect reasoning, policy learning, and control. Progress here depends on clearer definitions of action granularity, more reusable intermediate representations, and system interfaces that remain meaningful across tasks, embodiments, and control regimes. Recent work on action chunking, continuous decoding, and intermediate supervision shows that these design choices matter a great deal, yet a shared architectural basis has not been established.
\paragraph{System evaluation.}\label{sec:challenge_evaluation}
A fourth challenge concerns evaluation. Simulation makes controlled comparison possible, yet widely used benchmarks still capture only part of what matters in deployment. Real-world evaluation is closer to practice, though it is slower, harder to standardize, and more difficult to reproduce across platforms and laboratories. When a small set of benchmark suites becomes dominant, its local assumptions about task format, observation setting, and success criteria begin to shape model development more broadly. The result is a familiar distortion: systems become easier to compare on paper, while properties that matter in the wild remain weakly measured. Robustness under shift, recovery from intermediate failure, long-horizon consistency, sensitivity to control error, and execution efficiency can all fall outside the main evaluation signal. A more balanced evaluation framework is needed, one that combines reproducibility with deployment relevance and diagnostic value. That includes more stable task and success abstractions, richer measurements of failure and recovery, and scalable real-world protocols for logging, scoring, and cross-site replication. Recent benchmarks have made evaluation more rigorous, yet a balanced and widely usable framework has not emerged.
\paragraph{Trustworthiness.}\label{sec:challenge_safety}
A fifth challenge concerns trustworthiness once embodied manipulation systems move into shared human environments. At that point, task success is no longer a sufficient proxy for readiness. A policy may perform well on standard inputs and still fail under rare conditions, hidden triggers, or interactions that were underrepresented during training. Failures are also hard to diagnose because perception, reasoning, and action are tightly coupled, especially in end-to-end systems. Even when intermediate reasoning is exposed, it does not necessarily reveal why a particular action was chosen. Trustworthiness therefore has to be treated as a first-class concern. That means building evaluation procedures that can surface hidden failure modes, interpretability tools that make the causal basis of decisions more legible, and interaction designs that let people anticipate system behavior and intervene safely when needed. Recent work has started to address safety alignment, stress testing, interpretability, and shared autonomy, yet a general framework for trustworthy deployment is still missing.
Taken together, these challenges suggest that progress in embodied manipulation depends as much on shared foundations across embodiment, data, architecture, evaluation, and deployment as it does on better models. Recent work clearly points in this direction, yet the common infrastructure needed to support cumulative progress is still being built.

\section{Conclusion}

This survey has analyzed embodied manipulation intelligence through the organizational chain of \emph{embodiment $\to$ data $\to$ learning $\to$ evaluation}. We showed that hardware constraints---end-effector morphology, platform kinematics, middleware latency---propagate directly into data design, which in turn shapes what learning paradigms can achieve and how deployment must be engineered. The three-tier dataset taxonomy (skill/task/foundation) provides a principled basis for comparing data investments; the cognitive planning / policy learning dual pathway with Type~I/II/III integration taxonomy supersedes the monolithic VLA narrative by revealing structural trade-offs between semantic generality and physical precision; and the three-stage deployment lifecycle (offline/online/runtime) formalizes the engineering constraints that bridge laboratory models to real-world operation. The open challenges identified in \S\ref{sec:challenges}---representation bottlenecks, deployment safety, and ecosystem fragmentation---arise precisely at the interfaces between these stages, suggesting that future progress depends less on improving individual components than on tightening the coupling across the entire chain.

\section{Acknowledgments}

Identification of funding sources and other support, and thanks to
individuals and groups that assisted in the research and the
preparation of the work should be included in an acknowledgment
section, which is placed just before the reference section in your
document.

This section has a special environment:
\begin{verbatim}
  \begin{acks}
  ...
  \end{acks}
\end{verbatim}
so that the information contained therein can be more easily collected
during the article metadata extraction phase, and to ensure
consistency in the spelling of the section heading.

Authors should not prepare this section as a numbered or unnumbered {\verb|\section|}; please use the ``{\verb|acks|}'' environment.

%%
%% The next two lines define the bibliography style to be used, and
%% the bibliography file.
\bibliographystyle{ACM-Reference-Format}
\bibliography{sample-base}

\end{document}

